% !TeX root = ../Book1.tex
% 纲要标题：# 第一编　运算、群与对称
% 纲要位置：计划纲要.md，第 129 行。
% 本编导读按阅读目标列出关键结果及可以暂缓的内容；第8章交换化练习需回读第7章的交换化泛性质。

\BookPart{运算、群与对称}{b1:part:01}
\BookPartGuide

整数相加、置换复合和几何图形的对称变换看似是不同的计算，
却都涉及一个共同的问题：若干对象经过运算以后，仍能否留在原来的集合中，
运算的次序又怎样影响结果？群的概念提取了其中一类最常见的规律。
它既容纳具体计算，也使我们能够讨论同态、商结构和分解，
比较表面不同的对象是否具有相同的运算结构。

\ProseParagraphBreak
本编从熟悉的运算出发引入群的定义。
\chapref{b1:ch:01}分别讨论结合性、单位元和可逆性，
并以字的拼接和商集说明构造一个新对象时需要检查什么。
\chapref{b1:ch:02}建立群、子群、同态及商群的基本理论；
\chapref{b1:ch:03}将这些概念用于循环群、置换群和几何对称。
读到这里，应当能够从一张乘法表或一组明确的变换出发，
判断哪些结构已经确定，哪些结论仍需证明。

\ProseParagraphBreak
有限群的结构往往可以通过它在别的集合上的作用来研究。
\chapref{b1:ch:04}把计数转化为轨道与稳定子群的关系，
\chapref{b1:ch:05}由此建立 Sylow 理论。
这些定理给出很强的整除与共轭限制，却不能单凭一个群的阶就完成任意分类。
\chapref{b1:ch:06}转向直积、半直积和群扩张，
解释已知的较小群怎样组合成新群；
\chapref{b1:ch:07}则借正规列比较这些结构的逐层分解，
证明合成因子的唯一性，并讨论可解性、幂零性和交错群的单性。

\ProseParagraphBreak
最后两章回答另外两类基本问题。
给定生成元以后，只要求群公理，能得到怎样的群？
再加入指定关系，又怎样得到所需的商群？
这些问题在\chapref{b1:ch:08}中通过约化字、展示和自由积处理。
\chapref{b1:ch:09}回到交换情形，以整数关系和整除运算证明有限生成交换群的结构定理。
这里先用群论方法证明；第二卷第四章“主理想整环上的模与标准形”将从模论的角度重新解释。

\begin{center}
\begin{tikzcd}[column sep=large,row sep=large]
 & \begin{gathered}\text{第4--5章}\\\text{群作用与Sylow理论}\end{gathered}\arrow[dr] & \\
 \begin{gathered}\text{第1--3章}\\\text{基本概念与例子}\end{gathered}
   \arrow[ur]\arrow[r]\arrow[dr]\arrow[ddr]
 & \begin{gathered}\text{第6章}\\\text{乘积与扩张}\end{gathered}\arrow[r]
 & \begin{gathered}\text{第7章}\\\text{正规列与结构}\end{gathered}\\
 & \begin{gathered}\text{第8章}\\\text{自由群与展示}\end{gathered} & \\
 & \begin{gathered}\text{第9章}\\\text{有限生成交换群}\end{gathered} &
\end{tikzcd}
\end{center}

图中箭头概括各部分的主要依赖；具体阅读所需的结果见\cref{b1:tab:part01-reading}。
第八章和第九章可以在掌握基本构造后分别阅读。
其中第八章关于交换化的讨论与习题还需\cref{b1:cor:abelianization-universal}，
不能把图中的分支理解为第八章完全不使用第七章。
第八章第四节带选读星号，用图的方法证明自由群的子群仍是自由群。
该节将介绍所需的图论概念，不要求预先学习代数拓扑。

\ProseParagraphBreak
本编预备知识为集合、映射、数学归纳法，以及整数的整除、
最大公因子和带余除法；矩阵运算主要用于提供具体例子。
阅读时应特别留意映射的定义域和值域、左陪集与右陪集的区别，
以及“选定一个代表元”以后所得结果是否与选择无关。
遇到有限生成、有限阶或正规性等假设，可同时寻找去掉该假设后的反例。
各章习题及解答用于检验这些判断，并练习把同一个对象放在不同构造中计算。

\ProseParagraphBreak
本编的基础概念可参阅席南华的《基础代数》第一卷\cite[第2章, 第5章]{Xi2016BasicAlgebraI}，或柯斯特利金的《代数学引论》第一卷：基础代数（张英伯译）\cite[第4章]{Kostrikin2006AlgebraIChinese}。
群论的系统学习可选 Dummit 与 Foote 的《Abstract Algebra》\cite[Chapters 1--6]{DummitFoote2003}和 Rotman 的《An Introduction to the Theory of Groups》\cite[Chapters 1--7, 10--11]{Rotman1995TheoryGroups}。
Lang 的《Algebra》\cite[Chapter I, \S\S 1--8, 10, 12]{Lang2002}适合在熟悉基本例子后查阅结构性的论述。
自由群、正规列及交换群的专题读物分别放在相应章首，读者不必同时通读所有参考书。

\begingroup
\renewcommand{\thetable}{I.\arabic{table}}
\renewcommand{\theHtable}{part-overview.1.\arabic{table}}
\begin{longtable}{@{}>{\raggedright\arraybackslash}p{0.16\linewidth}>{\raggedright\arraybackslash}p{0.47\linewidth}>{\raggedright\arraybackslash}p{0.29\linewidth}@{}}
\caption{第一编的阅读目标与关键结果}\label{b1:tab:part01-reading}\\
\toprule
阅读目标 & 需要掌握的结果 & 可以暂缓的内容 \\
\midrule
\endfirsthead
\toprule
阅读目标 & 需要掌握的结果 & 可以暂缓的内容 \\
\midrule
\endhead
商结构与同态方法 & 字母映射的延拓与商映射的下降（\cref{b1:thm:free-monoid}、\cref{b1:thm:monoid-quotient}）；第一、第二、第三同构定理及子群对应（\secref{b1:sec:0204}） & 群化可作为综合应用回读；不需要无限基数或超限归纳。 \\[.7ex]
有限群的结构 & 置换模型，轨道—稳定子群定理；Sylow 三定理（\cref{b1:thm:sylow-existence}、\cref{b1:thm:sylow-conjugacy}、\cref{b1:thm:sylow-count}）；内半直积判别与正规列 & 自由群子群的图论证明不参与这部分的论证。 \\[.7ex]
根式理论的群论准备 & 同构定理、可解性的子群与商群及扩张性质（\cref{b1:thm:solvable-closure}）；交错群单性及对称群可解性判据（\cref{b1:thm:alternating-simple}、\cref{b1:cor:symmetric-nonsolvable}） & 幂零群专题、非交换核的一般因子集计算不是根式可解性判据的前提。 \\[.7ex]
生成关系与交换群分类 & 自由群及展示的映射性质（\cref{b1:thm:free-group-universal}、\cref{b1:prop:presentation-homomorphism}）；有限秩子群定理、Smith 引理和交换群结构定理（\cref{b1:thm:finite-rank-free-abelian-subgroup}、\cref{b1:lem:smith-integer}、\cref{b1:thm:fga-structure}） & 第九章分类并不依赖第八章的子群自由性；两章的一般无限秩证明都可暂缓。 \\[.7ex]
组合群论初步 & 约化字的唯一性、自由积标准形；生成树给出的闭路径自由基及 Schreier 秩公式（\cref{b1:lem:graph-loop-free}、\cref{b1:thm:schreier-rank}） & 不必先完成 Smith 计算。只读有限陪集图时，可暂缓一般无限图的选择论证。 \\
\bottomrule
\end{longtable}
\endgroup

表中列出的暂缓内容仍可按兴趣阅读；遇到正文明确调用的结果，再按链接补齐即可。
例如，学习第八章的交换化习题时，只需先读导出子群及交换化的映射性质，
无须因此读完合成列和全部幂零群理论。

% !TeX root = ../../Book1.tex
% 纲要标题：## 第1章　代数对象、映射与基本构造
% 纲要位置：计划纲要.md，第 131 行。

\chapter{代数对象、映射与基本构造}
\label{b1:ch:01}
\BookChapterNumber{b1:ch:01}

\begin{chapterabstract}
本章从熟悉的加法、复合与乘法出发，说明运算公理如何决定可以进行的推理。自由字、商集和生成子幺半群分别回答三个问题：如何记录运算表达式，如何把应当相同的表达式合并，以及如何用少量元素描述一个结构。映射的存在性与唯一性把这些构造联系起来，为下一章的群和商群准备共同语言。
\end{chapterabstract}

\begin{learninggoals}
\begin{itemize}
\item 区分结合律、交换律、单位元、逆元和消去律的作用，能用具体反例说明它们并不相互替代。
\item 构造自由幺半群和商幺半群，逐项核验运算的良定义及同态的延拓。
\item 分清生成元决定映射的唯一性与生成关系限制映射的存在性，并用泛性质证明指定同构的唯一性。
\end{itemize}
\end{learninggoals}

本章只要求集合、映射、整数和基本矩阵运算。读到新结构时，先找出它的元素、运算和相等标准，再问哪些映射保持这些资料。特别应亲自核对字的拼接与商集的计算；这些看似朴素的步骤将在商群、商环和模的构造中反复出现。

% !TeX root = ../../Book1.tex
% 纲要标题：### 1.1 运算与结构
% 纲要位置：计划纲要.md，第 133 行。

\section{运算与结构}
\label{b1:sec:0101}

% TODO: 按以下纲要要点撰写本节。

% 纲要内容（仅作写作计划，尚非教材正文）：
%
% * 集合上的一元、二元运算及其公理
% * 从整数加法、置换复合到矩阵乘法
% * 结合律、交换律、单位元与逆元的不同作用
%

\subsection{一元、二元运算及其公理}
\label{b1:subsec:0101:01}

待撰写。

\subsection{整数、置换与矩阵上的运算}
\label{b1:subsec:0101:02}

待撰写。

\subsection{结合律、交换律、单位元与逆元}
\label{b1:subsec:0101:03}

待撰写。

% !TeX root = ../../Book1.tex
% 纲要标题：### 1.2 半群与幺半群
% 纲要位置：计划纲要.md，第 139 行。

\section{半群与幺半群}
\label{b1:sec:0102}

% 纲要内容（仅作写作计划，尚非教材正文）：
%
% * 半群、幺半群及同态
% * 自由幺半群与字的拼接
% * 变换幺半群、消去律与群化的动机
%

\subsection{半群、幺半群及同态}
\label{b1:subsec:0102:01}

把结合律作为共同条件，便能同时讨论字的拼接、映射复合和矩阵乘法。由于这些运算通常不交换，本节使用乘法记号 \(ab\)，但不把因子顺序视为可以任意更换。

\begin{definition}\label{b1:def:monoid}
带有结合的二元运算的非空集合称为\term[banqun]{半群}[semigroup]。具有单位元的半群称为\term[yaobanqun]{幺半群}[monoid]，其单位元一般记为 \(1\)。幺半群 \(M,N\) 之间的\term[tongtai]{同态}[homomorphism]是满足
\[
f(ab)=f(a)f(b),\qquad f(1_M)=1_N
\]
的映射 \(f:M\rightarrow N\)。半群同态只要求乘法等式。
\end{definition}

单位保持条件不能从乘法保持条件自动得到。例如从单元素幺半群到整数乘法幺半群、把唯一元素送到 \(0\) 的映射保持乘法，却不保持单位。今后“幺半群同态”总包含保持单位的要求；这使空乘积与非空乘积具有相同的映射规则。

\ProseParagraphBreak
\((\mathbb N,+)\) 是以 \(0\) 为单位元的交换幺半群；正整数加法构成半群，却没有单位元。若 \(M\) 是幺半群，递归定义 \(a^0=1\)、\(a^{n+1}=a^na\)。结合律给出 \(a^{m+n}=a^ma^n\)：固定 \(m\) 对 \(n\) 归纳即可。于是任取 \(a\in M\)，映射 \(n\mapsto a^n\) 都是从 \((\mathbb N,+)\) 到 \(M\) 的同态。

\begin{proposition}\label{b1:prop:monoid-maps}
幺半群同态的复合仍是同态；恒等映射是同态。若同态 \(f:M\rightarrow N\) 是双射，则逆映射 \(f^{-1}\) 也保持乘法和单位。
\end{proposition}
\begin{proof}
对于同态 \(f:M\rightarrow N\)、\(g:N\rightarrow P\)，有
\((gf)(ab)=g(f(a)f(b))=(gf)(a)(gf)(b)\)，并且
\((gf)(1_M)=1_P\)。恒等映射的结论直接由定义得出。若 \(f\) 双射，写 \(x=f(a),y=f(b)\)，则
\(f^{-1}(xy)=f^{-1}(f(ab))=ab=f^{-1}(x)f^{-1}(y)\)，且
\(f^{-1}(1_N)=1_M\)。
\end{proof}

这样的双射同态称为\term[tonggou]{同构}[isomorphism]。它说明两个幺半群的乘法结构相同，只是元素的名称不同。若只比较底集合的大小，就会丢掉乘法信息；后面将反复用元素阶、交换性等性质区别同样大小的结构。

\subsection{自由幺半群与字的拼接}
\label{b1:subsec:0102:02}

给定集合 \(X\)，把它的元素看成字母。长度为 \(n\) 的\term[zi]{字}[word]是有序序列 \((x_1,\ldots,x_n)\)，通常写成 \(x_1\cdots x_n\)。长度零的字只有一个，称为空字，记为 \(\varepsilon\)。不同长度的字始终视为不同对象，即使 \(X\) 自身含有数或序列。

\begin{definition}\label{b1:def:free-monoid}
集合 \(X\) 上的\term[ziyou-yaobanqun]{自由幺半群}[free monoid]是所有有限字的集合
\[
X^*\coloneqq\coprod_{n\geq0}X^n,
\]
运算为前后拼接，单位元为 \(\varepsilon\)。记 \(\iota:X\rightarrow X^*\) 为把字母送到长度一的字的映射。
\end{definition}
\symidx{\(X^*\)：自由幺半群}

符号 \(\coprod\) 表示不交并，即保留每个字的长度标记。拼接后各位置的字母就是先列出第一个字，再列出第二个字。因此两次拼接的结果不依赖括号；空字不增加任何位置，确为单位元。例如 \(X=\{a,b\}\) 时，\((ab)(baa)=abbaa\)，而 \((baa)(ab)=baaab\)。所以当至少有两个字母时，自由幺半群并不交换。单字母的情形只有 \(\varepsilon,a,a^2,\ldots\)，按长度同构于 \((\mathbb N,+)\)。

\begin{theorem}[字母映射的延拓]\label{b1:thm:free-monoid}
任给幺半群 \(M\) 和集合映射 \(u:X\rightarrow M\)，存在唯一幺半群同态
\(\widehat u:X^*\rightarrow M\) 满足 \(\widehat u\iota=u\)。它由
\[
\widehat u(\varepsilon)=1_M,\qquad
\widehat u(x_1\cdots x_n)=u(x_1)\cdots u(x_n)
\]
给出。
\end{theorem}
\begin{proof}
一个非空字有唯一的字母序列，右端乘积由结合律不依赖括号，所以公式定义了映射。两个字拼接时，右端乘积恰是两个原乘积的乘积，故它保持乘法；空字的指定值保证保持单位。每个非空字都是长度一的字的乘积，任何满足要求的同态都只能取上述值；在空字上也只能取单位元，所以延拓唯一。
\end{proof}

例如给 \(a,b\) 指定两个同阶方阵 \(A,B\)，每个字便成为按既定次序进行的矩阵乘积。可能有不同字取同一矩阵，但在 \(X^*\) 中它们仍是不同的字。“自由”所排除的是额外规定的字之间的等式，并不要求每次代入后的取值都不同。

\subsection{变换幺半群、消去律与群化的动机}
\label{b1:subsec:0102:03}

集合 \(X\) 的所有自映射在复合下构成\term[bianhuan-yaobanqun]{变换幺半群}[transformation monoid]，记为 \(\operatorname{End}_{\mathrm{Set}}(X)\)，单位元为恒等映射。前节的移位映射说明，其中可有左逆而无右逆。由自映射到置换的限制，正是从一般操作转向可逆操作。

\begin{definition}\label{b1:def:cancellation}
幺半群满足\term[xiaoqulv]{消去律}[cancellation law]，是指
\(ax=ay\Rightarrow x=y\) 与 \(xa=ya\Rightarrow x=y\) 都成立。
\end{definition}

自由幺半群满足消去律，因为等字有相同的逐位字母，去掉相同前缀或后缀仍得等字。然而非空字没有逆元：拼接后的长度是两字长度之和，不可能成为零。因此消去律本身不等于可逆性。

\ProseParagraphBreak
对交换幺半群，可以像从自然数构造整数那样加入形式差。以下给出消去律充分发挥作用的一种情形。为了使公式清晰，暂以加法写 \(M\)，单位元记为 \(0\)。

\begin{proposition}[交换消去幺半群的群化]\label{b1:prop:commutative-completion}
设 \(M\) 是满足消去律的交换幺半群。在 \(M\times M\) 上规定
\[
(a,b)\sim(c,d)\quad\Longleftrightarrow\quad a+d=c+b.
\]
其等价类集合 \(K(M)\) 在
\([(a,b)]+[(c,d)]=[(a+c,b+d)]\) 下是交换幺半群，每个元素都有逆元 \([(b,a)]\)。映射
\(j(a)=[(a,0)]\) 是单射同态。
\end{proposition}
\begin{proof}
自反性和对称性直接成立。若 \(a+d=c+b\) 且 \(c+f=e+d\)，相加后消去 \(c+d\)，得到 \(a+f=e+b\)，故有传递性。若两对代表元分别等价，则把相应两条等式相加可知它们的逐项和等价，所以加法良定义。结合律与交换律来自 \(M\)，零元为 \([(0,0)]\)，而
\([(a+b,b+a)]=[(0,0)]\)，给出所述逆元。最后，\(j(a)=j(c)\) 当且仅当 \(a=c\)，并且 \(j\) 保持加法与零元。
\end{proof}

这里“等价类”的一般性质将在下一节系统说明；上述证明已经逐项给出了本构造需要的检验。对 \(M=\mathbb N\)，映射 \([(a,b)]\mapsto a-b\) 是到 \(\mathbb Z\) 的双射并保持加法。一般情形也有同样的映射原则：若 \(u:M\rightarrow A\) 是到一个每个元素都可逆的交换幺半群的同态，则
\([(a,b)]\mapsto u(a)-u(b)\) 是唯一延拓 \(u\) 的同态。良定义来自
\(a+d=c+b\) 蕴含 \(u(a)+u(d)=u(c)+u(b)\)；唯一性来自
\([(a,b)]=j(a)-j(b)\)。这说明新加的形式差除了使元素可逆，没有强加不必要的关系。

\ProseParagraphBreak
交换性在这里允许整理求和，消去律则用于等价关系的传递性和嵌入判断。非交换情形不能直接把同一公式中的加号换成乘号；一般半群嵌入群的问题有额外障碍，本节不作这种推广。

% !TeX root = ../../Book1.tex
% 纲要标题：### 1.3 等价关系与商集
% 纲要位置：计划纲要.md，第 145 行。

\section{等价关系与商集}
\label{b1:sec:0103}

% TODO: 按以下纲要要点撰写本节。

% 纲要内容（仅作写作计划，尚非教材正文）：
%
% * 等价关系、划分与商映射
% * 运算何时能够下降到商集
% * 同余关系与良定义性的逐项检验
%

\subsection{等价关系、划分与商映射}
\label{b1:subsec:0103:01}

待撰写。

\subsection{商集上的诱导运算}
\label{b1:subsec:0103:02}

待撰写。

\subsection{同余关系与良定义性}
\label{b1:subsec:0103:03}

待撰写。

% !TeX root = ../../Book1.tex
% 纲要标题：### 1.4 生成与泛性质
% 纲要位置：计划纲要.md，第 151 行。

\section{生成与泛性质}
\label{b1:sec:0104}

% 纲要内容（仅作写作计划，尚非教材正文）：
%
% * 生成子对象与生成元
% * 自由对象的映射延拓性质
% * 存在性、唯一性与典范同构的区别
%
% ---
%

\subsection{生成子对象与生成元}
\label{b1:subsec:0104:01}

要描述一个很大的代数对象，直接列出所有元素通常并不可行。更有效的办法是先挑选少量元素，再说明从这些元素可以进行哪些运算。生成问题因而分成两个层次：哪些元素能够得到，以及同一个元素有哪些不同的表达式。

\begin{definition}\label{b1:def:generated-monoid}
幺半群 \(M\) 的子集 \(N\) 若包含 \(1_M\) 且对乘法封闭，称为\term[zi-yaobanqun]{子幺半群}[submonoid]。给定 \(S\subseteq M\)，包含 \(S\) 的最小子幺半群称为 \(S\) \term[shengcheng]{生成}[generate]的子幺半群，记为 \(\langle S\rangle_{\mathrm{mon}}\)。若它等于 \(M\)，则称 \(S\) 为一个生成元集合。
\end{definition}
\symidx{\(\langle S\rangle_{\mathrm{mon}}\)：生成子幺半群}

\begin{proposition}\label{b1:prop:generated-monoid}
\(\langle S\rangle_{\mathrm{mon}}\) 存在，既等于所有包含 \(S\) 的子幺半群的交，又等于 \(S\) 中元素的所有有限乘积组成的集合，其中包括空乘积 \(1_M\)。
\end{proposition}
\begin{proof}
包含 \(S\) 的子幺半群至少有 \(M\) 自身。它们的交含有单位元；若两个元素属于交，它们属于每个子幺半群，因而乘积仍属于每一个，也属于交。这给出最小子幺半群。另一方面，所有有限乘积的集合包含 \(S\) 和单位元，并且两个有限乘积的乘积仍为有限乘积；任意包含 \(S\) 的子幺半群都必须包含这些乘积。故两种描述相同。
\end{proof}

例如在 \((\mathbb N,+)\) 中，\(\{2,3\}\) 生成
\(\{0,2,3,4,5,\ldots\}\)。因为 \(1\) 不能写成非负整数系数的 \(2,3\) 之和，它不在其中；每个偶数至少为 \(2\) 时是若干个 \(2\) 的和，每个奇数至少为 \(3\) 时是 \(3\) 加一个偶数。这里 \(6=2+2+2=3+3\) 显示生成表达式一般不唯一。

\begin{proposition}\label{b1:prop:maps-on-generators}
若 \(S\) 生成幺半群 \(M\)，则两个同态 \(f,g:M\rightarrow N\) 在 \(S\) 上相等便在 \(M\) 上相等。
\end{proposition}
\begin{proof}
对非空乘积 \(s_1\cdots s_r\)，同态的值分别为
\(f(s_1)\cdots f(s_r)\) 与 \(g(s_1)\cdots g(s_r)\)，逐因子相等；空乘积的像都是 \(1_N\)。这些元素已穷尽 \(M\)。
\end{proof}

这个命题只保证唯一性，不保证任意指定生成元的像都能形成同态。例如上例中，若想把 \(2\) 送到 \(1\in\mathbb N\)、把 \(3\) 送到 \(1\in\mathbb N\)，则 \(6\) 按两种表达式应同时送到 \(3\) 与 \(2\)，产生矛盾。生成关系正是延拓必须遵守的限制。

\subsection{自由对象的映射延拓性质}
\label{b1:subsec:0104:02}

\begin{definition}\label{b1:def:universal-property}
用对所有目标对象的映射存在性和唯一性来刻画一个构造的条件，称为该构造的\term[fan-xingzhi]{泛性质}[universal property]。例如 \(X^*\) 连同字母映射 \(\iota:X\rightarrow X^*\) 的泛性质是：到任意幺半群的字母映射都能唯一延拓为同态。
\end{definition}

\cref{b1:thm:free-monoid}已经证明这一性质。它同时处理所有目标幺半群，比列举若干可代入的数或矩阵更强。自由对象包含的资料是“对象加指定生成映射”，不是仅仅一个没有标记的幺半群。

\begin{proposition}[生成元与关系]\label{b1:prop:monoid-generators-relations}
设 \(R\) 是 \(X^*\) 中若干对字组成的集合。存在包含这些字对的最小同余关系 \(\equiv_R\)。商幺半群 \(X^*/{\equiv_R}\) 具有如下性质：它到 \(M\) 的同态，恰好对应那些使每对 \((v,w)\in R\) 都满足
\(\widehat u(v)=\widehat u(w)\) 的映射 \(u:X\rightarrow M\)。
\end{proposition}
\begin{proof}
把所有包含 \(R\) 的同余关系相交。这样的关系至少有把所有字视为等价的关系；交仍是等价关系，并且仍与左右拼接相容，故给出最小同余。若 \(u\) 满足所有指定等式，则 \(v\sim w\Longleftrightarrow\widehat u(v)=\widehat u(w)\) 是包含 \(R\) 的同余，因而包含 \(\equiv_R\)。由商幺半群的映射性质，\(\widehat u\) 唯一下降。反过来，商上的同态与商映射复合，在每对指定关系上必取相同值；限制到字母上恢复 \(u\)。这两个过程互为逆过程。
\end{proof}

例如在 \(\{a\}^*\) 中施加 \(a^2=a\)，任何正长度的字都等价于 \(a\)，所以商最多有 \(\varepsilon,a\) 两个元素。它们不会进一步相等：把 \(a\) 送到两元素幺半群 \(\{1,e\}\) 中满足 \(e^2=e\ne1\) 的 \(e\)，可区分二者。于是商正好是这一个两元素幺半群。向任意 \(M\) 指定它的同态，相当于在 \(M\) 中挑选一个满足 \(m^2=m\) 的元素。

\subsection{存在性、唯一性与典范同构的区别}
\label{b1:subsec:0104:03}

构造问题中，“存在”“唯一”和“典范”常被放在一句话里，但它们承担不同职责。存在性要求实际给出对象或论证其可取得；唯一性必须说明相对于哪些固定资料；典范性则强调映射由这些资料强制决定，不依赖额外选择。

\begin{theorem}[泛性质的唯一性]\label{b1:thm:universal-uniqueness}
若幺半群 \(F,F'\) 连同映射 \(i:X\rightarrow F\)、\(i':X\rightarrow F'\) 都满足自由幺半群的映射延拓性质，则存在唯一同构 \(\alpha:F\rightarrow F'\) 满足 \(\alpha i=i'\)。
\end{theorem}
\begin{proof}
由 \(F\) 的泛性质，\(i'\) 唯一延拓为同态 \(\alpha:F\rightarrow F'\)；由 \(F'\) 的泛性质，\(i\) 唯一延拓为同态 \(\beta:F'\rightarrow F\)。有
\((\beta\alpha)i=i\)，而恒等同态也在 \(i\) 上取值为 \(i\)。唯一性因此给出
\(\beta\alpha=\operatorname{id}_F\)。交换两个对象得到
\(\alpha\beta=\operatorname{id}_{F'}\)，所以 \(\alpha\) 是同构。任何满足指定等式的同构首先是延拓 \(i'\) 的同态，故只能是 \(\alpha\)。
\end{proof}

这个唯一同构称为由给定资料确定的\term[dianfan-tonggou]{典范同构}[canonical isomorphism]。它不意味着 \(F\) 与 \(F'\) 是同一个集合，也不意味着任意同构都唯一。例如两个字母的自由幺半群有交换两字母的自同构；只有要求分别保持指定字母时，允许的自同构才只剩恒等映射。

\ProseParagraphBreak
同样的双向构造方法适用于商幺半群和交换消去幺半群的群化：先用一方的泛性质构造映射，再用另一方反向构造，最后由唯一性证明复合为恒等。与猜测两个对象的元素对应相比，这种方法往往更短，也更能说明同构为何自然。使用时仍须先证明有关泛性质；把“显然具有泛性质”当作构造的代替，并不能解决良定义或存在性问题。


\begin{chaptersummary}
结合律使有限乘积可以省略括号，单位元处理空乘积，逆元则允许撤销运算。半群与幺半群把这些性质从具体算术中抽离出来，同时保留足够多的例子说明它们的差别。

生成提供表达式，自由对象允许任意指定生成元的像，同余关系规定哪些表达式必须取同一值。处理一个新商构造时，应先检查等价关系，再检查运算下降，最后检查所需映射的良定义、存在性与唯一性。下一章把每个元素均可逆这一条件加入幺半群，得到群。
\end{chaptersummary}

伴读时可参看 Dummit 与 Foote 的《Abstract Algebra》\cite{DummitFoote2003}及 Lang 的《Algebra》\cite{Lang2002}中关于代数结构与映射的讨论。以下习题练习构造新的运算、检验商映射，以及用反例辨别各项公理的作用。

\BookExercises

\begin{exercise}\label{b1:ex:01-affine}
固定实数 \(\alpha,\beta\)，在 \(\mathbb R\) 上定义 \(x*y=\alpha x+\beta y\)。求使运算满足结合律的全部 \((\alpha,\beta)\)，并在这些情形中判断何时有单位元、何时满足交换律。
\end{exercise}
\begin{solution}
展开得
\((x*y)*z=\alpha^2x+\alpha\beta y+\beta z\)，
\(x*(y*z)=\alpha x+\alpha\beta y+\beta^2z\)。对任意实数成立，当且仅当
\(\alpha^2=\alpha\)、\(\beta^2=\beta\)，故二者各为 \(0\) 或 \(1\)。若有双边单位元 \(e\)，从 \(e*y=y\) 比较 \(y\) 的系数得 \(\beta=1\)，从 \(x*e=x\) 得 \(\alpha=1\)，继而 \(e=0\)。交换律要求 \(\alpha=\beta\)，所以结合情形中只有 \((0,0)\) 与 \((1,1)\) 满足。前者处处取零而没有单位元，说明交换与结合仍不足以保证存在单位。
\end{solution}

\begin{exercise}\label{b1:ex:01-odd}
证明奇整数在乘法下构成幺半群，求其中所有有逆的元素。证明其消去律成立，但不能仅靠这些逆元生成整个幺半群。
\end{exercise}
\begin{solution}
奇数的乘积为奇数，\(1\) 是奇数，结合律来自整数乘法。若奇整数 \(a,b\) 满足 \(ab=1\)，则 \(a=b=1\) 或 \(a=b=-1\)，故有逆元素只有 \(\pm1\)。每个奇数非零，所以 \(ax=ay\) 在整数中推出 \(x=y\)，右消去同样成立。但 \(\pm1\) 的任意有限乘积仍为 \(\pm1\)，不能得到 \(3\)。因此一个消去幺半群既可能含有非平凡的可逆元素，也可能包含许多不能由这些可逆元素生成的元素。
\end{solution}

\begin{exercise}\label{b1:ex:01-length-mod}
设非空字母集 \(X\) 与整数 \(m\geq2\) 给定，规定两字等价当且仅当长度模 \(m\) 相同。描述所得商幺半群，并求任一类的逆元。
\end{exercise}
\begin{solution}
长度模 \(m\) 相同是等价关系，因为整数同余具有自反、对称和传递性。又因长度在拼接时相加，两对字分别等价时，它们的拼接也等价，故这是\cref{b1:def:congruence}中的同余关系。由\cref{b1:prop:operation-descends}，拼接能下降为商上的乘法。所有余数 \(0,\ldots,m-1\) 都能由同一个固定字母重复相应次数取得，因此商恰有 \(m\) 个元素，运算为余数相加再取模 \(m\)。长度余数为 \(r\) 的类的逆元是余数为 \(m-r\) 的类，若 \(r=0\) 则为自身。商映射可能把长度正的字送到单位类；这不表示自由幺半群中的原字已有逆元，而是施加关系后产生了新的可逆性。
\end{solution}

\begin{exercise}\label{b1:ex:01-transform-cancel}
设 \(X\) 至少含两个元素，\(f:X\rightarrow X\)。证明：对所有自映射 \(u,v\)，\(fu=fv\Rightarrow u=v\) 当且仅当 \(f\) 单射；\(uf=vf\Rightarrow u=v\) 当且仅当 \(f\) 满射。
\end{exercise}
\begin{solution}
若 \(f\) 单射，\(f\bigl(u(x)\bigr)=f\bigl(v(x)\bigr)\) 逐点推出 \(u(x)=v(x)\)。若不单射，取不同的 \(a,b\) 而 \(f(a)=f(b)\)，以 \(a,b\) 为值的两个常映射给出反例。若 \(f\) 满射，则任意 \(y\in X\) 可写成 \(f(x)\)，从 \(uf=vf\) 得 \(u(y)=v(y)\)。若不满射，取 \(z\notin f(X)\) 及不同的 \(a,b\in X\)，令 \(u\) 恒取 \(a\)，而 \(v\) 只在 \(z\) 处取 \(b\)、其余处取 \(a\)。于是 \(uf=vf\) 却 \(u\ne v\)。左右次序不同，所需的单射、满射条件也不同。
\end{solution}

\begin{exercise}\label{b1:ex:01-free-commutative}
设 \(X\) 为任意集合。令 \(C(X)\) 为所有仅在有限多个点非零的函数 \(a:X\rightarrow\mathbb N\)，运算为逐点相加。证明它具有交换幺半群中的自由映射延拓性质，并解释它与 \(X^*\) 的关系。
\end{exercise}
\begin{solution}
两个有限非零集合的并仍有限，所以加法封闭，零函数为单位，结合律和交换律逐点成立。记 \(\delta_x\) 为在 \(x\) 处取 \(1\)、其余取零的函数。对交换幺半群 \(M\) 和映射 \(u:X\rightarrow M\)，定义
\(\widehat u(a)=\prod_{a(x)\ne0}u(x)^{a(x)}\)。有限性使乘积存在，交换性使取点次序无关；指数相加公式证明它是同态。每个 \(a\) 唯一写成有限和 \(\sum_x a(x)\delta_x\)，故延拓唯一。把字送到各字母出现次数的函数得到满同态 \(X^*\rightarrow C(X)\)；其相等关系就是忽略字母次序、保留字母个数。
\end{solution}

\begin{exercise}\label{b1:ex:01-truncated-addition}
固定 \(n\geq1\)，在 \(T_n=\{0,1,\ldots,n\}\) 上定义 \(a\star b=\min(n,a+b)\)。证明这是幺半群，说明它是 \((\mathbb N,+)\) 的商，并证明任何到所有元素都有逆元的幺半群的同态都取常值单位元。
\end{exercise}
\begin{solution}
等式 \(\min\bigl(n,\min(n,a+b)+c\bigr)=\min(n,a+b+c)\) 证明结合律，零为单位。映射 \(q(k)=\min(n,k)\) 满，且 \(q(k+\ell)=q(k)\star q(\ell)\)、\(q(0)=0\)，故为幺半群同态。其同值关系 \(k\sim\ell\Longleftrightarrow q(k)=q(\ell)\) 与加法相容；由\cref{b1:thm:monoid-quotient}，公式 \([k]\mapsto q(k)\) 给出商到 \(T_n\) 的同态。它由同值关系的定义为单射，由 \(q\) 满射而满射，所以确实是所需同构。设 \(f:T_n\rightarrow G\) 保持单位与乘法，且 \(G\) 中元素都有逆。因为 \(n\star a=n\)，有 \(f(n)f(a)=f(n)\)，左乘 \(f(n)^{-1}\) 得 \(f(a)=1\)。尤其 \(T_n\) 不能单射嵌入这样的 \(G\)。失败源于截断运算丢失了消去律。
\end{solution}

\begin{exercise}\label{b1:ex:01-finite-cancellation}
证明满足左消去律的有限幺半群中，每个元素都有双边逆。给出无限情形的反例。
\end{exercise}
\begin{solution}
固定 \(a\)。左乘映射 \(x\mapsto ax\) 由左消去律为单射，有限性使其满射，故存在 \(b\) 使 \(ab=1\)。又有 \(a(ba)=(ab)a=a=a1\)，再用左消去得 \(ba=1\)。这为每个 \(a\) 提供双边逆。无限幺半群 \((\mathbb N,+)\) 满足左右消去律，但 \(1\) 没有加法逆元；其平移映射 \(x\mapsto1+x\) 单射而不满射。有限性的作用恰是把单射升级为满射。
\end{solution}

\begin{exercise}\label{b1:ex:01-infinitely-generated-submonoid}
在 \((\mathbb N^2,+)\) 中考虑
\(M=\bigl\{(0,0)\bigr\}\cup\bigl\{\,(a,b)\mid a\geq1,\ b\geq0\,\bigr\}\)。证明 \(M\) 是子幺半群但不能由有限多个元素生成，尽管 \(\mathbb N^2\) 可以。
\end{exercise}
\begin{solution}
零元属于 \(M\)；两个非零元素的和第一坐标至少为 \(2\)，仍在 \(M\)，所以它是子幺半群。每个 \((1,b)\) 若写成若干非零元素的和，由第一坐标相加可知只能有一个加数，且这个加数就是 \((1,b)\)。因而任意生成元集合都必须包含无限多个不同的 \((1,b)\)。另一方面，\(\mathbb N^2\) 由 \((1,0),(0,1)\) 生成。可见有限生成性不必传给子幺半群。
\end{solution}

% !TeX root = ../../Book1.tex
% 纲要标题：## 第2章　群、子群与同态
% 纲要位置：计划纲要.md，第 159 行。

\chapter{群、子群与同态}
\label{b1:ch:02}
\BookChapterNumber{b1:ch:02}

% TODO: 撰写本章摘要、学习目标与章导读。

% !TeX root = ../../Book1.tex
% 纲要标题：### 2.1 群与子群
% 纲要位置：计划纲要.md，第 161 行。

\section{群与子群}
\label{b1:sec:0201}

% 纲要内容（仅作写作计划，尚非教材正文）：
%
% * 群的基本恒等式
% * 子群判别、生成子群与循环子群
% * 子群格与交、并的不同性质
%

可逆操作之所以适合描述对称，是因为一次操作完成后，仍能通过另一次同类操作回到原状。整数加法、可逆矩阵的乘法和集合的置换都具有这种特征。群的公理把运算的闭合性、结合性、单位元和可逆性集中列出，使这些例子可以使用同一种推理。

\ProseParagraphBreak
下面先给出完整定义，再研究群内哪些子集仍能独立进行乘法和取逆。辨认子群时，尤其要区分集合的交、集合的并，以及由一个集合生成的子群；它们承担的作用并不相同。

\subsection{群的基本恒等式}
\label{b1:subsec:0201:01}

\begin{definition}\label{b1:def:group}
设 \(G\) 是非空集合，给定乘法 \(G\times G\rightarrow G\)，记 \((a,b)\) 的像为 \(ab\)。若满足以下三条公理，则称 \(G\) 连同此乘法为\term[qun]{群}[group]：
\begin{enumerate}
\item 对任意 \(a,b,c\in G\)，有 \((ab)c=a(bc)\)；
\item 存在 \(1\in G\)，使任意 \(g\in G\) 都满足 \(1g=g1=g\)；
\item 对每个 \(g\in G\)，存在 \(h\in G\)，使 \(gh=hg=1\)。
\end{enumerate}
乘法的值域取为 \(G\)，已包含运算的闭合性。若还对任意 \(a,b\in G\) 有 \(ab=ba\)，则称为\term[abel-qun]{阿贝尔群}[abelian group]或交换群。
\end{definition}
前两条公理给出一个幺半群，第三条要求其中每个元素都有双边逆元。因此，群也可以等价地说成“每个元素都可逆的幺半群”；这个说法有助于比较结构，而上面的公理可以独立用来检验一个例子。

\ProseParagraphBreak
群 \(G\) 的元素个数称为\term[qun-de-jie]{群的阶}[order of a group]，记为 \(\lvert G\rvert\)；元素个数有限时称有限群。
\symidx{\(\lvert G\rvert\)：群的阶}

单位元一般写为 \(1\)，逆元写为 \(g^{-1}\)。交换群使用加法记号时，二者分别写为 \(0\)、\(-g\)。加法记号并不改变定义，只提醒读者此时元素彼此交换。\((\mathbb Z,+)\) 是无限阿贝尔群，非零实数在乘法下是阿贝尔群，\(\operatorname{Sym}(X)\) 在复合下是群。

\begin{proposition}\label{b1:prop:group-identities}
在群中，逆元唯一，且
\[
(ab)^{-1}=b^{-1}a^{-1},\qquad (a^{-1})^{-1}=a.
\]
方程 \(ax=b\) 与 \(ya=b\) 分别有唯一解 \(x=a^{-1}b\)、\(y=ba^{-1}\)。特别地，左右消去律都成立。
\end{proposition}
\begin{proof}
若 \(e,e'\) 都是单位元，则 \(e=ee'=e'\)。若 \(b,c\) 都是 \(a\) 的双边逆元，则 \(b=b(ac)=(ba)c=c\)，故单位元和逆元都唯一。这也是\cref{b1:prop:identity-inverse-unique}在群中的具体应用。计算
\((ab)(b^{-1}a^{-1})=a(bb^{-1})a^{-1}=1\)，反向乘积也为 \(1\)，所以乘积逆元必须是所给式子。\(a\) 本身是 \(a^{-1}\) 的双边逆，故第二式成立。若 \(ax=b\)，左乘 \(a^{-1}\) 强制得到 \(x=a^{-1}b\)，而代入确实给出解；另一方程相同地由右乘求解。
\end{proof}

乘积取逆时次序必须反转。把 \((ab)^{-1}\) 写成 \(a^{-1}b^{-1}\)，一般相当于偷偷使用 \(ab=ba\)。同样，解 \(ax=b\) 时把 \(a^{-1}\) 写在 \(b\) 的右边也会改变结果。

\ProseParagraphBreak
对整数 \(n\) 定义 \(g^n\)：非负指数沿用幺半群中的定义，负指数取
\(g^{-n}=(g^{-1})^n\ (n>0)\)。先消去相邻的 \(gg^{-1}\)，再应用非负指数的归纳公式，得到对任意整数 \(r,s\) 都有
\(g^{r+s}=g^rg^s\)、\((g^r)^s=g^{rs}\)。这里所有因子都是同一元素或其逆，因此彼此交换；这并不能推出一般的 \((ab)^n=a^nb^n\)。

\subsection{子群判别、生成子群与循环子群}
\label{b1:subsec:0201:02}

\begin{definition}\label{b1:def:subgroup}
群 \(G\) 的子集 \(H\) 若包含单位元，且对乘法与取逆封闭，就称为 \(G\) 的\term[ziqun]{子群}[subgroup]，记为 \(H\leq G\)。子群使用 \(G\) 的原运算。
\end{definition}
\symidx{\(H\leq G\)：子群关系}

\begin{proposition}[子群判别]\label{b1:prop:subgroup-test}
非空子集 \(H\subseteq G\) 是子群，当且仅当任取 \(a,b\in H\) 都有 \(ab^{-1}\in H\)。
\end{proposition}
\begin{proof}
子群显然满足该条件。反过来，先取 \(a=b\) 得 \(1\in H\)，再取 \(a=1\) 得 \(b^{-1}\in H\)。既然 \(b^{-1}\) 也属于 \(H\)，把条件用于 \(a,b^{-1}\)，便得 \(ab\in H\)。结合律从 \(G\) 限制到 \(H\)，其他群公理也随之成立。
\end{proof}

非空条件不能删掉：空集使“任取两个元素”的条件空真，却不是群。在整数加法群中，\(m\mathbb Z\coloneqq\{mk\mid k\in\mathbb Z\}\) 是子群，因为两个倍数之差仍为倍数。正整数集合虽然对加法封闭，却不含零元，也不对取负封闭。

\begin{definition}\label{b1:def:generated-subgroup}
包含 \(S\subseteq G\) 的所有子群的交称为 \(S\) 生成的子群，记为 \(\langle S\rangle\)。单个元素 \(g\) 生成的子群称为\term[xunhuan-ziqun]{循环子群}[cyclic subgroup]，记为 \(\langle g\rangle\)。若存在 \(g\in G\)，使每个元素都能写为 \(g\) 的某个整数次幂，即 \(G=\{g^n\mid n\in\mathbb Z\}\)，则称 \(G\) 为\term[xunhuanqun]{循环群}[cyclic group]。下面的命题将说明，这等价于由一个元素生成。
\end{definition}
\symidx{\(\langle S\rangle\)：生成子群}

\begin{proposition}\label{b1:prop:generated-subgroup-words}
\(\langle S\rangle\) 恰由有限乘积 \(s_1^{\varepsilon_1}\cdots s_r^{\varepsilon_r}\) 组成，其中 \(s_i\in S\)、\(\varepsilon_i\in\{1,-1\}\)，并允许空乘积。特别地，\(\langle g\rangle=\{g^n\mid n\in\mathbb Z\}\)。
\end{proposition}
\begin{proof}
任一含 \(S\) 的子群都必须包含这些乘积。反之，这些乘积的集合含单位元，对拼接封闭；取逆时把因子次序反转并把每个指数变号，仍是这样的乘积，所以它是子群。因此它恰是所求的交。
\end{proof}

群生成与幺半群生成的区别在于允许取逆。例如 \(1\) 在 \((\mathbb Z,+)\) 中生成整个群，却只生成幺半群 \(\mathbb N\)。使用尖括号时，本章起默认表示群生成，幺半群生成仍显式带下标 \(\mathrm{mon}\)。

\subsection{子群格与交、并的不同性质}
\label{b1:subsec:0201:03}

子群的包含关系形成 \(G\) 的\term[ziqunge]{子群格}[subgroup lattice]。这里的“格”表达任意两个子群都有最大的公共子群和最小的共同上界：前者是交 \(H\cap K\)，后者是 \(\langle H\cup K\rangle\)。后者通常不是集合的并。

\begin{proposition}\label{b1:prop:subgroup-union}
任意族子群的交仍是子群，空族的交约定为 \(G\)。两个子群 \(H,K\) 的并是子群，当且仅当 \(H\subseteq K\) 或 \(K\subseteq H\)。若非空子群族按包含关系全序排列，则它们的并也是子群。
\end{proposition}
\begin{proof}
交包含单位元，任意两个交中元素的 \(ab^{-1}\) 属于每一个子群，故属于交；由\cref{b1:prop:subgroup-test}，交是子群。对于并，若有包含关系，并就等于较大的那个子群。若二者互不包含，取
\(h\in H\setminus K\)、\(k\in K\setminus H\)。倘若 \(hk\in H\)，则 \(k=h^{-1}(hk)\in H\)，矛盾；倘若 \(hk\in K\)，则 \(h=(hk)k^{-1}\in K\)，亦矛盾，故并不封闭。最后，该子群族非空，其并包含单位元；在按包含关系全序排列的子群族中，任取并中的 \(a,b\)，各自所属的两个子群中总有一个包含另一个，于是二者同属一个子群，\(ab^{-1}\) 也在并中。再次应用\cref{b1:prop:subgroup-test}便得其并为子群。
\end{proof}

在 \(\mathbb Z\) 中，\(2\mathbb Z\cup3\mathbb Z\) 包含 \(2,3\)，却不包含 \(2+3=5\)；而它们生成的子群含 \(3-2=1\)，所以是整个 \(\mathbb Z\)。交则是 \(6\mathbb Z\)。这一例子同时展示了交、并和由并生成的子群三者的不同。

\ProseParagraphBreak
当两个子群相互交换或有正规性条件时，可以用较简单的集合乘积描述它们生成的子群。但在当前阶段，\(\langle H\cup K\rangle\) 的一般元素仍须允许任意多次交替使用 \(H\) 与 \(K\) 的元素；不能未经检验就把它缩成一个 \(hk\)。

% !TeX root = ../../Book1.tex
% 纲要标题：### 2.2 陪集与指数
% 纲要位置：计划纲要.md，第 167 行。

\section{陪集与指数}
\label{b1:sec:0202}

% TODO: 按以下纲要要点撰写本节。

% 纲要内容（仅作写作计划，尚非教材正文）：
%
% * 左陪集、右陪集与等价关系
% * Lagrange 定理及其局限
% * 群同态的核与像
%

待撰写。

% !TeX root = ../../Book1.tex
% 纲要标题：### 2.3 正规子群与商群
% 纲要位置：计划纲要.md，第 173 行。

\section{正规子群与商群}
\label{b1:sec:0203}

% TODO: 按以下纲要要点撰写本节。

% 纲要内容（仅作写作计划，尚非教材正文）：
%
% * 正规性及商群构造
% * 商群的普遍性质
% * 单群概念与简单例子
%

\subsection{正规性及商群构造}
\label{b1:subsec:0203:01}

待撰写。

\subsection{商群的普遍性质}
\label{b1:subsec:0203:02}

待撰写。

\subsection{单群概念与简单例子}
\label{b1:subsec:0203:03}

待撰写。

% !TeX root = ../../Book1.tex
% 纲要标题：### 2.4 同构定理
% 纲要位置：计划纲要.md，第 179 行。

\section{同构定理}
\label{b1:sec:0204}

% 纲要内容（仅作写作计划，尚非教材正文）：
%
% * 第一、第二、第三同构定理
% * 子群对应定理
% * 同构定理的图解与计算
%
% ---
%

同构定理把几个常见操作联系起来：把映射的核商掉、把商映射限制到子群、以及连续取两次商。所得公式虽然相似，背后的映射却各不相同。计算时先写清一个元素怎样映射，再求核和像，通常比记忆商式可靠。

\ProseParagraphBreak
对于商群中的子群，还需要知道它在原群中对应什么。子群对应定理把这个问题统一处理；本节最后用交换图和整数加法群的例子核对各个映射的方向。

\subsection{三个同构定理}
\label{b1:subsec:0204:01}

\term[diyi-tonggou-dingli]{第一同构定理}[first isomorphism theorem]说明任意同态的像就是定义域商去核之后的结构。
\begin{theorem}[第一同构定理]\label{b1:thm:first-isomorphism}
对于同态 \(f:G\rightarrow K\)，映射
\[
G/\ker f\longrightarrow\operatorname{im}f,\qquad
g\ker f\longmapsto f(g)
\]
是同构。
\end{theorem}

\begin{proof}
由\cref{b1:prop:normal-basic-examples}，\(\ker f\trianglelefteq G\)。把 \(f\) 的目标限制为子群 \(\operatorname{im}f\)，核仍为 \(\ker f\)。在\cref{b1:thm:quotient-universal}中取 \(N=\ker f\)，其核包含条件成立，故 \(g\ker f\mapsto f(g)\) 良定义且为同态。它按像的定义为满射；若 \(f(g)=f(h)\)，则 \(h^{-1}g\in\ker f\)，故
\(g\ker f=h\ker f\)，得到单射。群同态也是幺半群同态，故\cref{b1:prop:monoid-maps}保证其双射逆映射也为同态，因此它是同构。
\end{proof}

\term[dier-tonggou-dingli]{第二同构定理}[second isomorphism theorem]比较一个子群在商映射下的像与该子群自身的商。
\begin{theorem}[第二同构定理]\label{b1:thm:second-isomorphism}
设 \(H\leq G\)、\(N\trianglelefteq G\)。则 \(HN=\{hn\mid h\in H,n\in N\}\) 是子群，\(H\cap N\trianglelefteq H\)，且存在同构
\[
H/(H\cap N)\longrightarrow HN/N,\qquad
h(H\cap N)\longmapsto hN.
\]
\end{theorem}

\begin{proof}
集合 \(HN\) 含单位元。对 \(h_1n_1,h_2n_2\in HN\)，有
\[
(h_1n_1)(h_2n_2)=h_1h_2(h_2^{-1}n_1h_2)n_2\in HN,
\]
而 \((hn)^{-1}=n^{-1}h^{-1}=h^{-1}(hn^{-1}h^{-1})\in HN\)，故是子群。\(N\) 在 \(HN\) 中仍正规。把商映射限制到 \(H\)，得同态 \(h\mapsto hN\)，其核是 \(H\cap N\)，像是 \(HN/N\)，因为 \(hnN=hN\)。由\cref{b1:prop:normal-basic-examples}，这个限制同态的核 \(H\cap N\) 正规于 \(H\)；再对该同态应用\cref{b1:thm:first-isomorphism}，便得所述同构及其元素公式。
\end{proof}

这里也回答了\secref{b1:sec:0201}末尾关于生成子群的问题。\(HN\) 是同时包含 \(H,N\) 的子群，而任一同时包含二者的子群都含有全部乘积 \(hn\)，所以
\[
HN=\langle H\cup N\rangle.
\]
由 \(N\) 的正规性，对每个 \(h\in H\) 有 \(hN=Nh\)，取并即得 \(HN=NH\)。这里是两个集合相等，不要求 \(h\) 与 \(n\) 逐元素交换。

\ProseParagraphBreak
\term[disan-tonggou-dingli]{第三同构定理}[third isomorphism theorem]处理先后两次取商。
\begin{theorem}[第三同构定理]\label{b1:thm:third-isomorphism}
若 \(N\leq K\) 且 \(N,K\trianglelefteq G\)，则 \(K/N\trianglelefteq G/N\)，且
\[
(G/N)/(K/N)\longrightarrow G/K,\qquad
(gN)(K/N)\longmapsto gK
\]
是同构。
\end{theorem}

\begin{proof}
同态 \(G\rightarrow G/K\) 的核为 \(K\)，包含 \(N\)，故\cref{b1:thm:quotient-universal}诱导出满同态
\(G/N\rightarrow G/K\)，\(gN\mapsto gK\)。它的核恰为
\(\{gN\mid g\in K\}=K/N\)。由\cref{b1:prop:normal-basic-examples}，\(K/N\) 正规于这个映射的定义域 \(G/N\)；再对该满同态应用\cref{b1:thm:first-isomorphism}，即得 \((G/N)/(K/N)\cong G/K\) 及所述元素公式。特别地，两层陪集公式的良定义已经由这两个具体映射保证，无须在嵌套括号中猜测代表元。
\end{proof}

\subsection{子群对应定理}
\label{b1:subsec:0204:02}

给定正规子群 \(N\)，不仅可以形成一个商群，也可以同时描述商群的全部子群。相应的\term[ziqun-duiying-dingli]{子群对应定理}[subgroup correspondence theorem]要求原群一侧的子群必须包含 \(N\)，这项条件不可省略。

\begin{theorem}[子群对应]\label{b1:thm:subgroup-correspondence}
设 \(N\trianglelefteq G\)，\(q:G\rightarrow G/N\) 是商映射。映射
\[
H\longmapsto H/N=q(H),\qquad U\longmapsto q^{-1}(U)
\]
在含 \(N\) 的子群 \(H\leq G\) 与 \(G/N\) 的子群 \(U\) 之间给出互逆的、保持包含关系的对应。此外：
\begin{enumerate}
\item \(H\trianglelefteq G\) 当且仅当 \(H/N\trianglelefteq G/N\)；
\item 对 \(N\leq K\leq H\leq G\)，有 \([H:K]=[H/N:K/N]\)；
\item 对应保持任意交，以及非空子群族所生成的子群。空族的交在两侧分别为 \(G\) 与 \(G/N\)；商群中由空族生成的单位子群 \(\{N\}\)，对应于原群中的 \(N\)。
\end{enumerate}

\end{theorem}

\begin{proof}
把\cref{b1:prop:kernel-fibers}用于限制同态 \(q|_H\)，可知 \(q(H)\) 是子群。同态下子群的原像含单位元，并对 \(ab^{-1}\) 封闭，故由\cref{b1:prop:subgroup-test}也是子群。原像必含 \(\ker q=N\)。若 \(H\supseteq N\)，则
\(q^{-1}\bigl(q(H)\bigr)=H\)：确实，\(q(g)=q(h)\) 蕴含 \(h^{-1}g\in N\subseteq H\)，从而 \(g\in H\)。另一方面，\(q\) 满射给出
\(q\bigl(q^{-1}(U)\bigr)=U\)。这证明互逆；包含保持由像与原像的基本性质得到。

若 \(H\) 正规，则任意商群元素 \(gN\) 对 \(H/N\) 的共轭仍落在其中。反之，若 \(H/N\) 正规，则对 \(h\in H,g\in G\)，有
\(q(ghg^{-1})\in H/N\)，取原像得 \(ghg^{-1}\in H\)，故 \(H\) 正规。指数公式来自左陪集的双射
\(hK\mapsto(hN)(K/N)\)：由\cref{b1:prop:coset-equivalence}，左边的相等条件为 \(h_2^{-1}h_1\in K\)，右边的相等条件为 \(q(h_2^{-1}h_1)\in q(K)\)；因 \(q^{-1}\bigl(q(K)\bigr)=K\)，二者等价。每个右边的陪集都有这种代表，故该映射确为双射。

交与生成的保持可以用原像公式统一说明。对 \(G/N\) 的任意子群族 \((U_i)_{i\in I}\)，有
\[
q^{-1}\biggl(\bigcap_{i\in I}U_i\biggr)
=\bigcap_{i\in I}q^{-1}(U_i).
\]
当 \(I\) 为空时，这正是 \(q^{-1}(G/N)=G\)。对生成子群，则应写成
\[
q^{-1}\biggl(\Bigl\langle\bigcup_{i\in I}U_i\Bigr\rangle\biggr)
=\Bigl\langle N\cup\bigcup_{i\in I}q^{-1}(U_i)\Bigr\rangle.
\]
为验证第二式，记右侧为 \(L\)。左侧是包含 \(N\) 和全部 \(q^{-1}(U_i)\) 的子群，因而包含 \(L\)。反过来，\(q(L)\) 包含所有 \(U_i\)，故包含它们生成的子群；取原像并用 \(q^{-1}\bigl(q(L)\bigr)=L\)，即得反向包含。若 \(I\) 非空，每个 \(q^{-1}(U_i)\) 已包含 \(N\)，右侧可以省去额外写出的 \(N\)。若 \(I\) 为空，等式两侧都是 \(N\)，而不是 \(\{1\}\)。这证明了第三项。
\end{proof}

特别地，对包含 \(N\) 的两个子群 \(H_1,H_2\)，上述结论给出
\[
\begin{aligned}
q(H_1\cap H_2)&=q(H_1)\cap q(H_2),\\
q\Bigl(\bigl\langle H_1\cup H_2\bigr\rangle\Bigr)
&=\bigl\langle q(H_1)\cup q(H_2)\bigr\rangle.
\end{aligned}
\]
第一式中包含核的条件很具体：若 \(q(h_1)=q(h_2)\)，其中 \(h_i\in H_i\)，则 \(h_2^{-1}h_1\in N\subseteq H_2\)，故 \(h_1\in H_2\)。所以这个公共的像能由 \(H_1\cap H_2\) 中的同一个元素取得。任意映射的直接像并不自动保持交；这里用到了两子群都包含核。

\ProseParagraphBreak
例如要找 \(\mathbb Z/12\mathbb Z\) 中由 \([4]\) 生成的子群，先在 \(\mathbb Z\) 中取包含 \(12\mathbb Z\) 与 \(4\) 的子群 \(4\mathbb Z\)。它的像是 \(\bigl\{[0],[4],[8]\bigr\}\)，指数为 \([\mathbb Z:4\mathbb Z]=4\)。若原子群不含核，例如 \(5\mathbb Z\)，其像的原像会变大：模 \(12\) 时 \([5]\) 生成全部余数，所以原像是 \(\mathbb Z\)，不是 \(5\mathbb Z\)。

\subsection{同构定理的图解与计算}
\label{b1:subsec:0204:03}

第一同构定理可排列成下面的交换图。图中两条从 \(G\) 到 \(\operatorname{im}f\) 的路径给出同一个映射；诱导映射 \(\overline f\) 由等式 \(\overline f q=f\) 唯一确定。
\[
\begin{tikzcd}
G \arrow[r,"f"] \arrow[d,"q"'] & \operatorname{im}f \\
G/\ker f \arrow[ur,"\overline f"',"\cong"] &
\end{tikzcd}
\]
这个图解释“典范”的含义：在商映射 \(q\) 与原同态 \(f\) 已经固定后，\(\overline f\) 没有选择余地。但若只知道两个抽象群同构而没有指定 \(f\)，并不能由此得到唯一同构。

\ProseParagraphBreak
第二同构定理的一个可跟算例子是加法群 \(G=\mathbb Z\)，取
\(H=4\mathbb Z\)、\(N=6\mathbb Z\)。有
\(H\cap N=12\mathbb Z\)，且
\(H+N=2\mathbb Z\)：任意和都是偶数，而 \(6-4=2\) 生成全部偶数。限制模 \(6\) 的映射到 \(4\mathbb Z\)，得到
\[
4\mathbb Z/12\mathbb Z\longrightarrow2\mathbb Z/6\mathbb Z,
\qquad4k+12\mathbb Z\longmapsto4k+6\mathbb Z.
\]
左边三个类的代表可取 \(0,4,8\)，它们分别送到右边的 \(0,4,2\) 类，确实一一对应。这个同构并不是把“分子与分母同时约去一个数”；它来自指定的模 \(6\) 映射。

\ProseParagraphBreak
第三同构定理中，\(12\mathbb Z\subseteq4\mathbb Z\subseteq\mathbb Z\) 给出满同态
\[
\mathbb Z/12\mathbb Z\longrightarrow\mathbb Z/4\mathbb Z,
\qquad a+12\mathbb Z\longmapsto a+4\mathbb Z.
\]
核是三个元素组成的 \(4\mathbb Z/12\mathbb Z\)，再把这三个元素组成的子群商掉，便得到四元素群。与前例并列时尤其应注意：两处虽然都出现 \(4,12\)，映射的定义域、核和目标不同，不能只凭商式外观套用名称。


% TODO: 撰写本章小结、习题与逐题详解。

% !TeX root = ../../Book1.tex
% 纲要标题：## 第3章　循环群、置换群与几何对称
% 纲要位置：计划纲要.md，第 187 行。

\chapter{循环群、置换群与几何对称}
\label{b1:ch:03}
\BookChapterNumber{b1:ch:03}

% TODO: 撰写本章摘要、学习目标与章导读。

% !TeX root = ../../Book1.tex
% 纲要标题：### 3.1 循环群
% 纲要位置：计划纲要.md，第 192 行。

\section{循环群}
\label{b1:sec:0301}

% 纲要内容（仅作写作计划，尚非教材正文）：
%
% * 有限及无限循环群的分类
% * 子群、商群与同态的显式计算
% * 循环群的自同构群
%

循环群中的每个元素都是一个固定元素的整数次幂，因此群内的等式可以转化为指数之间的整除关系。从整数加法群的子群出发，第一同构定理同时给出有限循环群和无限循环群的模型。

\ProseParagraphBreak
分类完成以后，还要能实际计算：哪些幂生成同一个子群，商群有多大，生成元可以映到哪里。下面始终区分“选择一个生成元”和“这个生成元是否由群本身唯一确定”。

\subsection{循环群的分类}
\label{b1:subsec:0301:01}

\begin{lemma}\label{b1:lem:integer-subgroups}
\(\mathbb Z\) 的每个加法子群都唯一写成 \(n\mathbb Z\)，其中 \(n\geq0\) 为整数。约定 \(0\mathbb Z=\{0\}\)。
\end{lemma}
\begin{proof}
零子群对应 \(n=0\)。非零子群含一个正整数，取其中最小者 \(n\)。任给子群中的 \(a\)，除法算法给出 \(a=qn+r\)，\(0\leq r<n\)。因
\(r=a-qn\) 仍在子群，最小性迫使 \(r=0\)。故子群包含于 \(n\mathbb Z\)，反向包含由 \(n\) 属于子群得出。正生成元是该子群最小的正元素，因而唯一。
\end{proof}

\begin{theorem}[循环群分类]\label{b1:thm:cyclic-classification}
每个无限循环群都同构于 \((\mathbb Z,+)\)；每个 \(n\) 阶有限循环群都同构于 \((\mathbb Z/n\mathbb Z,+)\)。因此两个循环群同构当且仅当它们的阶相同，包括无限阶情形。
\end{theorem}
\begin{proof}
设 \(G=\langle g\rangle\)。映射 \(\pi:\mathbb Z\rightarrow G\)，\(k\mapsto g^k\)，是满同态。核是 \(\mathbb Z\) 的子群，故\cref{b1:lem:integer-subgroups}给出唯一 \(n\geq0\)，使 \(\ker\pi=n\mathbb Z\)。对这个像为全体 \(G\) 的同态应用\cref{b1:thm:first-isomorphism}，得到同构 \(\mathbb Z/n\mathbb Z\rightarrow G\)，\([k]\mapsto g^k\)。若 \(n=0\)，商就是 \(\mathbb Z\)；若 \(n>0\)，\(0,\ldots,n-1\) 为商的全部不同代表，所以群阶为 \(n\)。同构保持元素个数，因而也得到最后的判别。
\end{proof}

记 \(n\) 阶循环群的一个标准模型为 \(C_n\coloneqq\mathbb Z/n\mathbb Z\)，其中 \(n\geq1\)，\(C_1\) 是平凡群。需要强调乘法表示时，也写 \(C_n=\langle g\rangle\)、\(g^n=1\)，同时声明 \(g\) 的阶恰为 \(n\)。仅写 \(g^n=1\) 还不够，因为阶可能是 \(n\) 的真因子。
\symidx{\(C_n\)：\(n\) 阶循环群}

\ProseParagraphBreak
上述证明给出的同构依赖生成元 \(g\) 的选择。若把生成元换成 \(g^a\)，由整数指数得到的同构一般也改变。这与上一章典范同构的讨论一致：有时群的同构类型唯一，但某个具体同构并不由群本身唯一确定。

\subsection{子群、商群与同态的计算}
\label{b1:subsec:0301:02}

\begin{proposition}\label{b1:prop:power-order}
设 \(G\) 为群，\(g\in G\)。若 \(g\) 的有限阶为正整数 \(n\)，则对任意 \(a\in\mathbb Z\)，有
\[
\operatorname{ord}(g^a)=\frac{n}{\gcd(n,a)}.
\]
这里 \(\gcd(n,a)\) 取正值，特别地，\(\gcd(n,0)=n\)，所以公式也包括 \(g^0=1\) 的情形。若 \(g\) 为无限阶且 \(a\in\mathbb Z\setminus\{0\}\)，则 \(g^a\) 也为无限阶。
\end{proposition}
\begin{proof}
有限情形中，\((g^a)^k=g^{ak}\)，由\eqref{b1:eq:element-order-divisibility}，它等于 \(1\) 当且仅当 \(n\mid ak\)。写
\(n=dn'\)、\(a=da'\)，其中 \(d=\gcd(n,a)\)、\(\gcd(n',a')=1\)。整除条件等价于 \(n'\mid a'k\)，再由互素性等价于 \(n'\mid k\)，所以最小正指数为 \(n'\)。无限情形中若正整数 \(k\) 满足 \(g^{ak}=1\)，则 \(ak\ne0\)，与无限阶矛盾。
\end{proof}

互素时的整除步骤可直接用 Bézout 等式验证：若 \(un'+va'=1\)，则
\(k=un'k+va'k\)，右边两项都被 \(n'\) 整除。这样的整数等式可由 Euclid 除法反复回代获得，本节只使用这一初等算术事实。

\begin{theorem}\label{b1:thm:cyclic-subgroups}
设 \(G\) 为循环群。则 \(G\) 的每个子群以及关于任意正规子群的商群均为循环群。若 \(G=\langle g\rangle\) 的有限阶为正整数 \(n\)，则对每个正因子 \(d\mid n\)，恰有一个 \(d\) 阶子群，即 \(\langle g^{n/d}\rangle\)，\(G\) 关于这个子群的商群阶为 \(n/d\)。
\end{theorem}
\begin{proof}
对于子群 \(H\leq G\)，取满同态 \(\pi:\mathbb Z\rightarrow G\)，\(k\mapsto g^k\)。原像 \(\pi^{-1}(H)\) 对整数之差封闭且含零，因而是加法子群。\cref{b1:lem:integer-subgroups}给出 \(\pi^{-1}(H)=m\mathbb Z\)；由 \(\pi\) 满射，\(\pi\bigl(\pi^{-1}(H)\bigr)=H\)，故 \(H=\langle g^m\rangle\)，所以循环。有限情形中原像包含 \(n\mathbb Z\)，于是 \(m\mid n\)，由\cref{b1:prop:power-order}，\(\langle g^m\rangle\) 的阶为 \(n/\gcd(n,m)=n/m\)，由阶唯一确定 \(m\)。群 \(G\) 交换，所以每个子群正规；商群的每个元素都是 \(gH\) 的幂，故商群循环。在有限情形中，其阶由\cref{b1:thm:lagrange}给出，即 \(\lvert G/H\rvert=\lvert G\rvert/\lvert H\rvert\)。
\end{proof}

例如 \(C_{18}=\langle g\rangle\) 中 \(g^6\) 的阶为 \(3\)，\(g^4\) 的阶为 \(9\)，且
\(\langle g^4\rangle=\langle g^2\rangle\)。商去 \(\langle g^6\rangle\) 得到六阶循环群。此时“子群唯一”是指给定阶的子群唯一，并非生成元唯一：同一个三阶子群由 \(g^6\) 或 \(g^{12}\) 生成。

\ProseParagraphBreak
同态同样可以显式计算。
\begin{proposition}[循环群之间的同态]\label{b1:prop:cyclic-homomorphisms}
设 \(m,n\) 为正整数，\(C_m=\langle x\rangle\)、\(C_n=\langle y\rangle\)。任一同态由 \(x\) 的像 \(y^a\) 唯一决定；指定这个像能够给出同态，当且仅当 \(n\mid ma\)。若 \(d=\gcd(m,n)\)，则模 \(n\) 的全部可取指数为
\[
a=0,n/d,2n/d,\ldots,(d-1)n/d,
\]
所以同态共有 \(d\) 个。
\end{proposition}
\begin{proof}
同态必须把 \(x^k\) 送到 \(y^{ak}\)，故至多有一个。因为 \(x^m=1\)，其存在要求 \(y^{am}=1\)，由\eqref{b1:eq:element-order-divisibility}等价于 \(n\mid ma\)。反过来，若此整除条件成立，则 \(x^k=x^\ell\) 给出 \(m\mid k-\ell\)，进而 \(n\mid a(k-\ell)\)，所以 \(y^{ak}=y^{a\ell}\)。公式 \(x^k\mapsto y^{ak}\) 因而良定义；指数相加立即给出同态性。写 \(m=dm'\)、\(n=dn'\)，则 \(n\mid ma\) 等价于 \(n'\mid m'a\)。因 \(m',n'\) 互素，Bézout 等式使它等价于 \(n'\mid a\)，给出所列的 \(d\) 个模 \(n\) 指数。
\end{proof}

比如 \(C_6\rightarrow C_9\) 只有三种，由生成元分别送到 \(1,y^3,y^6\) 给出；其中两种非平凡同态的像都为同一个三阶子群。

\subsection{循环群的自同构群}
\label{b1:subsec:0301:03}

\begin{definition}\label{b1:def:automorphism}
设 \(G\) 为群。将群同态 \(G\rightarrow G\) 称为 \(G\) 的\term[zitongtai]{自同态}[endomorphism]，将群同构 \(G\rightarrow G\) 称为 \(G\) 的\term[zitonggou]{自同构}[automorphism]。所有自同构在复合下构成\term[zitonggou-qun]{自同构群}[automorphism group]，记为 \(\operatorname{Aut}(G)\)。
\end{definition}
\symidx{\(\operatorname{Aut}(G)\)：自同构群}

复合封闭来自同态复合与双射复合；单位元是恒等映射；逆映射仍为同构。这说明自同构本身也有可研究的代数结构，它描述了群内部允许的重新命名方式。

\begin{theorem}[循环群的自同构]\label{b1:thm:cyclic-automorphisms}
设 \(n\geq2\) 为整数，\(C_n=\langle g\rangle\)。对每个满足 \(\gcd(a,n)=1\) 的整数 \(a\)，令
\[
\alpha_a:C_n\longrightarrow C_n,\qquad
\alpha_a(g^k)=g^{ak}\quad(k\in\mathbb Z).
\]
映射 \(\alpha_a\) 只依赖 \(a\) 的模 \(n\) 剩余类，且这些映射恰为 \(C_n\) 的全部自同构。复合公式为 \(\alpha_a\alpha_b=\alpha_{ab}\)，所以
\(\operatorname{Aut}(C_n)\) 同构于模 \(n\) 的可逆剩余类在乘法下组成的群，记为 \((\mathbb Z/n\mathbb Z)^\times\)。
\end{theorem}
\symidx{\((\mathbb Z/n\mathbb Z)^\times\)：可逆剩余类群}
\begin{proof}
在\cref{b1:prop:cyclic-homomorphisms}中取定义域与目标的阶都为 \(n\)，得到全部自同态 \(g^k\mapsto g^{ak}\)，其中任意模 \(n\) 指数 \(a\) 都满足存在条件 \(n\mid na\)。它满射当且仅当 \(g^a\) 仍生成
\(C_n\)，由\cref{b1:prop:power-order}，\(\operatorname{ord}(g^a)=n/\gcd(n,a)\)，故这当且仅当 \(\gcd(a,n)=1\)；有限集合上满射即双射。若互素，Bézout 等式给出 \(ab\equiv1\pmod n\)，所以 \(\alpha_b\) 是逆映射。反之，剩余类有乘法逆时 \(ab-1\) 被 \(n\) 整除，任何公因子都整除 \(1\)，故互素。逐点计算
\(\alpha_a\bigl(\alpha_b(g^k)\bigr)=g^{abk}\)，得到复合公式及所述群同构。
\end{proof}

定理的证明为了计算而选定了生成元 \(g\)，但所得自同构就是 \(C_n\) 上的幂映射 \(x\mapsto x^a\)，所以映射本身不依赖生成元的选择。事实上，若改取生成元 \(h=g^u\)，则
\(\alpha_a(h^k)=\alpha_a(g^{uk})=g^{auk}=h^{ak}\)。因此，可逆剩余类与幂自同构之间的对应不依赖生成元；与之不同，前面由生成元 \(g\) 给出的模型同构 \([k]\mapsto g^k\) 一般会随 \(g\) 改变。

\ProseParagraphBreak
对正整数 \(n\)，将模 \(n\) 的可逆剩余类的个数记为 \(\varphi(n)\)，并称 \(\varphi\) 为\term[euler-hanshu]{Euler 函数}[Euler totient function]。例如 \(n=8\) 时可取 \(1,3,5,7\)，它们的平方模 \(8\) 都为 \(1\)。于是 \(\operatorname{Aut}(C_8)\) 有四个元素，但不是循环群，因为没有四阶元素。一个循环群的自同构群不一定循环。

\ProseParagraphBreak
无限循环群只有两个自同构：在 \(\mathbb Z\) 的加法模型中，同态必为
\(k\mapsto ak\)，满射当且仅当 \(a=1\) 或 \(a=-1\)。因此
\(\operatorname{Aut}(\mathbb Z)\) 是二阶群。有限与无限情形的共同原则是“生成元必须送到生成元”，差别只在于整数指数何时仍能生成全部元素。

% !TeX root = ../../Book1.tex
% 纲要标题：### 3.2 对称群与交错群
% 纲要位置：计划纲要.md，第 195 行。

\section{对称群与交错群}
\label{b1:sec:0302}

% TODO: 按以下纲要要点撰写本节。

% 纲要内容（仅作写作计划，尚非教材正文）：
%
% * 轮换分解与置换的阶
% * 换位、符号同态与交错群
% * 共轭类型与分拆
%

\subsection{轮换分解与置换的阶}
\label{b1:subsec:0302:01}

待撰写。

\subsection{换位、符号同态与交错群}
\label{b1:subsec:0302:02}

待撰写。

\subsection{共轭类型与分拆}
\label{b1:subsec:0302:03}

待撰写。

% !TeX root = ../../Book1.tex
% 纲要标题：### 3.3 二面体群与四元数群
% 纲要位置：计划纲要.md，第 201 行。

\section{二面体群与四元数群}
\label{b1:sec:0303}

% TODO: 按以下纲要要点撰写本节。

% 纲要内容（仅作写作计划，尚非教材正文）：
%
% * 二面体群的表示、子群与中心
% * 四元数群与非交换小群
% * 同阶非同构群的比较
%

\subsection{二面体群的表示、子群与中心}
\label{b1:subsec:0303:01}

待撰写。

\subsection{四元数群与非交换小群}
\label{b1:subsec:0303:02}

待撰写。

\subsection{同阶非同构群的比较}
\label{b1:subsec:0303:03}

待撰写。

% !TeX root = ../../Book1.tex
% 纲要标题：### 3.4 正多面体群
% 纲要位置：计划纲要.md，第 207 行。

\section{正多面体群}
\label{b1:sec:0304}

% 纲要内容（仅作写作计划，尚非教材正文）：
%
% * 正四面体、立方体与正八面体
% * 正二十面体与正十二面体的构造及对称
% * 顶点、边、面上的置换作用
% * 旋转群与含反射的全对称群的区别
%
% ---
%

几何图形的对称必须连同所在空间一起说明。正多边形在平面中的旋转与反射已经给出循环群和二面体群；进入三维以后，顶点、边、面和体对角线都可以作为被置换的对象。同一个空间旋转在这些集合上可能具有不同的置换符号。

\ProseParagraphBreak
以下把多面体中心放在原点，以保持图形的正交线性变换为对称。行列式为 \(1\) 的称为旋转，行列式为 \(-1\) 的反转空间定向。确定一个图形的对称群时，我们依次构造对称、证明没有其他对称，再通过具体同态识别群结构；本节的计数不预先借用下一章的群作用定理。

\begin{lemma}[三维旋转的固定轴]\label{b1:lem:rotation-axis}
设 \(T\) 是三阶实正交矩阵，且 \(\det T=1\)。若 \(T\ne I\)，则
\(\dim\ker(T-I)=1\)。这条固定直线的正交补在 \(T\) 下保持不变，\(T\) 在该平面上是一个非恒等的平面旋转。
\end{lemma}
\begin{proof}
由 \(T^{-1}=T^{\mathsf T}\) 及行列式的乘法性质，
\[
\begin{aligned}
\det(T-I)
&=\det T\det(I-T^{-1})
 =\det(I-T^{\mathsf T})\\
&=\det(I-T)=-\det(T-I).
\end{aligned}
\]
所以 \(\det(T-I)=0\)，存在非零固定向量 \(v\)。若 \(x\perp v\)，则
\(\langle Tx,v\rangle=\langle Tx,Tv\rangle=\langle x,v\rangle=0\)，故 \(v^\perp\) 保持不变。
在以 \(v/\lVert v\rVert\) 为首向量的正交基下，\(T\) 分成一个 \(1\) 与一个行列式为 \(1\) 的二阶正交块。
后一块的第一列可写成 \((\cos\theta,\sin\theta)^{\mathsf T}\)，第二列由正交性及定向唯一确定，故它是平面旋转矩阵。

若另有与 \(v\) 线性无关的固定向量，减去它在 \(v\) 上的投影，就得到 \(v^\perp\) 内的非零固定向量。
平面旋转固定一个非零向量便是恒等变换，于是整个 \(T=I\)，与假设矛盾。因此固定空间恰为一维，平面上的旋转也非恒等。
\end{proof}

为了从顶点坐标确定多面体，先说明本节所需的几个几何概念。有限点集 \(V=\{v_1,\ldots,v_m\}\) 的\term[tubao]{凸包}[convex hull]是所有凸组合
\[
\operatorname{conv}V
=\left\{\,\sum_{i=1}^m t_i v_i\ \middle|\ t_i\geq0,\ \sum_{i=1}^m t_i=1\,\right\}.
\]
若所有顶点满足 \(\langle a,v_i\rangle\leq b\)，其中 \(a\ne0\)，且至少一个顶点取等号，
则平面 \(\langle a,x\rangle=b\) 称为凸包的\term[zhichi-pingmian]{支持平面}[supporting plane]。
凸组合取等号时只能使用取等号的顶点，所以凸包与该平面的交恰是这些顶点的凸包。
交集若为二维多边形，便是多面体的一个面；例如恰有三个不共线顶点取等号时，交集是三角面。
指向 \(\langle a,x\rangle>b\) 一侧的法向称为该面的\term[wai-faxiang]{外法向}[outward normal]。

\ProseParagraphBreak
下面还会用截面判断面是否已经列全。用一个平面在某顶点附近截下一个小角，并让其余顶点都位于截面的另一侧，截面就是一个凸多边形。
经过原顶点的边给出截面的顶点，经过原顶点的面给出截面的边：在每个面内，截面切出连接相邻两条边的线段。
凸多边形等于各边所在直线的内侧闭半平面之交。因此，若几个已经找到的面在截面上围成一个多边形，截面既包含这个多边形，又被这些面的内侧半平面限制在其中，两者便相等。
原顶点与每条截面边唯一确定对应面的平面，故截面没有额外的边，就说明没有遗漏经过原顶点的面。

\subsection{正四面体、立方体与正八面体}
\label{b1:subsec:0304:01}

有限多面体的任意欧氏对称会置换顶点，所以保持所有顶点的平均值，也就是这里取作原点的中心。故把中心移到原点后，只需研究正交线性变换。它们的复合仍正交、逆仍正交，保持给定图形的那些变换形成其\term[jihe-duichengqun]{几何对称群}[symmetry group]，其中的旋转形成\term[xuanzhuanqun]{旋转群}[rotation group]。

\begin{proposition}[正四面体]\label{b1:prop:tetrahedral-group}
正四面体的全对称群在四个顶点上的置换给出与 \(S_4\) 的同构；旋转群对应 \(A_4\)。
\end{proposition}
\begin{proof}
取顶点
\begin{equation}\label{b1:eq:tetrahedron-vertices}
v_1=(1,1,1),\quad v_2=(1,-1,-1),\quad
v_3=(-1,1,-1),\quad v_4=(-1,-1,1).
\end{equation}
它们等长，任意两个不同顶点的内积为 \(-1\)，故构成正四面体。前三个向量线性无关，且 \(v_1+v_2+v_3+v_4=0\)。任意顶点排列都唯一延拓为线性变换：先规定前三个向量的像，和为零的关系自动规定第四个的像。由于基向量的两两内积不变，这个线性变换正交。因此每个顶点排列都实现为唯一几何对称。将对称送到其顶点排列的映射还保持复合，因为 \((TU)v_i=T(Uv_i)\)，故这个双射是群同构。

为判断定向，把四列 \((1,v_i)\) 排成一个四阶矩阵，其行列式非零。顶点排列 \(\sigma\) 使这个行列式乘以 \(\operatorname{sgn}(\sigma)\)；相应线性变换 \(T\) 则使它乘以 \(\det T\)。所以
\(\det T=\operatorname{sgn}(\sigma)\)，旋转恰对应偶置换。
\end{proof}

\begin{proposition}[立方体与正八面体]\label{b1:prop:cubic-rotation-group}
立方体与正八面体的旋转群同构于 \(S_4\)，各有 \(24\) 个元素；各自全对称群有 \(48\) 个元素。
\end{proposition}
\begin{proof}
取立方体顶点为 \((\pm1,\pm1,\pm1)\)。对称置换六个面的外法向
\(\pm e_1,\pm e_2,\pm e_3\)，所以矩阵的每列、每行恰有一个非零元，非零元为 \(\pm1\)。反过来，每个这样的带符号坐标置换都保持立方体。共有 \(3!\cdot2^3=48\) 个，其中行列式正负各半，故旋转共 \(24\) 个。同一组六个法向本身构成正八面体的顶点集，因此两个图形有相同的正交对称群。

立方体的四条体对角线可由\eqref{b1:eq:tetrahedron-vertices}中的 \(v_1,v_2,v_3,v_4\) 表示，每条线忽略向量的正负。旋转置换这四条线，给出到 \(S_4\) 的同态。若旋转 \(T\) 固定每条线，则
\(Tv_i=\varepsilon_i v_i\)，\(\varepsilon_i\in\{1,-1\}\)。这四向量的唯一线性关系由和为零生成，所以
\(\sum_i\varepsilon_i v_i=0\) 迫使四个 \(\varepsilon_i\) 相同。故
\(T=I\) 或 \(T=-I\)；后一矩阵在三维中行列式为 \(-1\)，不是旋转。核因而平凡，\cref{b1:prop:kernel-fibers}保证这个同态单射；两边都有 \(24\) 个元素，所以是同构。
\end{proof}

\subsection{正二十面体与正十二面体}
\label{b1:subsec:0304:icosahedron-dodecahedron}

正二十面体也可以从坐标构造。这里先确认顶点、边和面的完整结构，再计算对称群；正十二面体则由这些面的中心构造。令
\(\phi=(1+\sqrt5)/2\)，取十二个顶点
\[
V=\bigl\{(0,\pm1,\pm\phi),\ (\pm1,\pm\phi,0),\
(\pm\phi,0,\pm1)\bigr\}.
\]
记 \(P=\operatorname{conv}V\)。下面验证它是正二十面体。
各顶点的长度的平方为 \(\phi+2\)，最短的顶点间距离为 \(2\)。以
\(v=(0,1,\phi)\) 为例，距离为 \(2\) 的五个点依次记为 \(w_1,\ldots,w_5\)：
\[
(0,-1,\phi),\ (\phi,0,1),\ (1,\phi,0),\
(-1,\phi,0),\ (-\phi,0,1).
\]
它们与 \(v\) 的内积都为 \(\phi\)，在垂直于 \(v\) 的同一平面中构成正五边形：相邻点距离为 \(2\)，非相邻点距离为 \(2\phi\)。
以下下标按模 \(5\) 计算。五个三角形 \(\operatorname{conv}\{v,w_i,w_{i+1}\}\) 都是等边三角形。
对于其中任意一个，记其三个顶点为 \(u,v,w\)，则
\[
\langle u+v+w,u\rangle
=\langle u+v+w,v\rangle
=\langle u+v+w,w\rangle=3\phi+2.
\]
不同顶点间的内积至多为 \(\phi\)，故对其余顶点 \(z\)，
\(\langle u+v+w,z\rangle\leq3\phi\)。因此每个候选三角形确实是 \(P\) 的面。

\ProseParagraphBreak
为排除遗漏，在 \(v\) 附近取截面
\[
\langle v,x\rangle=\phi+2-\varepsilon,
\qquad 0<\varepsilon<2.
\]
除 \(v\) 外的所有顶点在此平面的另一侧。它与 \([v,w_i]\) 的交点为
\[
p_i=\left(1-\frac{\varepsilon}{2}\right)v
     +\frac{\varepsilon}{2}w_i.
\]
五个已证三角面与截面相交成五条线段 \([p_i,p_{i+1}]\)，它们围成一个正五边形。
截面包含这五个点的凸包，又位于这五条边的各自内侧，所以恰是这个正五边形。
经过 \(v\) 的边和面因而各有五个，正是已经列出的那些。

\ProseParagraphBreak
坐标循环置换和任意两个坐标同时变号保持 \(V\)，并能把 \(v\) 送到任意顶点，所以每个顶点都有同样的邻接结构。
每条边有两个端点，每个三角面有三个顶点，故
\[
E=\frac{12\cdot5}{2}=30,
\qquad F=\frac{12\cdot5}{3}=20.
\]
这些面全是全等的等边三角形，每个顶点恰有五个面相交，得到正二十面体。

\begin{proposition}[正二十面体与正十二面体]\label{b1:prop:icosahedral-order}
正二十面体与正十二面体各有 \(60\) 个旋转对称，\(120\) 个全对称。
\end{proposition}
\begin{proof}
坐标循环置换与任意两个坐标同时变号都是行列式为 \(1\) 的对称，它们能把上面的 \(v\) 送到十二个顶点中的任意一个。固定 \(v\) 的旋转必须置换其五个邻点，并在垂直于 \(v\) 的平面中保持定向，故只能是该正五边形的五个旋转之一。这五个旋转确实保持整个顶点集：它们固定 \(v,-v\)，置换五个邻点，也置换这五点的相反点，而这些点恰好构成全部十二个顶点。

对每个目标顶点 \(w\)，固定一个把 \(v\) 送到 \(w\) 的旋转 \(t_w\)。所有把 \(v\) 送到 \(w\) 的旋转恰为 \(t_w\) 与那五个固定 \(v\) 的旋转的复合：左乘 \(t_w^{-1}\) 即可验证这一一对应。所以总数为 \(12\cdot5=60\)。中心反演 \(-I\) 保持顶点集且反转定向，与它复合建立两类定向之间的双射，故全对称群有 \(120\) 个元素。

再构造正十二面体。对每个三角面 \(F\) 的三个顶点 \(u,w,z\)，令
\[
c_F=\frac{u+w+z}{3},\qquad
Q=\operatorname{conv}\{\,c_F\mid F\text{ 是 }P\text{ 的三角面}\,\}.
\]
这里 \(c_F\) 是等边三角面的中心。所有这些中心满足
\(\lVert c_F\rVert^2=\phi+\frac23\)。这些面中心组成的点集关于原点对称。
固定原顶点 \(v\)，若 \(v\in F\)，则
\[
\langle v,c_F\rangle
=\frac{(\phi+2)+\phi+\phi}{3}=\phi+\frac23;
\]
若 \(v\notin F\)，则 \(\langle v,c_F\rangle\leq\phi\)。特别地，若 \(F\ne F'\)，取 \(v\in F\setminus F'\)，这两种内积取值便说明 \(c_F\ne c_{F'}\)。
因此，经过 \(v\) 的五个三角面的中心恰在 \(Q\) 的支持平面
\(\langle v,x\rangle=\phi+\frac23\) 上。
绕 \(v\) 轴的五次旋转循环置换这五点，所以它们围成正五边形面，记为 \(Q_v\)。
坐标对称把各个 \(v\) 互相变换，故这十二个五边形全等。
一个五边形及其相反面不在同一平面内，故 \(Q\) 是三维凸多面体。

还须证明 \(Q\) 没有其他面。若原来的边 \([u,w]\) 连接两三角面 \(F,F'\)，则同时满足
\(\langle u,x\rangle=\langle w,x\rangle=\phi+\frac23\)
的面中心恰为 \(c_F,c_{F'}\)。因此
\(Q_u\cap Q_w=[c_F,c_{F'}]\)，这是两个五边形共有的边。
在一个固定的 \(c_F\) 处，原三角面 \(F\) 的三条边给出这样的三条边，它们两两同属 \(F\) 的一个顶点所对应的五边形。

由于各面中心等长且两两不同，线性函数 \(x\mapsto\langle c_F,x\rangle\) 在这些点中仅于 \(c_F\) 取最大值，故每个 \(c_F\) 确为 \(Q\) 的顶点。
取 \(\eta>0\) 小于所有正数
\(\lVert c_F\rVert^2-\langle c_F,c_{F'}\rangle\)，其中 \(F'\ne F\)，用
\(\langle c_F,x\rangle=\lVert c_F\rVert^2-\eta\) 作截面。
上述三条边给出三个截点。每一对截点分别位于同一个正五边形在 \(c_F\) 处的两条不同邻边上，故两两不同。
它们也不共线：否则这条直线同时位于两个五边形的平面内，只能是它们的公共边所在直线；但那条直线经过 \(c_F\)，与截面只有一个交点。
所以三个已证五边形在截面上给出一个三角形的三条边。
截面包含这个三角形，又被这三条边的内侧半平面限制在其中，故恰是三角形。
所以每个 \(c_F\) 处只有这三个面；任何面都含有顶点，十二个 \(Q_v\) 因而已经列全。
新多面体有 \(20\) 个顶点、\(20\cdot3/2=30\) 条边及 \(12\) 个全等正五边形面，即正十二面体。

保持 \(P\) 的正交变换置换其三角面，也就置换面中心并保持 \(Q\)。
反过来，保持 \(Q\) 的正交变换置换已列全的十二个面，其单位外法向恰为
\(v/\sqrt{\phi+2}\)，其中 \(v\in V\)。因此它也保持 \(V\)，从而保持 \(P\)。
两图形的正交对称群相同，旋转群也相同，所述群阶随之得到。
\end{proof}

这里直接求出了两种群的阶，并没有把阶相同当作同构证明。识别正二十面体旋转群与 \(A_5\) 的进一步关系，需要另选合适的五个几何对象或补充群结构论证，本节不使用这一识别作为前提。

\subsection{顶点、边与面上的置换作用}
\label{b1:subsec:0304:02}

每个几何对称都置换顶点、边和面，且复合对称所得的置换是原置换的复合。因此选择这些有限集合，就得到从几何对称群到某个对称群的同态。这是把连续空间中的几何运动化为有限代数计算的方法；下一章将把这种做法统一称为群作用。

\ProseParagraphBreak
对本节的正多面体，顶点置换能够区分全部对称：固定全部顶点的线性变换固定一组空间基，因而是恒等变换。边置换也能区分对称，因为一个顶点是所有与它相接的边的唯一公共点；若每条边作为集合均被固定，每个顶点也被固定。类似地，一个顶点由与它相接的面的公共交点确定，所以固定每个面的对称也是恒等对称。

\ProseParagraphBreak
同一对称在不同集合上可以有很不一样的轮换分解。例如立方体绕 \(z\) 轴逆时针转过 \(\pi/2\)，映射为
\(T(x,y,z)=(-y,x,z)\)。在八个顶点上，它分别循环置换上面与下面的四个顶点，轮换类型为 \(4+4\)，因此是偶置换。在六个面上，上下面固定，四个侧面循环，轮换类型为 \(4+1+1\)，因此是奇置换。在四条体对角线上，按\eqref{b1:eq:tetrahedron-vertices}中的 \(v_i\) 编号，它给出 \((1\,2\,4\,3)\)，也是奇置换。

\ProseParagraphBreak
这说明“空间旋转对应偶置换”并非普遍规则。对于四面体顶点它成立，是因为已经证明三维行列式等于该顶点排列的符号；换成另一套被置换对象后，必须重新判断。前面的立方体旋转群同构 \(S_4\)，也正是四条体对角线的置换所给出的，而非八个顶点的全部置换群。

\begin{proposition}\label{b1:prop:cube-lines-kernel}
立方体全对称群在四条体对角线上的置换同态的核为 \(\{I,-I\}\)，而在八个顶点上的置换同态的核为 \(\{I\}\)。
\end{proposition}
\begin{proof}
设 \(T\) 固定四条体对角线。沿用\eqref{b1:eq:tetrahedron-vertices}的四个等长向量，正交性给出 \(Tv_i=\varepsilon_i v_i\)，其中 \(\varepsilon_i\in\{1,-1\}\)。由于 \(v_1,v_2,v_3\) 线性无关且 \(v_4=-v_1-v_2-v_3\)，等式 \(\sum_i Tv_i=0\) 变为
\[
(\varepsilon_1-\varepsilon_4)v_1+
(\varepsilon_2-\varepsilon_4)v_2+
(\varepsilon_3-\varepsilon_4)v_3=0.
\]
所以四个符号相同，\(T=I\) 或 \(T=-I\)。这重用了\cref{b1:prop:cubic-rotation-group}证明中的线性关系，但此处允许反转定向，所以两者都保留；它们也确实固定每条体对角线。若顶点置换恒等，则 \(T\) 固定基向量 \(v_1,v_2,v_3\)，只能是 \(I\)。
\end{proof}

所以顶点、边、面或对角线的选择会改变同态的核。较小的置换集合通常使计算更简洁，但必须检查它是否仍能分辨全部对称，或者明确记录哪些对称被视为相同。

\subsection{旋转群与全对称群}
\label{b1:subsec:0304:03}

行列式是从全对称群到 \(\{1,-1\}\) 的同态，旋转群恰为其核。只要存在一个反转定向的对称，该行列式同态就满射到二元素群；由\cref{b1:thm:first-isomorphism}，全对称群商去旋转群同构于 \(\{1,-1\}\)，故旋转群的指数为 \(2\)。五种正多面体都有这样的对称：四面体由奇顶点排列实现，其余四种由中心反演实现。因此全对称群的元素数都是旋转群的两倍。\cref{b1:tab:regular-polyhedra}汇总了本节的计算。

\begin{table}[htbp]
\centering
\caption[正多面体的对称计数]{正多面体的对称计数。旋转与全对称均指三维空间中的等距对称。}
\label{b1:tab:regular-polyhedra}
\begin{tabular}{lrrrrr}
\toprule
多面体 & 顶点数 & 边数 & 面数 & 旋转群阶 & 全对称群阶 \\
\midrule
正四面体 & 4 & 6 & 4 & 12 & 24 \\
立方体 & 8 & 12 & 6 & 24 & 48 \\
正八面体 & 6 & 12 & 8 & 24 & 48 \\
正十二面体 & 20 & 30 & 12 & 60 & 120 \\
正二十面体 & 12 & 30 & 20 & 60 & 120 \\
\bottomrule
\end{tabular}
\end{table}

其中正四面体、立方体和正八面体的旋转群已分别识别为 \(A_4,S_4,S_4\)；正十二面体与正二十面体在本节完成了对称群阶数的计算及两者之间的对应，旋转群与 \(A_5\) 的同构尚未在这里证明。

\ProseParagraphBreak
中心对称的三维图形还有一个更具体的分解。若 \(R\) 是其旋转群，则任意全对称都唯一写成 \((-I)^\varepsilon r\)，其中 \(r\in R\)、\(\varepsilon\in\{0,1\}\)。存在性由行列式决定是否乘一次 \(-I\)，唯一性来自 \(-I\) 不是旋转；又因 \(-I\) 与所有线性变换交换，乘法为
\[
\bigl((-I)^\varepsilon r\bigr)\bigl((-I)^\delta s\bigr)
=(-I)^{\varepsilon+\delta}(rs).
\]
因此立方体、正八面体、正十二面体、正二十面体的全对称群都可看作一个旋转及一个二值定向标记组成的群，乘法逐项进行。后面讨论群的直积时将给这种结构统一命名。

\ProseParagraphBreak
四面体没有中心反演对称：\(-v_1\) 不是它的顶点。因此它的全对称群虽也比旋转群大一倍，却不能直接沿用上面的分解。实际上全对称群 \(S_4\) 的中心平凡：若某置换与每个换位交换，\eqref{b1:eq:cycle-conjugation}说明它固定每个二元素子集。对不同的 \(i,j,k\)，单点集 \(\{i\}=\{i,j\}\cap\{i,k\}\) 因而也被固定，故它固定每个点。若四面体全对称群也有那样独立且可交换的二值标记，就会出现一个非平凡中心元素，与此矛盾。相同的阶数比，并不意味着相同的扩张方式。


% TODO: 撰写本章小结、习题与逐题详解。

% !TeX root = ../../Book1.tex
% 纲要标题：## 第4章　群作用与计数
% 纲要位置：计划纲要.md，第 215 行。

\chapter{群作用与计数}
\label{b1:ch:04}
\BookChapterNumber{b1:ch:04}

% TODO: 撰写本章摘要、学习目标与章导读。

% !TeX root = ../../Book1.tex
% 纲要标题：### 4.1 作用、轨道与稳定子
% 纲要位置：计划纲要.md，第 217 行。

\section{作用、轨道与稳定子}
\label{b1:sec:0401}

% TODO: 按以下纲要要点撰写本节。

% 纲要内容（仅作写作计划，尚非教材正文）：
%
% * 群作用与到置换群的同态
% * 轨道—稳定子定理
% * 传递作用与陪集空间
%

\subsection{群作用与到置换群的同态}
\label{b1:subsec:0401:01}

待撰写。

\subsection{轨道—稳定子定理}
\label{b1:subsec:0401:02}

待撰写。

\subsection{传递作用与陪集空间}
\label{b1:subsec:0401:03}

待撰写。

% !TeX root = ../../Book1.tex
% 纲要标题：### 4.2 共轭作用
% 纲要位置：计划纲要.md，第 223 行。

\section{共轭作用}
\label{b1:sec:0402}

上一节允许群作用在任意集合上，现在取这个集合为群自身。左乘把所有元素连成一个轨道，共轭却保留了更细的区别：单位始终不动，其他元素则按共轭类分组。计算这些轨道的大小，可以把群阶分解成若干子群指数之和；对素数幂阶群，这一等式尤其有力。

% 本节正文按下列原纲要要点撰写。

% 纲要内容（仅作写作计划，尚非教材正文）：
%
% * 共轭类、中心化子与正规化子
% * 类方程
% * 有限 p 群的中心
%

\subsection{共轭类、中心化子与正规化子}
\label{b1:subsec:0402:01}

令 \(g\cdot x=gxg^{-1}\)，则 \(1\cdot x=x\)，并且
\[
g\cdot(h\cdot x)=ghx(gh)^{-1}=(gh)\cdot x,
\]
故这确为群作用。

\begin{definition}\label{b1:def:centralizer-normalizer}
设 \(G\) 为群，\(x\in G\)。将 \(x\) 在 \(G\) 的共轭作用下的轨道称为它的\term[gonge-lei]{共轭类}[conjugacy class]，记为
\(\operatorname{Cl}_G(x)=\{\,gxg^{-1}\mid g\in G\,\}\)。
元素 \(x\) 在该作用下的稳定子群称为它的\term[zhongxinhua-zi]{中心化子}[centralizer]，记为
\[
C_G(x)\coloneqq\{\,g\in G\mid gx=xg\,\}.
\]
对任意子集 \(A\subseteq G\)，其中心化子为
\[
C_G(A)\coloneqq\{\,g\in G\mid ga=ag\text{ 对每个 }a\in A\,\}
=\bigcap_{a\in A}C_G(a).
\]
前面定义的中心可以写为 \(Z(G)=C_G(G)\)。
子群 \(H\leq G\) 的\term[zhengguihua-zi]{正规化子}[normalizer]为
\[
N_G(H)\coloneqq\{\,g\in G\mid gHg^{-1}=H\,\}.
\]
\symidx{\(\operatorname{Cl}_G(x)\)：元素的共轭类}
\symidx{\(C_G(x)\)：元素的中心化子}
\symidx{\(C_G(A)\)：子集的中心化子}
\symidx{\(N_G(H)\)：子群的正规化子}
\end{definition}

这里 \(C_G(x)\) 固定的是元素 \(x\)。若考察整个共轭类作为集合是否保持不变，答案则不同：对任意 \(g\in G\)，都有
\[
g\operatorname{Cl}_G(x)g^{-1}
=\Bigl\{\,(gh)x(gh)^{-1}\mid h\in G\,\Bigr\}
=\operatorname{Cl}_G(x),
\]
因为 \(h\mapsto gh\) 是 \(G\) 上的双射。因此，保持这个共轭类整体不变的群元素是 \(G\) 的全部元素，而不只是 \(C_G(x)\) 中的元素。

\ProseParagraphBreak
正规化子对应于另一个作用，其作用集合是 \(G\) 的全部子群组成的集合 \(\mathcal S(G)\)。对 \(g\in G\) 和 \(H\in\mathcal S(G)\)，定义
\[
g\cdot H=gHg^{-1}.
\]
这首先确实给出一个子群：\(1\in gHg^{-1}\)，而对 \(u,v\in H\)，有
\[
(gug^{-1})(gvg^{-1})^{-1}=g(uv^{-1})g^{-1}\in gHg^{-1},
\]
故\cref{b1:prop:subgroup-test}适用。其次，\(1\cdot H=H\)，并且
\[
g\cdot(h\cdot H)=g(hHh^{-1})g^{-1}=(gh)H(gh)^{-1}=(gh)\cdot H.
\]
因此这定义了 \(G\) 在 \(\mathcal S(G)\) 上的作用。点 \(H\) 的稳定子群恰为 \(N_G(H)\)，所以正规化子是子群。由\cref{b1:thm:orbit-stabilizer}的陪集与轨道双射，还得到
\[
\bigl|\{\,gHg^{-1}\mid g\in G\,\}\bigr|=[G:N_G(H)].
\]
此处不要求 \(G\) 有限；等号表示共轭子群集合与相应左陪集集合具有相同基数。

\ProseParagraphBreak
中心化子 \(C_G(H)\) 要求共轭固定 \(H\) 的每个元素，正规化子 \(N_G(H)\) 则只要求共轭保持 \(H\) 整体不变。每个 \(h\in H\) 都满足 \(hHh^{-1}=H\)，故 \(H\leq N_G(H)\)；再按正规化子的定义，便有 \(H\trianglelefteq N_G(H)\)。此外，\(H\trianglelefteq G\) 当且仅当 \(N_G(H)=G\)。中心也是正规子群：\(Z(G)=\bigcap_{x\in G}C_G(x)\) 是子群，而且若 \(z\in Z(G)\)，则 \(gzg^{-1}=z\)，故每个共轭仍在中心中。

\ProseParagraphBreak
以 \(S_3\) 为例，\(C_{S_3}\Bigl((12)\Bigr)=\Bigl\{1,(12)\Bigr\}\)，因为与 \((12)\) 交换的置换必须保持其唯一固定点 \(3\)，并只能在 \(1,2\) 之间置换。因此换位的共轭类有三个元素。三个换位组成的集合被整个 \(S_3\) 保持，而其中单个换位的中心化子只有两个元素。另一方面，\(H=\Bigl\langle(123)\Bigr\rangle\) 的中心化子为 \(H\)，而正规化子是整个 \(S_3\)；换位把 \((123)\) 变成其逆，保持子群，却不逐点固定子群。

\subsection{类方程}
\label{b1:subsec:0402:02}

对共轭作用使用\cref{b1:thm:orbit-stabilizer}，其中点 \(x\) 的稳定子群为 \(C_G(x)\)，得到
\[
\bigl|\operatorname{Cl}_G(x)\bigr|=\bigl[G:C_G(x)\bigr].
\]
此式对任意群都成立，因为该定理给出了两个集合之间的双射。共轭类大小为 \(1\) 当且仅当 \(x\in Z(G)\)。以下设 \(G\) 有限，把这些单元素轨道先分离出来，再把其他共轭类的大小相加，就得到\term[lei-fangcheng]{类方程}[class equation]。

\begin{proposition}[类方程]\label{b1:prop:class-equation}
设 \(G\) 为有限群，从每个非中心共轭类中选一个代表 \(x_1,\ldots,x_r\)。则
\begin{equation}\label{b1:eq:class-equation}
|G|=\bigl|Z(G)\bigr|+\sum_{i=1}^{r}\bigl[G:C_G(x_i)\bigr].
\end{equation}
\end{proposition}
\begin{proof}
共轭类把 \(G\) 分割成不交并。中心元素各自构成一个单元素类。对每个非中心代表 \(x_i\)，\cref{b1:thm:orbit-stabilizer}给出其共轭类大小为 \(\bigl[G:C_G(x_i)\bigr]\)。对这些不交的类计数，就得到\eqref{b1:eq:class-equation}。
\end{proof}

例如 \(S_3\) 的类方程为 \(6=1+3+2\)：单位、三个换位、两个三轮换恰好给出三个共轭类。此式中的 \(1\) 表明 \(Z(S_3)=\{1\}\)。类方程中的求和按共轭类取代表，而不是对每个非中心元素求和；否则同一个轨道会重复计算。

\ProseParagraphBreak
对 \(D_8=\langle r,s\rangle\)，其中 \(r\) 是正方形四分之一周旋转，\(s\) 是反射，有 \(srs^{-1}=r^{-1}\)。旋转 \(r^a\) 与 \(s\) 交换当且仅当 \(r^a=r^{-a}\)，即 \(a\) 为偶数；反射 \(r^as\) 不与 \(r\) 交换。因此 \(Z(D_8)=\{1,r^2\}\)。由同一关系可得
\[
sr^as^{-1}=r^{-a},\qquad
r^b(r^as)r^{-b}=r^{a+2b}s,\qquad
s(r^as)s^{-1}=r^{-a}s.
\]
旋转彼此交换，而用 \(s\) 共轭将旋转变成其逆，所以非中心旋转恰组成 \(\{r,r^3\}\)。对反射，用生成元 \(r,s\) 共轭都保持指数的奇偶性，故任意群元素的共轭也保持这一奇偶性；用 \(r\) 共轭又将指数加 \(2\)，从而遍历同一奇偶类。于是全部共轭类为
\[
\{1\},\quad \{r^2\},\quad \{r,r^3\},\quad
\{s,r^2s\},\quad \{rs,r^3s\},
\]
类方程为 \(8=2+2+2+2\)。它既反映中心，也区分了正方形两种几何位置不同的反射。

\subsection{有限 \texorpdfstring{\(p\)}{p} 群的中心}
\label{b1:subsec:0402:03}

\begin{definition}\label{b1:def:p-group}
设 \(p\) 为素数，\(a\geq0\) 为整数。如果有限群 \(G\) 的阶为 \(p^a\)，则将 \(G\) 称为有限 \(p\)-群。当 \(a\geq1\) 时，\(G\) 非平凡。
\end{definition}

由\cref{b1:thm:orbit-stabilizer}，有限 \term[p-qun]{\(p\)-群}[p-group] \(P\) 的每个轨道大小为某个稳定子群的指数 \([P:P_x]\)，因而整除 \(|P|\)，是 \(p\) 的幂。这给出一个将在 Sylow 理论中反复使用的计数方法。

\ProseParagraphBreak
对于任意群 \(G\) 在集合 \(X\) 上的作用，把被所有群元素固定的点称为\term[gongtong-budongdian]{共同不动点}[common fixed point]，这些点组成的集合记为
\[
X^G\coloneqq\{\,x\in X\mid gx=x\text{ 对每个 }g\in G\,\}.
\]
\symidx{\(X^G\)：群作用的共同不动点集}
共同不动点恰是单元素轨道中的点。作用群为 \(P\) 时，这个集合相应地记为 \(X^P\)。
\symidx{\(X^P\)：\(p\)-群作用的共同不动点集}

\begin{lemma}\label{b1:lem:p-action-fixed}
设 \(p\) 为素数。若有限 \(p\)-群 \(P\) 作用在有限集合 \(X\) 上，记
\(X^P=\{\,x\in X\mid gx=x\text{ 对所有 }g\in P\,\}\)，则
\[
|X|\equiv |X^P|\pmod p.
\]
\end{lemma}
\begin{proof}
大小为 \(1\) 的轨道恰由 \(X^P\) 中的点组成；其他轨道大小是大于 \(1\) 的 \(p\) 幂，均被 \(p\) 整除。把轨道大小相加再模 \(p\) 即得结论。
\end{proof}
\begin{theorem}\label{b1:thm:p-group-center}
设 \(p\) 为素数，\(P\) 为非平凡有限 \(p\)-群。则 \(Z(P)\) 非平凡，而且 \(|Z(P)|\) 被 \(p\) 整除。
\end{theorem}
\begin{proof}
记所给非平凡有限 \(p\)-群为 \(P\)，令它作用在有限集合 \(X=P\) 上，作用为 \(g\cdot x=gxg^{-1}\)。一个元素 \(x\) 被所有 \(g\in P\) 固定，当且仅当它与所有 \(g\) 交换，所以 \(X^P=Z(P)\)。对这个作用应用\cref{b1:lem:p-action-fixed}，其结论 \(|X^P|\equiv|X|\pmod p\) 在这里成为
\[
\bigl|Z(P)\bigr|\equiv|P|\equiv0\pmod p.
\]
末一个同余使用 \(P\neq\{1\}\)，即 \(|P|=p^a\)、\(a\geq1\)。中心包含单位，故其阶是正的 \(p\) 倍数，至少为 \(p\)，从而中心非平凡。
\end{proof}

这里“不动点存在”并非仅由作用非空得到，而是由同余和中心含单位共同得到。一般有限群的中心可以平凡，如 \(S_3\) 所示。

\ProseParagraphBreak
下面先记录一个不要求有限性的事实，供后续讨论中心商时使用。

\begin{proposition}\label{b1:prop:cyclic-central-quotient}
若群 \(G\) 的商群 \(G/Z(G)\) 是循环群，则 \(G\) 是交换群。
\end{proposition}
\begin{proof}
取一个生成陪集 \(aZ(G)\)。对每个 \(g\in G\)，存在整数 \(i\) 使 \(gZ(G)=a^iZ(G)\)，故 \(g=a^iz\)，其中 \(z\in Z(G)\)。任取两个元素 \(a^iz,a^jw\)，利用 \(z,w\) 与所有元素交换，得到
\[
(a^iz)(a^jw)=a^{i+j}zw=a^{j+i}wz=(a^jw)(a^iz).
\]
因此 \(G\) 中任意两个元素交换。
\end{proof}

\begin{corollary}\label{b1:cor:order-p-square-abelian}
设 \(p\) 为素数。每个阶为 \(p^2\) 的群都是交换群。
\end{corollary}
\begin{proof}
由\cref{b1:thm:p-group-center}，\(Z(G)\) 非平凡；又由\cref{b1:thm:lagrange}，其阶整除 \(p^2\)，所以只能为 \(p\) 或 \(p^2\)。若为 \(p^2\)，则 \(Z(G)=G\)，结论成立。若为 \(p\)，商群 \(G/Z(G)\) 的阶为 \(p\)，由\cref{b1:cor:finite-order-divides}中关于素数阶群的结论，它是循环群。再用\cref{b1:prop:cyclic-central-quotient}便知 \(G\) 交换；事实上这也排除了中心阶恰为 \(p\) 的可能。
\end{proof}

中心非平凡并不意味着整个 \(p\)-群交换；阶为 \(8\) 的 \(D_8\) 和 \(Q_8\) 已提供反例。上述推论利用的是商群阶恰为素数，不能把这一论证无条件移到 \(p^3\) 阶。

% !TeX root = ../../Book1.tex
% 纲要标题：### 4.3 轨道计数
% 纲要位置：计划纲要.md，第 229 行。

\section{轨道计数}
\label{b1:sec:0403}

当若干对象在对称变换下视为相同，所求数量就是轨道数。不能一概用对象总数除以群阶，因为对称性较强的对象有较小的轨道。直接列出每个轨道往往困难，固定一个变换后计算它保持不动的对象却容易得多。对这些不动点数取平均，将给出正确的计数公式。

% 本节正文按下列原纲要要点撰写。

% 纲要内容（仅作写作计划，尚非教材正文）：
%
% * Burnside 轨道计数引理
% * 染色与有限组合计数
% * 作用的核与忠实性
%

\subsection{Burnside 轨道计数引理}
\label{b1:subsec:0403:01}

下面的\term[burnside-guidao-jishu-yinli]{Burnside 轨道计数引理}[Burnside orbit-counting lemma]把轨道数与各个群元素的不动点数联系起来。

\begin{theorem}[Burnside 轨道计数]\label{b1:thm:burnside-count}
设有限群 \(G\) 作用在有限集合 \(X\) 上。记轨道集合为 \(X/G\)，并记
\[
X^g\coloneqq\{\,x\in X\mid gx=x\,\}.
\]
则
\[
|X/G|=\frac{1}{|G|}\sum_{g\in G}|X^g|.
\]
\symidx{\(X^g\)：一个群元素的不动点集}
\symidx{\(X/G\)：作用的轨道集合}
\end{theorem}
\begin{proof}
考察集合 \(E=\bigl\{\,(g,x)\in G\times X\mid gx=x\,\bigr\}\)。先固定 \(g\)，得到
\(|E|=\sum_g|X^g|\)；先固定 \(x\)，得到
\(|E|=\sum_x|G_x|\)。
对同一轨道 \(O\) 中的每个 \(x\)，\cref{b1:thm:orbit-stabilizer}给出 \(|G_x|=|G|/|Gx|=|G|/|O|\)，所以该轨道对后一和式的贡献为
\[
\sum_{x\in O}|G_x|=|O|\frac{|G|}{|O|}=|G|.
\]
每个轨道贡献一次 \(|G|\)，故 \(|E|=|X/G|\,|G|\)，除以 \(|G|\) 即得公式。
\end{proof}
公式取的是不动点数的平均值；即使某些轨道较小，也不需要另外为它们加权。直接把对象总数除以群阶只有在所有稳定子均平凡时才正确。一个完全对称的对象可能被很多变换固定，双重计数正好补偿这种轨道大小的差异。

\subsection{染色与有限组合计数}
\label{b1:subsec:0403:02}

给正方形四个顶点涂上 \(k\) 种有名称的颜色，允许重复使用颜色。把顶点集合记为 \(V\)，颜色集合记为 \(C\)，全部染色组成 \(X=C^V\)，即从 \(V\) 到 \(C\) 的函数集合。若 \(G\) 是所允许的对称群，正确的左作用是
\[
(g\cdot f)(v)=f(g^{-1}v).
\]
逆元使复合方向相容：\(g\cdot(h\cdot f)\) 在 \(v\) 的值为
\(f(h^{-1}g^{-1}v)=f\bigl((gh)^{-1}v\bigr)\)。染色 \(f\) 被 \(g\) 固定，当且仅当它在 \(g\) 的每个顶点轮换上取常值。因此，若 \(g\) 在 \(V\) 上有 \(c(g)\) 个轮换（包括不动顶点），则 \(|X^g|=k^{c(g)}\)。

\ProseParagraphBreak
若只允许旋转，群为 \(C_4\)。单位有四个轮换，两个四分之一周旋转各有一个轮换，半周旋转有两个轮换。把相应的不动点数 \(k^4,k,k,k^2\) 代入\cref{b1:thm:burnside-count}并除以群阶 \(4\)，得到染色轨道数
\[
\frac{k^4+k^2+2k}{4}.
\]
例如 \(k=2\) 时得到 \(6\) 类。颜色名称不随对称变换交换，所以全红与全蓝是两个不同的轨道。

\ProseParagraphBreak
若允许全部正方形对称，群为 \(D_8\)。两条对角线反射各有两个固定顶点和一个二轮换，因此有三个轮换；两条对边中点连线反射各有两个二轮换。在四个旋转的不动点总数之外再加上 \(2k^3+2k^2\)，\cref{b1:thm:burnside-count}要求这次除以 \(|D_8|=8\)，所以计数变为
\[
\frac{k^4+2k^3+3k^2+2k}{8}.
\]
当 \(k=3\) 时，旋转等价下有 \(24\) 类，全部对称等价下有 \(21\) 类。两种答案的不同来自等价关系不同，不是同一个计数问题的矛盾答案。

\ProseParagraphBreak
有额外限制时，应先把 \(X\) 换成满足限制的染色集合，再检查该集合是否被群保持。例如要求相邻顶点异色，这个条件在 \(D_8\) 下保持，但此时不能直接使用 \(k^{c(g)}\)：有的轮换会迫使相邻顶点同色，固定染色数便可能为零。轨道计数引理仍成立，需要修改的是逐个变换的不动点计算。

\subsection{作用的核与忠实性}
\label{b1:subsec:0403:03}

\begin{definition}\label{b1:def:faithful-action}
设群 \(G\) 作用在集合 \(X\) 上。使 \(X\) 中每一点都不动的群元素组成作用的核，记为
\[
K\coloneqq\{\,g\in G\mid gx=x\text{ 对所有 }x\in X\,\}.
\]
若只有单位元素固定全部的点，即 \(K=\{1\}\)，称此作用为\term[zhongshi-zuoyong]{忠实作用}[faithful action]。若对任意 \(x\in X\) 和 \(g\in G\)，等式 \(gx=x\) 都蕴含 \(g=1\)，称此作用为\term[ziyou-zuoyong]{自由作用}[free action]。
\end{definition}
由定义，\(K\) 等于对应同态 \(\rho:G\rightarrow S_X\) 的核，也等于所有稳定子的交 \(\bigcap_{x\in X}G_x\)。\cref{b1:prop:kernel-fibers}说明同态单射当且仅当核平凡，因此忠实等价于 \(\rho\) 单射；自由则按定义等价于每个稳定子都平凡。

\ProseParagraphBreak
自由作用一定忠实（这里假设 \(X\neq\varnothing\)），反过来不成立：\(S_3\) 在三个点上的自然作用忠实，但换位仍固定一个点。忠实性要求一个元素只要固定所有点就必须为单位；自由性要求一个元素只要固定任一点就必须为单位，二者的量词不同。

\begin{proposition}\label{b1:prop:quotient-action}
若 \(K\) 是作用的核，则公式 \((gK)x=gx\) 定义 \(G/K\) 在 \(X\) 上的忠实作用，且与原作用具有相同的轨道。
\end{proposition}
\begin{proof}
首先，\(K=\ker\rho\) 是同态的核，故由\cref{b1:prop:normal-basic-examples}正规，\cref{b1:thm:quotient-group}保证 \(G/K\) 确为群。若 \(gK=hK\)，写 \(g=hk\)、\(k\in K\)，则 \(gx=h(kx)=hx\)，所以定义与代表元无关。两条作用公理从原作用继承。若 \(gK\) 固定每个点，则 \(g\in K\)，所以这个陪集是单位。最后，任一 \(g\) 与其陪集对 \(X\) 的变换相同，故轨道不变。
\end{proof}
例如 \(D_8\) 对两条对角线的作用有核 \(\{1,r^2,s,r^2s\}\)，这里 \(s\) 取沿一条对角线的反射。商群 \(C_2\) 忠实地交换两条对角线。这个两点作用不能区分八个对称，却足以发现一个指数为 \(2\) 的正规子群。Burnside 公式不要求忠实性；若按商群求和，每个商元素的不动点数在原群中恰重复 \(|K|\) 次，平均值保持不变。

% !TeX root = ../../Book1.tex
% 纲要标题：### 4.4 作用的结构应用
% 纲要位置：计划纲要.md，第 235 行。

\section{作用的结构应用}
\label{b1:sec:0404}

% TODO: 按以下纲要要点撰写本节。

% 纲要内容（仅作写作计划，尚非教材正文）：
%
% * Cayley 定理
% * 低指数子群与正规性
% * 通过不同作用比较群的结构
%
% ---
%

\subsection{Cayley 定理}
\label{b1:subsec:0404:01}

待撰写。

\subsection{低指数子群与正规性}
\label{b1:subsec:0404:02}

待撰写。

\subsection{群作用与结构比较}
\label{b1:subsec:0404:03}

待撰写。


% TODO: 撰写本章小结、习题与逐题详解。

% !TeX root = ../../Book1.tex
% 纲要标题：## 第5章　Sylow 理论与有限群
% 纲要位置：计划纲要.md，第 243 行。

\chapter{Sylow 理论与有限群}
\label{b1:ch:05}
\BookChapterNumber{b1:ch:05}

% TODO: 撰写本章摘要、学习目标与章导读。

% !TeX root = ../../Book1.tex
% 纲要标题：### 5.1 Cauchy 定理与存在性
% 纲要位置：计划纲要.md，第 245 行。

\section{Cauchy 定理与存在性}
\label{b1:sec:0501}

% TODO: 按以下纲要要点撰写本节。

% 纲要内容（仅作写作计划，尚非教材正文）：
%
% * 素数阶元素的存在
% * Sylow p 子群
% * Sylow 第一定理
%

\subsection{素数阶元素的存在}
\label{b1:subsec:0501:01}

待撰写。

\subsection{Sylow \texorpdfstring{\(p\)}{p} 子群}
\label{b1:subsec:0501:02}

待撰写。

\subsection{Sylow 第一定理}
\label{b1:subsec:0501:03}

待撰写。

% !TeX root = ../../Book1.tex
% 纲要标题：### 5.2 共轭性与计数
% 纲要位置：计划纲要.md，第 251 行。

\section{共轭性与计数}
\label{b1:sec:0502}

% TODO: 按以下纲要要点撰写本节。

% 纲要内容（仅作写作计划，尚非教材正文）：
%
% * Sylow 子群的共轭性
% * Sylow 子群个数的同余和整除限制
% * p 子群与正规化子的关系
%

待撰写。

% !TeX root = ../../Book1.tex
% 纲要标题：### 5.3 小阶群的分析
% 纲要位置：计划纲要.md，第 257 行。

\section{小阶群的分析}
\label{b1:sec:0503}

% TODO: 按以下纲要要点撰写本节。

% 纲要内容（仅作写作计划，尚非教材正文）：
%
% * pq 阶群及适当条件下的分类
% * 平方自由阶的若干现象
% * 用 Sylow 约束排除单群
%

\subsection{两素数乘积阶群的分类}
\label{b1:subsec:0503:01}

待撰写。

\subsection{平方自由阶群}
\label{b1:subsec:0503:02}

待撰写。

\subsection{Sylow 约束与单群排除}
\label{b1:subsec:0503:03}

待撰写。

% !TeX root = ../../Book1.tex
% 纲要标题：### 5.4 综合方法
% 纲要位置：计划纲要.md，第 263 行。

\section{综合方法}
\label{b1:sec:0504}

% TODO: 按以下纲要要点撰写本节。

% 纲要内容（仅作写作计划，尚非教材正文）：
%
% * 类方程与 Sylow 理论的配合
% * 置换作用与单群检验
% * 不把 Sylow 定理误当作任意有限群的分类算法
%
% ---
%

待撰写。


% TODO: 撰写本章小结、习题与逐题详解。

% !TeX root = ../../Book1.tex
% 纲要标题：## 第6章　直积、半直积与群扩张
% 纲要位置：计划纲要.md，第 275 行。

\chapter{直积、半直积与群扩张}
\label{b1:ch:06}
\BookChapterNumber{b1:ch:06}

\begin{chapterabstract}
知道两个群后，怎样构造一个以它们为组成部分的新群？直积使两个因子独立运算，半直积加入一个自同构作用，群扩张则允许陪集代表的乘法产生额外因子。本章建立直积、直和与半直积的构造及其映射性质，再用截面和因子集描述一般群扩张，并通过二面体群、仿射群、四元数群和循环群说明分裂与非分裂的区别。
\end{chapterabstract}

\begin{learninggoals}
\begin{itemize}
\item 区分任意直积、有限支集直积和交换群直和的泛性质。
\item 验证外半直积的群公理，运用内半直积判别计算具体群。
\item 从短正合群列识别核、商及截面，说明扩张等价要求保持哪些映射，证明分裂与半直积分解等价。
\item 计算简单扩张的因子集，辨认依赖截面选择的数据和不变的结构。
\end{itemize}
\end{learninggoals}

本章需要用到\chapref{b1:ch:02}的正规子群和商群、\chapref{b1:ch:03}的自同构与具体有限群，以及\chapref{b1:ch:04}的共轭作用。半直积本身不依赖 Sylow 理论；\chapref{b1:ch:05}的 \(pq\) 阶群分类在本章只作为已完成的算例重新解释。

\ProseParagraphBreak
初读时先掌握\secref{b1:sec:0601}和\secref{b1:sec:0602}的直积、直和与半直积，以及\cref{b1:thm:split-extension-semidirect}的分裂判别和\cref{b1:cor:abelian-kernel-action}的交换核诱导作用；再用\subsecref{b1:subsec:0604:01}和\subsecref{b1:subsec:0604:02}的具体扩张检验这些概念。\subsecref{b1:subsec:0603:03}允许核不交换，把截面、共轭和结合律合在一起描述一般扩张，可作为综合训练；继续学习下一章的正规列与可解群时，可以暂缓其中的一般因子集计算。\subsecref{b1:subsec:0604:03}末尾的上同调解释是后续学习的预告，不是本章证明的前提。

\ProseParagraphBreak
直积与扩张可参阅 Rotman 的《An Introduction to the Theory of Groups》\cite[Chapters 6--7]{Rotman1995TheoryGroups}。若要继续研究分裂扩张与交换核扩张的上同调分类，可在学习模与基本同调代数后阅读 Brown 的《Cohomology of Groups》\cite[Chapter IV, \S\S~2--3, pp. 87--96]{Brown1982CohomologyGroups}；非交换核扩张的概述见同章\cite[Chapter IV, \S~6, pp. 104--106]{Brown1982CohomologyGroups}。本章先用具体乘法说明截面、半直积和因子集，不以群上同调为前提。

% !TeX root = ../../Book1.tex
% 纲要标题：### 6.1 直积与直和
% 纲要位置：计划纲要.md，第 273 行。

\section{直积与直和}
\label{b1:sec:0601}

% TODO: 按以下纲要要点撰写本节。

% 纲要内容（仅作写作计划，尚非教材正文）：
%
% * 有限直积及其投影
% * 任意群族的直积与限制直积式直和
% * 阿贝尔群范畴中的直和
% * 非交换群的有限支集直积是直积型构造，不是群范畴中的余积；后者由第8.3节的自由积承担
%

\subsection{有限直积及其投影}
\label{b1:subsec:0601:01}

待撰写。

\subsection{任意直积与有限支集直积}
\label{b1:subsec:0601:02}

待撰写。

\subsection{阿贝尔群的直和}
\label{b1:subsec:0601:03}

待撰写。

\subsection{有限支集直积与群的余积}
\label{b1:subsec:0601:04}

待撰写。

% !TeX root = ../../Book1.tex
% 纲要标题：### 6.2 半直积
% 纲要位置：计划纲要.md，第 280 行。

\section{半直积}
\label{b1:sec:0602}

% TODO: 按以下纲要要点撰写本节。

% 纲要内容（仅作写作计划，尚非教材正文）：
%
% * 自同构作用与外半直积
% * 内半直积判别
% * 二面体群、仿射群与具体计算
%

\subsection{自同构作用与外半直积}
\label{b1:subsec:0602:01}

待撰写。

\subsection{内半直积判别}
\label{b1:subsec:0602:02}

待撰写。

\subsection{二面体群、仿射群与具体计算}
\label{b1:subsec:0602:03}

待撰写。

% !TeX root = ../../Book1.tex
% 纲要标题：### 6.3 群扩张
% 纲要位置：计划纲要.md，第 286 行。

\section{群扩张}
\label{b1:sec:0603}

% 本节正文按下列原纲要要点撰写。

% 纲要内容（仅作写作计划，尚非教材正文）：
%
% * 短正合群列与分裂扩张
% * 截面是集合映射还是群同态
% * 扩张数据与半直积的关系
%

\subsection{短正合群列与分裂扩张}
\label{b1:subsec:0603:01}

商群把正规子群压缩为单位；群扩张提出相反的问题：已知正规子群和商群，能恢复哪些群？直积和半直积给出两种答案，但并非所有扩张都属于这两类。

\begin{definition}\label{b1:def:group-extension}
一列群同态
\[
1\longrightarrow N\xrightarrow{\iota}G
\xrightarrow{\pi}H\longrightarrow1
\]
称为\term[duan-zhenghe-qunlie]{短正合群列}[short exact sequence of groups]，如果 \(\iota\) 单射、\(\pi\) 满射，且 \(\operatorname{im}\iota=\ker\pi\)。它也称为 \(H\) 由 \(N\) 的一个\term[qun-kuozhang]{群扩张}[group extension]。若存在群同态 \(s:H\rightarrow G\) 满足 \(\pi s=\operatorname{id}_H\)，称扩张\term[fenlie-kuozhang]{分裂}[split extension]。
\end{definition}
这里的 \(1\) 表示平凡群。“正合”在每个位置要求前一个映射的像等于后一个映射的核，因此两端的条件正好是单射与满射。通过 \(\iota\)，可将 \(N\) 识别为 \(G\) 的正规子群；此时 \(G/N\cong H\)。但扩张还记住了指定映射 \(\iota,\pi\)，比只写一个抽象同构 \(G/N\cong H\) 包含更多信息。

\begin{theorem}\label{b1:thm:split-extension-semidirect}
短正合群列分裂，当且仅当 \(G\) 含有子群 \(L\)，使 \(\pi|_L:L\rightarrow H\) 为同构。存在这样的 \(L\) 时，\(G\) 是 \(N\) 与 \(L\) 的内半直积；若选择了同态截面 \(s\)，则
\[
G\cong N\rtimes_\alpha H,\qquad
\alpha_h(n)=s(h)ns(h)^{-1},
\]
同构由 \((n,h)\mapsto ns(h)\) 给出。
\end{theorem}
\begin{proof}
若有同态截面 \(s\)，则由 \(\pi s=\operatorname{id}\) 可知 \(s\) 单射，像 \(L=s(H)\) 是子群，且 \(\pi|_L\) 与 \(s\) 互逆。对任意 \(g\in G\)，令 \(h=\pi(g)\)，则 \(gs(h)^{-1}\in\ker\pi=N\)，所以 \(G=NL\)。若 \(s(h)\in N\)，应用 \(\pi\) 得 \(h=1\)，故 \(N\cap L=\{1\}\)。内半直积判别给出所需同构。

反过来，若 \(\pi|_L\) 是同构，取它的逆再接子群嵌入，得到同态截面。特别地，外半直积投影 \(N\rtimes H\rightarrow H\) 总有同态截面 \(h\mapsto(1,h)\)。
\end{proof}

\subsection{集合截面与同态截面}
\label{b1:subsec:0603:02}

满射 \(\pi:G\rightarrow H\) 的\term[jiemian]{截面}[section]，作为集合映射，仅要求为每个 \(h\) 选择一个原像 \(s(h)\)。有限群时可逐个选择，一般情形可借助选择公理。本节总把单位的代表选为单位，即 \(s(1)=1\)。集合截面存在不等于群扩张分裂，后者要求这个选择还能保持乘法。

\ProseParagraphBreak
考虑加法群的满射
\[
\mathbb Z/4\mathbb Z\longrightarrow\mathbb Z/2\mathbb Z,
\qquad [a]_4\longmapsto[a]_2.
\]
选择 \(s([0]_2)=[0]_4\)、\(s([1]_2)=[1]_4\) 得到集合截面，但
\[
s([1]_2+[1]_2)=[0]_4,\qquad
s([1]_2)+s([1]_2)=[2]_4,
\]
两者不同。另一种可能的非零陪集代表为 \([3]_4\)，其加倍仍为 \([2]_4\)，所以也不能给出同态截面。这里核与商群都同构于 \(C_2\)，中间群却是 \(C_4\)，并非 \(C_2\times C_2\)。

\ProseParagraphBreak
对任意归一化集合截面 \(s\)，每个 \(g\in G\) 仍有唯一表达 \(g=ns(h)\)：先令 \(h=\pi(g)\)，再令 \(n=gs(h)^{-1}\in N\)。因此 \(G\) 的底集合总可用 \(N\times H\) 的坐标描述；困难在于乘法，而非集合的分解。与半直积相比，新出现的项正好记录截面未能保持乘法的差额。

\subsection{扩张数据与半直积}
\label{b1:subsec:0603:03}

给定归一化集合截面，定义
\[
\alpha_h(n)=s(h)ns(h)^{-1},\qquad
c(h,k)=s(h)s(k)s(hk)^{-1}\in N.
\]
每个 \(\alpha_h\) 都是 \(N\) 的自同构，但映射 \(h\mapsto\alpha_h\) 未必是同态。元素 \(c(h,k)\) 称为这个截面的\term[yinzi-ji]{因子集}[factor set]；它等于单位当且仅当截面对相应的一对元素保持乘法。

\begin{proposition}\label{b1:prop:extension-factor-data}
在坐标 \(g=ns(h)\) 下，扩张中的乘法为
\[
(n,h)(m,k)=(n\alpha_h(m)c(h,k),hk).
\]
这些数据满足 \(\alpha_1=\operatorname{id}\)、\(c(1,h)=c(h,1)=1\)，并满足
\begin{align}
\alpha_h\alpha_k
&=\operatorname{Inn}_{c(h,k)}\alpha_{hk},\label{b1:eq:extension-action-defect}\\
c(h,k)c(hk,\ell)
&=\alpha_h(c(k,\ell))c(h,k\ell),\label{b1:eq:extension-factor-identity}
\end{align}
其中 \(\operatorname{Inn}_u(n)=unu^{-1}\) 表示由 \(u\) 给出的\term[nei-zitonggou]{内自同构}[inner automorphism]。
\end{proposition}
\begin{proof}
先把 \(m\) 移过 \(s(h)\)，再用 \(s(h)s(k)=c(h,k)s(hk)\)，得到
\[
ns(h)ms(k)=n\alpha_h(m)c(h,k)s(hk),
\]
这给出乘法公式。由 \(s(h)s(k)=c(h,k)s(hk)\)，两边对 \(N\) 作共轭便得到\eqref{b1:eq:extension-action-defect}。另一方面，将 \(s(h)s(k)s(\ell)\) 分别按左右结合展开，左结合为
\[
c(h,k)c(hk,\ell)s(hk\ell),
\]
右结合为
\[
\alpha_h(c(k,\ell))c(h,k\ell)s(hk\ell).
\]
消去右端共同因子，就得到\eqref{b1:eq:extension-factor-identity}。归一化条件直接来自 \(s(1)=1\)。
\end{proof}
若把截面改为 \(s'(h)=b(h)s(h)\)，其中 \(b(h)\in N\)、\(b(1)=1\)，则直接代入定义给出
\[
\alpha'_h=\operatorname{Inn}_{b(h)}\alpha_h,\qquad
c'(h,k)=b(h)\alpha_h(b(k))c(h,k)b(hk)^{-1}.
\]
计算第二式时，唯一需要移过截面的因子是 \(b(k)\)，由
\(s(h)b(k)=\alpha_h(b(k))s(h)\) 处理。可见因子集依赖截面的选择，同一个扩张可以有不同的 \(c\)。扩张分裂当且仅当能选到使全部 \(c(h,k)=1\) 的截面；这时\eqref{b1:eq:extension-action-defect}说明 \(\alpha\) 成为真正的同态，乘法恢复为半直积。对于非交换的 \(N\)，不能删掉内自同构项并把任意 \(\alpha\) 都当作群作用。

% !TeX root = ../../Book1.tex
% 纲要标题：### 6.4 扩张的例子
% 纲要位置：计划纲要.md，第 292 行。

\section{扩张的例子}
\label{b1:sec:0604}

% TODO: 按以下纲要要点撰写本节。

% 纲要内容（仅作写作计划，尚非教材正文）：
%
% * 中心扩张与非分裂例子
% * 四元数群与二面体群的对照
% * 第四卷第二群上同调的分类解释预告
%
% ---
%

\subsection{中心扩张与非分裂例子}
\label{b1:subsec:0604:01}

待撰写。

\subsection{四元数群与二面体群}
\label{b1:subsec:0604:02}

待撰写。

\subsection{群扩张与第二群上同调}
\label{b1:subsec:0604:03}

待撰写。


\begin{chaptersummary}
直积由进入各因子的同态刻画；交换群直和由从各因子出发的同态刻画。在一般群中，有限支集直积强制不同坐标交换，因而不具有群余积的泛性质。

\ProseParagraphBreak
半直积由两个群和一个自同构作用构成；指定扩张分裂，当且仅当其商映射具有同态截面。一般集合截面仍给出唯一坐标表达，但乘法中会出现因子集；反过来，满足相容恒等式的作用数据与因子集也能在 \(N\times H\) 上构造群扩张。更换截面只改变坐标数据，不改变原扩张；能否选择截面消去因子集，才决定扩张是否分裂。\(D_8\) 与 \(Q_8\) 的比较说明，即使核、商及其诱导作用相同，仍可能得到不同的扩张和不同的中间群。

\ProseParagraphBreak
扩张等价要求中间群的同构保持指定的核嵌入与商映射，比抽象群同构更严格。\chapref{b1:ch:07}将用正规列逐层分解群；本章说明了每一层的核与商怎样联系，以及为什么只知道各层因子通常不足以恢复整个群。
\end{chaptersummary}

\BookExercises

\begin{exercise}\label{b1:ex:06-1}
设 \(G\) 为群，\(\Delta=\bigl\{\,(g,g)\mid g\in G\,\bigr\}\leq G\times G\)。
\begin{enumerate}
\item 证明 \(\Delta\trianglelefteq G\times G\) 当且仅当 \(G\) 是交换群。
\item 设 \(G\neq\{1\}\)。证明 \(\Delta\) 的两个坐标投影均满射，但 \(\Delta\) 不是 \(G\times G\) 中两个坐标子群的直积。
\end{enumerate}
\end{exercise}
\begin{solution}
\textbf{（1）}
若 \(\Delta\) 正规，则对任意 \(a,g\in G\)，
\[
(a,1)(g,g)(a^{-1},1)=(aga^{-1},g)\in\Delta,
\]
故 \(aga^{-1}=g\)，所以 \(G\) 交换。反过来，若 \(G\) 交换，则 \(G\times G\) 交换，其中每个子群都正规，包括 \(\Delta\)。

\textbf{（2）}
当 \(G\neq\{1\}\) 时，\(\Delta\) 的两个坐标投影都是 \(G\)，但取 \(g\neq1\)，便有 \((g,1)\notin\Delta\)，所以 \(\Delta\neq G\times G\)。若 \(\Delta\) 是两个坐标子群的直积，这两个子群由投影都只能是 \(G\)，与上述严格不等矛盾。因此直积中的子群不必是两个因子中子群的坐标直积。
\end{solution}

\begin{exercise}\label{b1:ex:06-2}
设 \(m,n\) 为正整数。证明 \(C_m\times C_n\) 循环当且仅当 \(\gcd(m,n)=1\)，并求 \(C_4\times C_6\) 中元素阶的最大值。
\end{exercise}
\begin{solution}
若 \(a,b\) 分别生成 \(C_m,C_n\)，则 \((a,b)^d=1\) 当且仅当 \(m\mid d\) 且 \(n\mid d\)，故其阶为 \(\operatorname{lcm}(m,n)\)。任意元素的阶也整除这个最小公倍数。因此直积有阶为 \(mn\) 的生成元，当且仅当 \(\operatorname{lcm}(m,n)=mn\)，等价于 \(\gcd(m,n)=1\)。对 \(4,6\)，最大元素阶为 \(12\)，由两个生成元组成的元素达到，但直积阶为 \(24\)，所以不是循环群。
\end{solution}

\begin{exercise}\label{b1:ex:06-infinite-product-maps}
令 \(P=\prod_{r\geq1}\mathbb Z\)、\(S=\bigoplus_{r\geq1}\mathbb Z\)，并将 \(S\) 看成 \(P\) 的有限支集子群。记第 \(r\) 个坐标为 \(1\)、其他坐标为 \(0\) 的序列为 \(e_r\)。证明 \(S\) 由全部 \(e_r\) 生成，而 \(P\) 不由它们生成。再给出两个不同的同态 \(P\rightarrow P/S\)，使它们在每个坐标子群 \(\mathbb Z e_r\) 上的限制相同。
\end{exercise}
\begin{solution}
有限多个 \(e_r\) 的整数线性组合只有有限个非零坐标，故生成的子群包含于 \(S\)。反过来，若 \(a=(a_r)\in S\)，其支集 \(J\) 有限，就有 \(a=\sum_{r\in J}a_re_r\)。所以这些元素生成的子群恰为 \(S\)。常数序列 \(v=(1,1,\ldots)\) 不属于 \(S\)，因此它们不生成整个 \(P\)。

\(P\) 交换，故 \(S\trianglelefteq P\)，自然商映射 \(q:P\rightarrow P/S\) 与零同态 \(z:P\rightarrow P/S\) 都有定义。对每个 \(r\)，坐标子群 \(\mathbb Z e_r\) 包含于 \(S\)，所以两者在该子群上都为零。但 \(q(v)=v+S\neq S=z(v)\)，因而两个同态不同。从无限直积出发的同态，未必由它在各坐标子群上的限制决定。
\end{solution}

\begin{exercise}\label{b1:ex:06-3}
设 \(N\) 为加法交换群，\(H\) 为交换群，且 \(\alpha:H\rightarrow\operatorname{Aut}(N)\) 为群同态。令
\(N^H=\bigl\{\,n\in N\mid\alpha_h(n)=n\text{ 对所有 }h\in H\,\bigr\}\)。
证明 \(Z(N\rtimes_\alpha H)=N^H\times\ker\alpha\)，并计算非交换 \(21\) 阶群的中心。
\end{exercise}
\begin{solution}
元素 \((n,h)\) 与每个 \((m,1)\) 交换，当且仅当
\(n+\alpha_h(m)=m+n\) 对所有 \(m\) 成立，即 \(h\in\ker\alpha\)。它与每个 \((0,k)\) 交换，当且仅当 \(n=\alpha_k(n)\)，这里第二坐标交换用到了 \(H\) 交换。两种坐标子群生成整个半直积，所以这两个条件也充分。它们组成的群运算为直积运算，因为 \(\ker\alpha\) 作用平凡。对\cref{b1:thm:groups-order-pq}构造的非交换 \(21\) 阶群，\(N=\mathbb Z/7\mathbb Z\)、\(H=C_3\)，其生成元作用为乘以 \(2\)。共同固定元素必须满足 \(2n=n\)，所以只有 \(n=0\)；而三个作用分别为乘以 \(1,2,4\)，模 \(7\) 彼此不同，故 \(\ker\alpha=\{1\}\)。代入刚证明的中心公式，得到中心平凡。
\end{solution}

\begin{exercise}\label{b1:ex:06-4}
设整数 \(n\geq3\)。在 \(D_{2n}\) 中求旋转子群 \(\langle r\rangle\) 的全部补群。证明 \(n\) 为奇数时它们全共轭，\(n\) 为偶数时恰分成两个共轭类。
\end{exercise}
\begin{solution}
令 \(N=\langle r\rangle\)，\(\pi:D_{2n}\rightarrow D_{2n}/N\cong C_2\) 为商映射。若 \(H\) 是 \(N\) 的补群，则 \(NH=D_{2n}\) 说明 \(\pi|_H\) 满射，而 \(\ker(\pi|_H)=H\cap N=\{1\}\)。由\cref{b1:thm:first-isomorphism}，\(H\cong C_2\)，所以其非单位元素在旋转子群之外。\cref{b1:prop:dihedral-normal-form}说明所有这样的元素正是反射 \(r^is\)，且平方为单位。

反过来，每个 \(\langle r^is\rangle\) 与 \(N\) 交平凡，它在商群中的像含非单位元，故乘上 \(N\) 得到全群。因此全部补群为 \(\langle r^is\rangle\)，\(i\) 模 \(n\)，共有 \(n\) 个。

用 \(r^a\) 共轭使指数 \(i\) 变为 \(i+2a\)，用 \(s\) 共轭使指数变为 \(-i\)。若 \(n\) 奇，乘 \(2\) 在模 \(n\) 剩余类中可逆，任意两个指数都可由这种共轭互相变换。若 \(n\) 偶，这两种操作都保持指数奇偶；同奇偶指数又能通过加 \(2a\) 互相到达，因此恰有两个类。
\end{solution}

\begin{exercise}\label{b1:ex:06-5}
在实直线仿射群中，设 \(T\) 为平移子群。对 \(c\in\mathbb R\)，令 \(L_c\) 为固定点 \(c\) 的全部仿射变换。证明 \(L_c\) 是 \(T\) 的补群，并写出将 \(L_0\) 共轭到 \(L_c\) 的平移。
\end{exercise}
\begin{solution}
固定 \(c\) 的条件为 \(ac+b=c\)，故
\[
L_c=\Bigl\{\,\bigl((1-a)c,a\bigr)\mid a\in\mathbb R^\times\,\Bigr\}.
\]
它是点稳定子群，因而是子群；也可直接按半直积乘法验证。投影到 \(a\) 坐标在 \(L_c\) 上为同构，且 \(L_c\cap T=\{1\}\)。任意仿射元素有分解
\[
(b,a)=\bigl(b-(1-a)c,1\bigr)\bigl((1-a)c,a\bigr),
\]
所以 \(TL_c\) 为全群，\(L_c\) 是补群。令平移 \(\tau_c(x)=x+c\)，伸缩 \(f_a(x)=ax\)，则
\[
(\tau_c\circ f_a\circ\tau_c^{-1})(x)=a(x-c)+c.
\]
因此 \(\tau_c\) 把 \(L_0\) 共轭到 \(L_c\)。
\end{solution}

\begin{exercise}\label{b1:ex:06-complement-cocycle}
设 \(G=N\rtimes_\alpha H\)，并把 \(N\) 识别为 \(N\times\{1\}\)。对映射 \(b:H\rightarrow N\)，令
\[
L_b=\{\,\bigl(b(h),h\bigr)\mid h\in H\,\}.
\]
证明 \(L_b\) 是 \(N\) 在 \(G\) 中的补群，当且仅当
\[
b(1_H)=1_N,\qquad
b(hk)=b(h)\alpha_h\bigl(b(k)\bigr)
\quad(h,k\in H).
\]
并证明：用 \((n_0,1)\) 共轭 \(L_b\) 后所得补群对应的函数为
\[
b'(h)=n_0b(h)\alpha_h(n_0^{-1}).
\]
\end{exercise}
\begin{solution}
若题设两式成立，则
\[
\bigl(b(h),h\bigr)\bigl(b(k),k\bigr)
=\bigl(b(h)\alpha_h(b(k)),hk\bigr)
=\bigl(b(hk),hk\bigr).
\]
因此映射 \(h\mapsto\bigl(b(h),h\bigr)\) 是群同态，\(L_b\) 是其像，故为子群。第二坐标投影在 \(L_b\) 上与这个同态互逆，所以 \(L_b\cong H\) 且 \(L_b\cap N=\{1\}\)。任取 \((n,h)\in G\)，有
\[
(n,h)=\bigl(nb(h)^{-1},1\bigr)\bigl(b(h),h\bigr),
\]
故 \(NL_b=G\)，所以 \(L_b\) 是补群。

反过来，若 \(L_b\) 是补群，则第二坐标投影 \(L_b\rightarrow H\) 是同构：它由集合的定义满射，核为 \(L_b\cap N=\{1\}\)。其逆映射 \(h\mapsto\bigl(b(h),h\bigr)\) 必为同态；比较单位与乘积的第一坐标，就得到题设两式。

最后直接计算
\[
\begin{aligned}
(n_0,1)\bigl(b(h),h\bigr)(n_0^{-1},1)
&=\bigl(n_0b(h),h\bigr)(n_0^{-1},1)\\
&=\bigl(n_0b(h)\alpha_h(n_0^{-1}),h\bigr).
\end{aligned}
\]
因此共轭后的补群正是 \(L_{b'}\)，其中 \(b'\) 如题所述。这一区分了补群的存在、具体选择及不同补群是否共轭三个问题。
\end{solution}

\begin{exercise}\label{b1:ex:06-6}
设整数 \(m,n\geq2\)，采用 \(C_r=\mathbb Z/r\mathbb Z\) 的加法记号。考虑短正合列
\[
0\longrightarrow C_m\xrightarrow{\iota}C_{mn}\xrightarrow{\pi}C_n\longrightarrow0,
\]
其中 \(\iota\bigl([k]_m\bigr)=[nk]_{mn}\)，\(\pi\bigl([a]_{mn}\bigr)=[a]_n\)。证明这一扩张分裂，当且仅当 \(\gcd(m,n)=1\)。
\end{exercise}
\begin{solution}
先核对短正合性。改变 \(k\) 的代表会使 \(nk\) 改变 \(mn\) 的倍数，故 \(\iota\) 良定义；\(mn\mid nk\) 当且仅当 \(m\mid k\)，所以它单射。\(\pi\) 是满同态，其核恰由 \(n\) 的倍数的剩余类组成，也就是 \(\operatorname{im}\iota\)。

同态截面由生成元 \([1]_n\) 的像 \([a]_{mn}\) 确定，所需条件为
\[
a\equiv1\pmod n,\qquad na\equiv0\pmod{mn}.
\]
第二个条件等价于 \(m\mid a\)，写 \(a=mu\) 后，第一条件变为 \(mu\equiv1\pmod n\)。若有解，任何 \(m,n\) 的公因子都整除 \(mu-nv=1\)，所以二者互素；反过来，Bézout 等式在互素时给出这样的 \(u,v\)。取 \(a=mu\) 后，定义 \(s\Bigl([k]_n\Bigr)=[ka]_{mn}\)。若 \(k\) 改变 \(n\) 的倍数，\(ka\) 改变 \(na\) 的倍数，而 \(mn\mid na\)，所以该公式良定义；它保持加法，且 \(a\equiv1\pmod n\) 保证 \(\pi s\Bigl([k]_n\Bigr)=[k]_n\)。这就构造了同态截面。特别地，\(C_4\rightarrow C_2\) 不分裂，\(C_6\rightarrow C_3\) 分裂。
\end{solution}

\begin{exercise}\label{b1:ex:06-7}
设整数 \(n\geq2\)。对自然商映射 \(\mathbb Z\rightarrow\mathbb Z/n\mathbb Z\)，用 \(0,1,\ldots,n-1\) 作陪集代表，并通过 \(a\mapsto na\) 把核识别为 \(\mathbb Z\)。计算因子集，说明它怎样记录加法进位。
\end{exercise}
\begin{solution}
对整数 \(t\)，记 \(r_n(t)\in\{0,\ldots,n-1\}\) 为它除以 \(n\) 的余数；这是一个整数，与剩余类 \([t]_n\) 区分。对代表 \(0\leq i,j<n\)，有
\[
i+j=n\left\lfloor\frac{i+j}{n}\right\rfloor+r_n(i+j).
\]
核经 \(a\mapsto na\) 识别为 \(\mathbb Z\)，所以因子集为
\[
c(i,j)=\left\lfloor\frac{i+j}{n}\right\rfloor.
\]
其值为 \(0\) 或 \(1\)，记录一次加法是否产生进位。这里核和商群都采用加法记号，诱导作用平凡。\eqref{b1:eq:extension-factor-identity}具体成为
\[
\begin{aligned}
c(i,j)+c\Bigl(r_n(i+j),k\Bigr)
&=c(j,k)+c\Bigl(i,r_n(j+k)\Bigr)\\
&=\left\lfloor\frac{i+j+k}{n}\right\rfloor
\qquad(0\leq i,j,k<n).
\end{aligned}
\]
这表示先加哪两个数，都不改变最后的总进位。

在 \(\mathbb Z\times\{0,\ldots,n-1\}\) 上定义坐标运算
\[
(a,i)+(b,j)=\Bigl(a+b+c(i,j),r_n(i+j)\Bigr).
\]
整数除法说明 \((a,i)\mapsto na+i\) 是到 \(\mathbb Z\) 的双射；上面的余数等式又说明它把坐标运算变为整数加法。因此坐标运算确实给出一个群，并通过这个双射恢复原扩张。
\end{solution}

\begin{exercise}\label{b1:ex:06-8}
设 \(1\rightarrow A\rightarrow G\rightarrow H\rightarrow1\) 为中心扩张，\(H\) 为循环群。证明 \(G\) 交换，并说明这是否保证扩张分裂。另证明：任意满同态 \(\pi:G\rightarrow\mathbb Z\) 都有同态截面，这一结论不需要中心性假设。
\end{exercise}
\begin{solution}
取商群生成元的任一提升 \(x\in G\)。每个 \(g\in G\) 都可写为 \(ax^r\)，其中 \(a\in A\)、\(r\in\mathbb Z\)：商像写为生成元的幂后，将相应 \(x^r\) 除去，剩下的元素就在核中。因 \(A\leq Z(G)\)，对任意 \(a,b\in A\) 有
\[
(ax^r)(bx^s)=abx^{r+s}=bax^{s+r}=(bx^s)(ax^r).
\]
所以 \(G\) 交换。不过扩张不必分裂；\subsecref{b1:subsec:0603:02}已经验证模 \(2\) 的满射 \(C_4\rightarrow C_2\) 没有同态截面。这个例子表明，交换性只要求提升元素彼此交换，分裂还要求能选择满足有限循环商之阶数关系的提升。

对任意满同态 \(\pi:G\rightarrow\mathbb Z\)，取 \(x\in G\) 满足 \(\pi(x)=1\)，定义 \(s(k)=x^k\)。整数幂法则给出 \(s(k+\ell)=s(k)s(\ell)\)，而 \(\pi\bigl(s(k)\bigr)=k\)，所以 \(s\) 是同态截面。此处没有使用核的中心性。无限循环商没有有限阶关系需要提升元素满足；上面的不分裂例子来自有限循环商的这个额外要求。
\end{solution}

\begin{exercise}[综合题]\label{b1:ex:06-change-split-section}
设 \(C_4=\langle z\rangle\)、\(C_2=\{1,q\}\)，考虑第二坐标投影给出的扩张
\[
1\longrightarrow C_4\longrightarrow C_4\times C_2\longrightarrow C_2\longrightarrow1.
\]
\begin{enumerate}
\item 选择归一化集合截面 \(s(1)=(1,1)\)、\(s(q)=(z,q)\)。求其诱导作用和全部因子集值，并选取 \(b:C_2\rightarrow C_4\)，使新截面 \(s'(h)=(b(h),1)s(h)\) 的因子集恒为单位。解释这个例子为何不矛盾于分裂判别。
\item 令 \(N=S_3\)、\(H=C_2=\{1,q\}\)，考虑第二坐标投影 \(N\times H\rightarrow H\)。取 \(a=(1\,2\,3)\)，并选归一化集合截面
\[
s(1)=(1,1),\qquad s(q)=(a,q).
\]
计算相应的 \(\alpha,c\)，证明 \(\alpha:H\rightarrow\operatorname{Aut}(N)\) 不是群同态；再改变截面，使新的 \(\alpha'\) 平凡且 \(c'\equiv1\)。
\end{enumerate}
\end{exercise}
\begin{solution}
\textbf{（1）}
中间群 \(C_4\times C_2\) 交换，因此核上的诱导作用平凡。用第一坐标把核识别为 \(C_4\)，归一化给出
\[
c(1,1)=c(1,q)=c(q,1)=1,\qquad c(q,q)=z^2.
\]
最后一式来自 \(s(q)^2=(z^2,1)\)，故当前截面不是同态。

取 \(b(1)=1\)、\(b(q)=z^{-1}\)。将\eqref{b1:eq:factor-set-section-change}用于平凡作用，得到
\[
c'(h,k)=b(h)b(k)c(h,k)b(hk)^{-1}.
\]
有单位输入时，新因子集值仍为单位；唯一剩下的一对给出
\[
c'(q,q)=z^{-1}z^{-1}z^2=1.
\]
对应截面满足 \(s'(1)=(1,1)\)、\(s'(q)=(1,q)\)，正是同态截面。分裂要求存在这样的截面，并不要求所有集合截面都是同态；更换截面消去了这次坐标选择带来的非平凡因子集，扩张本身没有改变。

\textbf{（2）}
通过第一坐标把 \(N\) 识别为投影的核。截面对单位给出恒等自同构，而
\[
\alpha_q(n)=ana^{-1},
\qquad
c(1,1)=c(1,q)=c(q,1)=1,
\qquad
c(q,q)=a^2.
\]
最后一式来自 \(s(q)^2=(a^2,1)\)。这里
\[
\alpha_q^2=\operatorname{Inn}(a^2)\neq\operatorname{id}_N
=\alpha_{q^2};
\]
例如 \(\alpha_q^2\bigl((1\,2)\bigr)=(1\,3)\)。因此 \(\alpha\) 不是群同态。

取 \(b(1)=1\)、\(b(q)=a^{-1}\)，并令 \(s'(h)=(b(h),1)s(h)\)。于是
\[
s'(1)=(1,1),\qquad s'(q)=(1,q).
\]
新截面是第二坐标投影的标准同态截面；它的元素与 \(N\times\{1\}\) 逐元素交换，所以 \(\alpha'_q=\operatorname{id}_N\)，并且 \(s'(h)s'(k)=s'(hk)\) 给出 \(c'(h,k)=1\) 对所有 \(h,k\in H\) 成立。原扩张本来分裂，并不意味着任意集合截面产生的 \(\alpha\) 都是群作用。
\end{solution}

% !TeX root = ../../Book1.tex
% 纲要标题：## 第7章　正规列、可解群与幂零群
% 纲要位置：计划纲要.md，第 300 行。

\chapter{正规列、可解群与幂零群}
\label{b1:ch:07}
\BookChapterNumber{b1:ch:07}

\begin{chapterabstract}
群可以沿正规子群逐层分解，也可以用交换子逐步消去非交换性。本章证明合成因子的唯一性，建立可解群和幂零群的基本判据，并证明五次以上交错群的单性。这些结果把前面零散的有限群算例组织成结构理论，同时为后面的 Galois 理论准备可解性这一判别工具。
\end{chapterabstract}

\begin{learninggoals}
\begin{itemize}
\item 分清正规列与次正规列，掌握蝴蝶引理如何配对细化的因子。
\item 会计算常见群的导出列和中心列，并检查可解性、幂零性的稳定性质。
\item 能把有限幂零群拆为 Sylow 子群直积，理解交错群单性证明中的最小支持度方法。
\end{itemize}
\end{learninggoals}

本章使用\chapref{b1:ch:02}的同构定理、\chapref{b1:ch:04}的类方程及\chapref{b1:ch:05}的 Sylow 理论。\secref{b1:sec:0701}的合成列理论与后面的交换子方法可先分别阅读，再通过“有限可解群的合成因子”联系起来。伴读可选 Dummit 与 Foote 的《Abstract Algebra》\cite{DummitFoote2003}中群结构的有关内容；本章所用的结构定理均在正文证明。

% !TeX root = ../../Book1.tex
% 纲要标题：### 7.1 正规列与合成列
% 纲要位置：计划纲要.md，第 302 行。

\section{正规列与合成列}
\label{b1:sec:0701}

% 纲要内容（仅作写作计划，尚非教材正文）：
%
% * 次正规列、合成列及细化
% * Schreier 细化定理
% * Jordan–Hölder 定理
%

\subsection{次正规列、合成列及细化}
\label{b1:subsec:0701:01}

前两章用一个正规子群把群分成子群与商群。反复进行这种分解，便得到一列较简单的因子。不过，某个群可以有许多不同的正规子群；要使分解成为结构工具，必须说明不同选择之间有什么共同之处。

\begin{definition}\label{b1:def:subnormal-series}
群 \(G\) 的一个\term[cizhenggui-lie]{次正规列}[subnormal series]是有限链
\[
G=G_0\trianglerighteq G_1\trianglerighteq\cdots
\trianglerighteq G_r=1,
\]
其中只要求 \(G_{i+1}\) 在 \(G_i\) 中正规。若各 \(G_i\) 都在 \(G\) 中正规，称为\term[zhenggui-lie]{正规列}[normal series]。商群 \(G_i/G_{i+1}\) 称为列的因子；插入若干子群得到的新次正规列称为原列的\term[xihua]{细化}[refinement]。
\end{definition}
有些文献把本书的次正规列也称为正规列。这里区分二者，是因为“在前一项中正规”不推出“在整个群中正规”。例如
\[
S_4\trianglerighteq V_4\trianglerighteq
\langle(12)(34)\rangle\trianglerighteq1
\]
是次正规列，但中间的二阶子群不在 \(S_4\) 中正规。

\begin{definition}\label{b1:def:composition-series}
回顾单群是只有自身和单位子群两个正规子群的非平凡群。所有因子都是单群的严格次正规列称为\term[hecheng-lie]{合成列}[composition series]，其因子称为\term[hecheng-yinzi]{合成因子}[composition factor]。
\end{definition}
对应定理说明，合成列恰是不能再插入不同子群的严格次正规列。每个非平凡有限群都存在合成列：选一个极大真正规子群，再对它重复；子群的阶严格下降，过程必停。无限群未必如此，\(\mathbb Z\) 的任何非零子群都同构于 \(\mathbb Z\)，都还含有非零真子群，因而不可能在有限步后到达一个单的末项。

\ProseParagraphBreak
若删去平凡因子后，可以把两列的因子一一配对，使每对同构，就称它们为\term[dengjia-lie]{等价列}[equivalent series]；配对不要求保留原有次序。例如 \(C_6\) 可以先取二阶子群，也可以先取三阶子群，两次所得因子分别是 \(C_3,C_2\) 与 \(C_2,C_3\)。

\subsection{Schreier 细化定理}
\label{b1:subsec:0701:02}

比较两条列时，应把一条列中的子群逐个与另一条列相交，再乘回原来较低的一项。使这些商群成对同构的工具是\term[hudie-yinli]{蝴蝶引理}[butterfly lemma]，也称 Zassenhaus 引理。

\begin{lemma}[Zassenhaus]\label{b1:lem:butterfly}
设 \(A'\trianglelefteq A\)、\(B'\trianglelefteq B\) 是同一群中的子群。则
\[
\frac{A'(A\cap B)}{A'(A\cap B')}
\cong
\frac{B'(B\cap A)}{B'(B\cap A')}.
\]
式中分母分别正规于分子。
\end{lemma}
\begin{proof}
令 \(C=A\cap B\)，以及
\[
D=(A'\cap B)(A\cap B').
\]
两个乘积因子都正规于 \(C\)，所以 \(D\trianglelefteq C\)。记
\(K=A'(A\cap B)\)、\(L=A'(A\cap B')\)。映射 \(K\to A/A'\) 的像由 \(C\) 给出；\(L\) 的像是 \(A\cap B'\) 的像。由于 \(A\cap B'\) 正规于 \(C\)，这个像正规于 \(K/A'\)，从而 \(L\trianglelefteq K\)。

包含映射 \(C\to K\) 后接商映射给出满同态 \(C\to K/L\)。其核为 \(C\cap L=D\)：若 \(c=a'b'\in C\)，其中 \(a'\in A'\)、\(b'\in A\cap B'\)，则 \(a'=c(b')^{-1}\in A'\cap B\)；反向包含直接成立。于是 \(K/L\cong C/D\)。交换 \((A,A')\) 与 \((B,B')\)，另一个商群也同构于 \(C/D\)。
\end{proof}

\begin{theorem}[Schreier]\label{b1:thm:schreier-refinement}
同一群的任意两条次正规列都有等价的细化。
\end{theorem}
\begin{proof}
把两条列记为 \((G_i)_{i=0}^r\) 与 \((H_j)_{j=0}^s\)。在 \(G_i\) 和 \(G_{i+1}\) 之间插入
\[
G_{ij}=G_{i+1}(G_i\cap H_j),\qquad 0\le j\le s.
\]
它的首末项分别为 \(G_i\)、\(G_{i+1}\)。蝴蝶引理保证相邻项的正规性，因此把这些小段接起来得到第一条列的细化。同样在 \(H_j\) 与 \(H_{j+1}\) 之间插入
\[
H_{ji}=H_{j+1}(H_j\cap G_i),\qquad 0\le i\le r.
\]
对每对 \((i,j)\)，蝴蝶引理给出
\[
G_{ij}/G_{i,j+1}\cong H_{ji}/H_{j,i+1}.
\]
两个矩形数组按行和按列读取的次序不同，但格子的对应提供了所需配对。重复项只产生平凡因子，删掉即可。
\end{proof}
这就是\term[schreier-xihua-dingli]{Schreier 细化定理}[Schreier refinement theorem]。它比较的是商群，而不是断言两条细化中的子群相同。事实上，\(C_2\times C_2\) 中任取一个二阶子群都给出合成列，三种选择的中间子群互不相同。

\subsection{Jordan–Hölder 定理}
\label{b1:subsec:0701:03}

\begin{theorem}[Jordan--Hölder]\label{b1:thm:jordan-holder}
若群 \(G\) 有合成列，则它的任意两条合成列长度相同，并且合成因子在同构及排列的意义下相同。
\end{theorem}
\begin{proof}
对两条合成列应用 Schreier 细化定理。合成列中的相邻因子是单群；由对应定理，在相邻两项之间不能插入不同的新子群。因此细化只会重复原有项。删去这些重复项之后，细化的非平凡因子就是原列的全部因子。两条细化等价，结论随即成立。
\end{proof}
\term[jordan-holder-dingli]{Jordan--Hölder 定理}[Jordan--Hölder theorem]使合成因子的无序集合（计入重数）成为群的不变量。不过这些因子不足以重建原群：\(C_4\) 和 \(C_2\times C_2\) 都有两个 \(C_2\) 因子，却不相互同构。合成因子记录了分解后的简单部分，部分之间怎样结合仍是扩张问题。

\ProseParagraphBreak
以 \(S_4\) 为例，链
\[
S_4\trianglerighteq A_4\trianglerighteq V_4
\trianglerighteq\langle(12)(34)\rangle\trianglerighteq1
\]
给出因子 \(C_2,C_3,C_2,C_2\)，故是合成列。因子阶的乘积为 \(24\)，与 \(S_4\) 的阶相符。以后遇到一个有限群的候选合成列时，可以先用阶数乘积排查计算错误，再逐项检查正规性和单性；阶数一致本身并不能代替后两项检查。

% !TeX root = ../../Book1.tex
% 纲要标题：### 7.2 可解群
% 纲要位置：计划纲要.md，第 308 行。

\section{可解群}
\label{b1:sec:0702}

% TODO: 按以下纲要要点撰写本节。

% 纲要内容（仅作写作计划，尚非教材正文）：
%
% * 导出列与交换子
% * 可解性的子群、商群及扩张稳定性
% * 对称群和三角矩阵群的例子
%

\subsection{导出列与交换子}
\label{b1:subsec:0702:01}

待撰写。

\subsection{可解性的稳定性}
\label{b1:subsec:0702:02}

待撰写。

\subsection{对称群与三角矩阵群}
\label{b1:subsec:0702:03}

待撰写。

% !TeX root = ../../Book1.tex
% 纲要标题：### 7.3 幂零群
% 纲要位置：计划纲要.md，第 314 行。

\section{幂零群}
\label{b1:sec:0703}

% 纲要内容（仅作写作计划，尚非教材正文）：
%
% * 上、下中心列
% * 有限 p 群的幂零性
% * 有限幂零群与 Sylow 子群直积
%

\subsection{上、下中心列}
\label{b1:subsec:0703:01}

可解群允许每一步的商群交换，却不要求该商群与群中其他元素相容。若一层元素不仅彼此交换，而且在模去更低一层以后与整个群交换，就得到更强的条件。

\begin{definition}\label{b1:def:central-series}
群 \(G\) 的\term[xia-zhongxin-lie]{下中心列}[lower central series]定义为
\[
\gamma_1(G)=G,\qquad \gamma_{i+1}(G)=[G,\gamma_i(G)].
\]
其\term[shang-zhongxin-lie]{上中心列}[upper central series]定义为
\[
Z_0(G)=1,\qquad
Z_{i+1}(G)/Z_i(G)=Z\bigl(G/Z_i(G)\bigr).
\]
右式通过商映射的原像定义 \(Z_{i+1}(G)\)。若 \(\gamma_{c+1}(G)=1\)，称 \(G\) 为\term[mi-ling-qun]{幂零群}[nilpotent group]；最小这样的非负整数 \(c\) 是其幂零类。
\end{definition}
\symidx{\(\gamma_i(G)\)：下中心列}
\symidx{\(Z_i(G)\)：上中心列}
这些子群都是特征子群。对下中心列，这由交换子在自同构下的相容性归纳得出；对上中心列，自同构先诱导商群的自同构，再保持该商群的中心。上中心列的定义也可写为
\[
x\in Z_{i+1}(G)\quad\Longleftrightarrow\quad
[g,x]\in Z_i(G)\text{ 对每个 }g\in G\text{ 成立}.
\]

\begin{proposition}\label{b1:prop:central-series-equivalence}
对任意 \(c\ge0\)，以下条件等价：
\begin{enumerate}
\item \(\gamma_{c+1}(G)=1\)；
\item \(Z_c(G)=G\)；
\item 存在正规子群链 \(1=C_0\le C_1\le\cdots\le C_c=G\)，使 \([G,C_i]\le C_{i-1}\)。
\end{enumerate}
每个幂零群都可解。
\end{proposition}
\begin{proof}
第一条件成立时，取 \(C_i=\gamma_{c+1-i}(G)\)，即得第三条件。若第三条件成立，由 \(C_0=Z_0(G)\) 和上述中心的判别逐步得到 \(C_i\le Z_i(G)\)，所以第二条件成立。第二条件成立时，直接取 \(C_i=Z_i(G)\) 得到第三条件；由 \(\gamma_1(G)=C_c\) 及交换子的单调性可归纳推出 \(\gamma_j(G)\le C_{c+1-j}\)，于是得到第一条件。

第三条件的每个因子 \(C_i/C_{i-1}\) 都阿贝尔，因此 \(G\) 可解。
\end{proof}
非平凡阿贝尔群的幂零类为一。四元数群满足
\[
Q_8'=Z(Q_8)=\{1,-1\},\qquad [Q_8,\{1,-1\}]=1,
\]
所以幂零类为二。相反，\(S_3\) 虽然可解，中心却平凡；其上中心列永远停留在单位群，故不幂零。由交换子的包含关系还可得幂零性传给子群及商群，但“可解扩张仍可解”不能换成“幂零扩张仍幂零”：\(C_3\trianglelefteq S_3\) 的核与商都阿贝尔就是反例。

\subsection{有限 \texorpdfstring{\(p\)}{p} 群的幂零性}
\label{b1:subsec:0703:02}

\begin{theorem}\label{b1:thm:p-group-nilpotent}
每个有限 \(p\) 群都幂零。若群的阶为 \(p^a\)，则幂零类至多为 \(a\)。
\end{theorem}
\begin{proof}
\chapref{b1:ch:04}的类方程已经说明，非平凡有限 \(p\) 群的中心非平凡。如果 \(Z_i(G)\ne G\)，那么 \(G/Z_i(G)\) 仍是非平凡有限 \(p\) 群，故其中心非平凡，得到严格包含 \(Z_i(G)<Z_{i+1}(G)\)。每次严格增大至少使阶乘上 \(p\)；从一阶增大到 \(p^a\) 至多需要 \(a\) 步。因此 \(Z_a(G)=G\)。
\end{proof}
这里并非只使用“中心非平凡”一次，而是在每一层商群中重新应用它。一般非平凡中心并不足以保证幂零，例如 \(C_2\times S_3\) 的中心非平凡，但其商群含有不可幂零的 \(S_3\)。

\ProseParagraphBreak
对阶为 \(8\) 的二面体群 \(D_8=\langle r,s\mid r^4=s^2=1,\ srs^{-1}=r^{-1}\rangle\)，有
\[
[s,r]=r^{-2}=r^2,\quad
\gamma_2(D_8)=\langle r^2\rangle,\quad
\gamma_3(D_8)=1.
\]
它与 \(Q_8\) 同为二类幂零群，但两群不相互同构：前者有五个二阶元，后者只有一个。幂零类衡量的是交换子消失的层数，不是群的完全分类资料。

\subsection{有限幂零群与 Sylow 子群直积}
\label{b1:subsec:0703:03}

有限群的幂零性还可以从 Sylow 子群是否彼此独立看出来。证明先用一个对无限幂零群也成立的观察。

\begin{lemma}\label{b1:lem:nilpotent-normalizer}
幂零群的每个真子群 \(H<G\) 都严格包含在其正规化子 \(N_G(H)\) 中。
\end{lemma}
\begin{proof}
取最小的 \(i\) 使 \(Z_i(G)\not\le H\)，再取 \(z\in Z_i(G)\setminus H\)。此时 \(Z_{i-1}(G)\le H\)。对任意 \(h\in H\)，有 \([z,h]\in Z_{i-1}(G)\le H\)，于是 \(zhz^{-1}\in H\)。对 \(z^{-1}\in Z_i(G)\) 作同样论证，得到 \(zHz^{-1}=H\)。所以 \(z\in N_G(H)\setminus H\)。
\end{proof}

\begin{theorem}\label{b1:thm:finite-nilpotent-sylow}
对有限群 \(G\)，以下条件等价：\(G\) 幂零；\(G\) 的每个 Sylow 子群都正规；\(G\) 是其各个 Sylow 子群的内直积。
\end{theorem}
\begin{proof}
设 \(G\) 幂零，\(P\) 是一个 Sylow 子群，令 \(N=N_G(P)\)。在 \(N\) 中，\(P\) 是正规 Sylow 子群，故是唯一的对应素数的 Sylow 子群，从而是 \(N\) 的特征子群。任何正规化 \(N\) 的元素都保持 \(P\)，所以
\(N_G(N)\le N_G(P)=N\)。若 \(N<G\)，这与上一引理矛盾。因此 \(N=G\)，\(P\trianglelefteq G\)。

反过来，设各 Sylow 子群 \(P_1,\ldots,P_t\) 正规。对 \(i\ne j\)，正规性给出 \([P_i,P_j]\le P_i\cap P_j=1\)，所以不同因子逐元素交换。乘积映射
\[
P_1\times\cdots\times P_t\longrightarrow G,
\qquad (x_1,\ldots,x_t)\longmapsto x_1\cdots x_t
\]
是同态。其核平凡：若某一 \(x_i\) 等于其余因子乘积的逆，那么其阶一方面是对应素数的幂，另一方面整除其余因子群阶的乘积，只能为一。定义域阶等于 \(|G|\)，故此映射为同构。

最后，有限 \(p\) 群已经证明幂零，而直积的交换子逐分量计算，因此
\(\gamma_k(P_1\times\cdots\times P_t)=\prod_i\gamma_k(P_i)\)。取各因子幂零类的最大值，得到直积幂零。
\end{proof}
例如，阶为 \(12\) 的 \(C_3\times C_4\) 和 \(C_3\times V_4\) 均幂零，\(A_4\) 却不幂零，因为其 Sylow 三子群不唯一。判断有限群时，可先检查 Sylow 子群；计算无限群时，则应回到中心列的定义，不能直接搬用有限 Sylow 直积判据。

% !TeX root = ../../Book1.tex
% 纲要标题：### 7.4 交错群的单性
% 纲要位置：计划纲要.md，第 320 行。

\section{交错群的单性}
\label{b1:sec:0704}

% 纲要内容（仅作写作计划，尚非教材正文）：
%
% * 三轮换生成与正规子群分析
% * n≥5 时交错群的单性
% * 与一般方程不可根式求解的后续联系
%
% ---
%

\subsection{三轮换生成与正规子群}
\label{b1:subsec:0704:01}

一个非平凡正规子群一旦包含某种置换，就必须包含它的所有群内共轭。因此，证明交错群单性的关键是找到足够简单且共轭丰富的元素。三轮换最适合这个任务：它是移动字母数最少的非单位偶置换。

\begin{lemma}\label{b1:lem:three-cycles-generate}
三轮换生成 \(A_n\)。当 \(n\ge5\) 时，所有三轮换在 \(A_n\) 中共轭。
\end{lemma}
\begin{proof}
三轮换生成 \(A_n\) 已在\secref{b1:sec:0302}证明；这里只需补充群内共轭性。

任意两个三轮换 \(\alpha,\beta\) 在 \(S_n\) 中共轭，设 \(u\alpha u^{-1}=\beta\)。若 \(u\) 为偶置换，已经得到所需共轭。若 \(u\) 为奇置换，因 \(n\ge5\)，可在 \(\beta\) 的支集之外选两个字母构成换位 \(t\)。它与 \(\beta\) 交换，且 \(tu\) 为偶置换，于是 \((tu)\alpha(tu)^{-1}=\beta\)。
\end{proof}
因此只要证明非平凡正规子群包含一个三轮换，就能推出它等于整个 \(A_n\)。下面的交换子方法把任意非单位元改造成移动较少字母的元素。

\begin{definition}\label{b1:def:permutation-support}
置换 \(\sigma\) 的\term[zhichi]{支集}[support]是
\(\operatorname{supp}(\sigma)=\{i\mid\sigma(i)\ne i\}\)。其大小称为此处的支持度。
\end{definition}
若 \(N\trianglelefteq A_n\)、\(\sigma\in N\)、\(\tau\in A_n\)，则
\[
[\sigma,\tau]=\sigma\tau\sigma^{-1}\tau^{-1}\in N.
\]
如果 \(\tau\) 是三轮换，两个乘积因子的支集各有三个字母，故交换子的支持度至多为六。还必须确保交换子非平凡；下述证明在每次使用时都核对这一点。

\subsection{交错群的单性}
\label{b1:subsec:0704:02}

\begin{theorem}\label{b1:thm:alternating-simple}
当 \(n\ge5\) 时，交错群 \(A_n\) 是非阿贝尔单群。
\end{theorem}
\begin{proof}
设 \(1\ne N\trianglelefteq A_n\)，在 \(N\setminus\{1\}\) 中取支持度最小的元素 \(\sigma\)，记该支持度为 \(m\)。偶置换不可能恰好移动一个或两个字母，所以 \(m\ge3\)。

先证 \(m\le6\)。若 \(m>6\)，选 \(a\) 满足 \(\sigma(a)=b\ne a\)，再选不同于 \(a,b\) 的两个字母 \(c,d\)，令 \(\tau=(acd)\)。共轭 \(\sigma\tau\sigma^{-1}\) 的支集包含 \(b\)，而 \(\tau\) 的支集不包含 \(b\)，故二者不同，\([\sigma,\tau]\ne1\)。这个交换子属于 \(N\)，支持度至多六，违背 \(m\) 的最小性。

若 \(m=3\)，则 \(\sigma\) 已是三轮换。其余可能的偶置换轮换类型只有下面几种；字母均表示互不相同的点。
\begin{enumerate}
\item 若 \(m=4\)，必有 \(\sigma=(ab)(cd)\)。选支集外的字母 \(e\)（此处用 \(n\ge5\)），令 \(\tau=(cde)\)。于是
\(\sigma\tau\sigma^{-1}=(dce)=\tau^{-1}\)，所以
\([\sigma,\tau]=\tau^{-2}=\tau\) 是 \(N\) 中的三轮换。这也表明支持度最小元不能实际具有 \(m=4\)。
\item 若 \(m=5\)，偶性迫使 \(\sigma=(abcde)\)。取 \(\tau=(abc)\)，则共轭为 \((bcd)\)。二者支集不同，所以交换子非平凡；其支集包含于 \(\{a,b,c,d\}\)，违背最小性。
\item 若 \(m=6\) 且 \(\sigma=(abc)(def)\)，取 \(\tau=(abd)\)。共轭为 \((bce)\)，所以交换子非平凡且支集包含于 \(\{a,b,c,d,e\}\)，仍与最小性矛盾。
\item 若 \(m=6\) 且 \(\sigma=(abcd)(ef)\)，取 \(\tau=(abc)\)。共轭为 \((bcd)\)，产生一个非平凡且支持度至多四的交换子，再次矛盾。
\end{enumerate}
这些情形穷尽了支持度四至六的偶置换：六轮换、三个换位，以及三轮换乘一个换位都是奇置换；含固定点的轮换型应按实际支持度计入前面的情形。因此 \(N\) 含有三轮换。由前一引理及正规性，\(N\) 包含所有三轮换，故 \(N=A_n\)。最后，\((123)\) 与 \((124)\) 不交换，说明 \(A_n\) 非阿贝尔。
\end{proof}
\term[jiaocuo-qun-danxing]{交错群的单性}[simplicity of the alternating group]在 \(n=5\) 就成立；证明没有使用多于五个字母去处理 \(A_5\) 的双换位。低阶例外也应同时记住：\(A_3\cong C_3\) 是阿贝尔单群，\(A_4\) 有真正规子群 \(V_4\)，因而不单。

\subsection{交错群与根式不可解性}
\label{b1:subsec:0704:03}

\begin{corollary}\label{b1:cor:symmetric-nonsolvable}
当 \(n\ge5\) 时，\(A_n'=A_n\)，因而 \(A_n\) 及 \(S_n\) 均不可解。对 \(n\le4\)，\(S_n\) 可解。
\end{corollary}
\begin{proof}
导出子群正规。由于 \(A_n\) 非阿贝尔，其导出子群不平凡；单性迫使它等于 \(A_n\)。所以导出列不可能下降到单位群，\(A_n\) 不可解，而可解性传给子群表明 \(S_n\) 也不可解。低阶情形在\secref{b1:sec:0702}已经算出；\(S_1,S_2\) 则直接阿贝尔。
\end{proof}
这为一般方程的根式不可解性准备了群论部分。到本卷 Galois 理论中，我们将建立“用有限次开根表达根”与适当 Galois 群的可解性之间的联系，并证明相应一般多项式的 Galois 群是对称群。两项桥梁建立之后，\(S_n\) 在 \(n\ge5\) 时不可解才转化为一般五次及更高次方程没有统一根式公式。

\ProseParagraphBreak
当前结果并不说“每个五次方程都不可根式求解”。例如 \(x^5-2=0\) 的一个根本来就是 \(\sqrt[5]{2}\)，其余根再乘五次单位根即可表达。决定具体方程的是它自己的 Galois 群，而不是次数一个数字。这里先明确逻辑分工：本章完成可解群及非阿贝尔单群的事实，域扩张与根式之间的对应留到本卷后续章节证明。


\begin{chaptersummary}
Schreier 细化定理通过交与积比较不同的次正规列，Jordan--Hölder 定理则把合成因子确定为计入重数的群不变量。导出列检测群能否由有限层交换群逐步扩张得到；中心列提出更强要求，使每一层在相应商群中居于中心。有限幂零群恰是各 Sylow 子群的直积，但可解群一般不具备这种分离性质。\(A_n\) 在 \(n\ge5\) 时的非交换性与单性给出了导出列永不下降的基本例子。
\end{chaptersummary}
\BookExercises

\begin{exercise}\label{b1:ex:07-composition-dihedral}
设 \(D_{24}=\langle r,s\rangle\)，其中 \(r\) 的阶为 \(12\)。写出一条合成列，确定合成因子，并说明 \(C_{24}\) 与它有相同合成因子但不相互同构。
\end{exercise}
\begin{solution}
取
\[
D_{24}\trianglerighteq\langle r\rangle
\trianglerighteq\langle r^2\rangle
\trianglerighteq\langle r^4\rangle\trianglerighteq1.
\]
\cref{b1:prop:power-order}给出 \(r^2,r^4\) 的阶分别为 \(6,3\)，所以子群阶依次为 \(24,12,6,3,1\)。旋转子群的指数为二，由\cref{b1:prop:normal-basic-examples}正规；后续相邻项处于循环群内，而循环群是交换群，故也满足正规性。\cref{b1:thm:lagrange}给出相邻商的阶为 \(2,2,2,3\)，由\cref{b1:prop:abelian-simple}及素数阶群的循环性，它们分别是单群 \(C_2,C_2,C_2,C_3\)，因此所写链是合成列。\cref{b1:thm:cyclic-subgroups}保证 \(C_{24}\) 中存在阶为 \(12,6,3,1\) 的嵌套子群，取此列得到相同因子。前者非交换，因为 \(srs^{-1}=r^{-1}\ne r\)；后者交换，所以不相互同构。
\end{solution}

\begin{exercise}\label{b1:ex:07-derived-map}
设 \(f:G\to H\) 是满同态，\(N=\ker f\)。证明 \(H'\cong G'/(G'\cap N)\)。若 \(G'\cap N=1\) 且 \(H\) 交换，能否推出 \(G\) 交换？
\end{exercise}
\begin{solution}
由式\eqref{b1:eq:commutator-image}及 \(f(G)=H\)，限制同态 \(f|_{G'}:G'\to H'\) 的像是整个 \(H'\)，其核为 \(G'\cap\ker f=G'\cap N\)。对这个满同态应用\cref{b1:thm:first-isomorphism}，即得所求同构。若此外 \(H'=1\)、\(G'\cap N=1\)，则 \(G'\) 既包含在 \(N\) 中又与 \(N\) 交为一，因此 \(G'=1\)。\cref{b1:prop:derived-quotient}取正规子群为 \(1\)，便说明 \(G\) 交换。
\end{solution}

\begin{exercise}\label{b1:ex:07-normal-solvable-product}
设 \(M,N\trianglelefteq G\) 都可解。证明 \(MN\) 可解，并据此说明每个有限群有唯一的最大可解正规子群。
\end{exercise}
\begin{solution}
\cref{b1:thm:second-isomorphism}用于 \(M\trianglelefteq G\)、\(N\le G\)，给出 \(MN\) 是子群，\(M\trianglelefteq MN\)，以及 \(MN/M\cong N/(M\cap N)\)；这个同构把 \(n(M\cap N)\) 送到 \(nM\)。由\cref{b1:thm:solvable-closure}的商群结论，右侧作为 \(N\) 的商群可解；再将同一定理的扩张结论用于 \(M\trianglelefteq MN\)，得到 \(MN\) 可解。又因 \(M,N\) 均在 \(G\) 中正规，任意共轭保持它们的乘积，故 \(MN\trianglelefteq G\)。有限群只有有限多个子群，把它的全部可解正规子群依次相乘，反复应用刚证的两子群乘积结论，得到可解正规子群 \(R\)。它包含每个可解正规子群，所以最大；任何另一个具有此性质的子群必须既包含 \(R\) 又被 \(R\) 包含，故唯一。
\end{solution}

\begin{exercise}\label{b1:ex:07-dihedral-series}
对 \(n\ge3\)，计算 \(D_{2n}'\) 及下中心列，并证明 \(D_{2n}\) 幂零当且仅当 \(n\) 为二的幂。若 \(n=2^a\)，求其幂零类。
\end{exercise}
\begin{solution}
由\cref{b1:prop:dihedral-normal-form}中的 \(srs^{-1}=r^{-1}\)，计算得 \([s,r]=r^{-2}\)，故 \(\langle r^2\rangle\le D_{2n}'\)。共轭生成元 \(r,s\) 都保持 \(\langle r^2\rangle\)，所以这个子群正规。模去它以后，\(r\) 的像满足 \(r=r^{-1}\)，因而与 \(s\) 的像交换，商群交换；\cref{b1:prop:derived-quotient}给出反向包含，得到 \(D_{2n}'=\langle r^2\rangle\)。对任何整数 \(t,k\)，标准形 \(r^t\) 与 \(r^ts\) 分别满足
\[
[r^t,r^k]=1,\qquad [r^ts,r^k]=r^{-2k}.
\]
因此与 \(\langle r^k\rangle\) 取交换子恰生成 \(\langle r^{2k}\rangle\)，对下中心列的指标归纳即得
\[
\gamma_j(D_{2n})=\langle r^{2^{j-1}}\rangle\qquad(j\ge2).
\]
\cref{b1:prop:power-order}说明 \(r^{2^{j-1}}=1\) 当且仅当 \(n\mid2^{j-1}\)，所以该列终止当且仅当 \(n\) 是二的幂。若 \(n=2^a\)（\(a\ge2\)），则 \(\gamma_{a+1}=1\)、\(\gamma_a\ne1\)，幂零类为 \(a\)。所有这些二面体群的导出子群都是循环群，故其二次导出子群平凡，无论是否幂零都可解。
\end{solution}

\begin{exercise}\label{b1:ex:07-central-extension}
若 \(G/Z(G)\) 幂零类至多为 \(c\)，证明 \(G\) 幂零类至多为 \(c+1\)。举例说明把 \(Z(G)\) 换成任意交换正规子群会使结论失效。
\end{exercise}
\begin{solution}
对满射 \(\pi:G\to G/Z(G)\)，式\eqref{b1:eq:lower-central-functoriality}给出
\[
\pi\Bigl(\gamma_{c+1}(G)\Bigr)
=\gamma_{c+1}\Bigl(G/Z(G)\Bigr)=1.
\]
后一等号正是商群幂零类至多为 \(c\) 的假设。因此 \(\gamma_{c+1}(G)\le\ker\pi=Z(G)\)，再与 \(G\) 取交换子，得到 \(\gamma_{c+2}(G)\le\Bigl[G,Z(G)\Bigr]=1\)。反例取 \(G=S_3\)、\(N=A_3\)：\(N\cong C_3\) 与 \(G/N\cong C_2\) 均为交换群，但 \(S_3\) 的三个二阶子群都是 Sylow \(2\) 子群且互不相同，\cref{b1:thm:finite-nilpotent-sylow}因而排除其幂零性。
\end{solution}

\begin{exercise}\label{b1:ex:07-order-p2q}
设 \(p,q\) 是不同素数。确定阶为 \(p^2q\) 的全部幂零群，直到同构。
\end{exercise}
\begin{solution}
由\cref{b1:thm:finite-nilpotent-sylow}，所求群是 Sylow 子群的直积 \(P\times Q\)，其中 \(|P|=p^2,|Q|=q\)。素数阶群循环，故 \(Q\cong C_q\)；\cref{b1:cor:order-p-square-abelian}则说明 \(P\) 交换。若 \(P\) 有 \(p^2\) 阶元，该元素生成全群，所以 \(P\cong C_{p^2}\)。否则，由\cref{b1:cor:finite-order-divides}，每个非单位元都为 \(p\) 阶。取 \(a\ne1\) 及 \(b\notin\langle a\rangle\)，两个 \(p\) 阶循环子群的交只能为 \(1\)：非平凡交会生成双方，从而使它们相等。它们交换，故\cref{b1:prop:internal-direct-product}给出乘积子群同构于 \(C_p\times C_p\)；其阶已为 \(p^2\)，因此等于 \(P\)。答案为 \(C_{p^2}\times C_q\) 与 \(C_p\times C_p\times C_q\)，两者都由有限 \(p\) 群的直积组成，按同一 Sylow 判据确实幂零。两者不同构，因为第一种有 \(p^2\) 阶元，第二种任意元素的阶都整除 \(pq\)。
\end{solution}

\begin{exercise}\label{b1:ex:07-normal-symmetric}
利用交错群单性，证明当 \(n\ge5\) 时，\(S_n\) 的正规子群只有 \(1,A_n,S_n\)。
\end{exercise}
\begin{solution}
设 \(N\trianglelefteq S_n\)，则 \(N\cap A_n\trianglelefteq A_n\)；\cref{b1:thm:alternating-simple}说明这个交为 \(1\) 或 \(A_n\)。若 \(A_n\le N\)，由\cref{b1:thm:sign}和\cref{b1:thm:first-isomorphism}，符号同态给出 \(S_n/A_n\cong C_2\)。再由\cref{b1:thm:subgroup-correspondence}，\(N/A_n\) 只能为这个二阶商群的单位子群或全群，所以 \(N=A_n\) 或 \(S_n\)。若 \(N\cap A_n=1\)，符号同态在 \(N\) 上的核为 \(N\cap A_n\)，故由\cref{b1:prop:kernel-fibers}单射，得到 \(|N|\le2\)。若阶为二，由正规性，唯一非单位元在全部 \(S_n\) 的共轭下固定，因而位于 \(Z(S_n)\)。但 \(Z(S_n)=1\)：若中心元 \(z\) 移动 \(a\) 到 \(b\ne a\)，取不同于 \(a,b\) 的 \(c\)，式\eqref{b1:eq:cycle-conjugation}给出 \(z(ac)z^{-1}=\Bigl(b,z(c)\Bigr)\ne(ac)\)，因为左侧换位移动 \(b\)，右侧不移动 \(b\)，矛盾。因此 \(N=1\)。
\end{solution}

\begin{exercise}\label{b1:ex:07-infinite-unitriangular}
将实单位上三角群依次嵌入为 \(U_1\subset U_2\subset\cdots\)，嵌入在右下角补单位元。证明并群 \(U=\bigcup_m U_m\) 的每个有限生成子群都幂零，但 \(U\) 不可解。
\end{exercise}
\begin{solution}
有限个生成元都落在某个 \(U_m\) 中。\cref{b1:prop:unitriangular-filtration}取 \(a=1,b=k\)，给出 \([U_m,1+J_k]\le1+J_{k+1}\)。以 \(\gamma_1(U_m)=U_m=1+J_1\) 为起点归纳，得到 \(\gamma_k(U_m)\le1+J_k\)，所以 \(\gamma_m(U_m)=1\)，即 \(U_m\) 幂零。式\eqref{b1:eq:lower-central-functoriality}的子群包含式又说明，所给生成元生成的子群也幂零。为证明 \(U\) 不可解，以 \(E_{ij}\) 记只有 \((i,j)\) 位置为一的矩阵。对 \(i<j<k\)，由 \(E_{ij}^2=E_{jk}^2=0\)、\(E_{ij}E_{jk}=E_{ik}\)、\(E_{jk}E_{ij}=0\)，展开四个因子得到
\[
[1+E_{ij},1+E_{jk}]=1+E_{ik}.
\]
归纳可知 \(U^{(r)}\) 含有全部 \(1+E_{i,i+2^r}\)：\(r=0\) 时这些都是 \(U\) 的元素；若结论在第 \(r\) 步成立，则 \(1+E_{i,i+2^r}\) 和 \(1+E_{i+2^r,i+2^{r+1}}\) 都在 \(U^{(r)}\) 中，其交换子由刚算出的恒等式等于 \(1+E_{i,i+2^{r+1}}\)，所以属于 \(U^{(r+1)}\)。对每个 \(r\) 取足够大的矩阵阶数，即可得到这样的非单位元，故每个导出子群非平凡，\(U\) 不可解。
\end{solution}

% !TeX root = ../../Book1.tex
% 纲要标题：## 第8章　自由群与生成关系
% 纲要位置：计划纲要.md，第 328 行。

\chapter{自由群与生成关系}
\label{b1:ch:08}
\BookChapterNumber{b1:ch:08}

% TODO: 撰写本章摘要、学习目标与章导读。

% !TeX root = ../../Book1.tex
% 纲要标题：### 8.1 自由群的构造
% 纲要位置：计划纲要.md，第 330 行。

\section{自由群的构造}
\label{b1:sec:0801}

% TODO: 按以下纲要要点撰写本节。

% 纲要内容（仅作写作计划，尚非教材正文）：
%
% * 约化字及其乘法
% * 自由群的普遍性质
% * 基与秩
%

待撰写。

% !TeX root = ../../Book1.tex
% 纲要标题：### 8.2 群的呈示
% 纲要位置：计划纲要.md，第 336 行。

\section{群的呈示}
\label{b1:sec:0802}

% 纲要内容（仅作写作计划，尚非教材正文）：
%
% * 生成元、关系与正规闭包
% * 从呈示定义同态
% * 循环群、二面体群与多面体群的呈示
%

\subsection{生成元、关系与正规闭包}
\label{b1:subsec:0802:01}

自由群允许任意指定生成元的像。若希望生成元满足额外等式，便须在自由群中把相应的字压成单位。由于同态的核正规，不能只取这些字生成的普通子群，而必须取它们的正规闭包。

\begin{definition}\label{b1:def:group-presentation}
设 \(R\subseteq F(S)\)。包含 \(R\) 的最小正规子群记为 \(\langle\!\langle R\rangle\!\rangle\)，称为 \(R\) 的\term[zhenggui-bibao]{正规闭包}[normal closure]。商群
\[
\langle S\mid R\rangle\coloneqq F(S)/\langle\!\langle R\rangle\!\rangle
\]
称为由生成元 \(S\) 与关系字 \(R\) 给出的\term[qun-chengshi]{群的呈示}[group presentation]。\(S,R\) 都有限时，称为有限呈示。
\end{definition}
\symidx{\(\langle\langle R\rangle\rangle\)：关系集的正规闭包}
正规闭包既可写为包含 \(R\) 的全部正规子群之交，也可写为所有有限乘积
\[
\prod_{i=1}^k u_i r_i^{\varepsilon_i}u_i^{-1},
\qquad u_i\in F(S),\ r_i\in R,\ \varepsilon_i\in\{1,-1\}.
\]
后一集合是子群，在共轭下保持，并包含 \(R\)；任何包含 \(R\) 的正规子群又必须包含这些乘积，所以两种描述一致。等式关系 \(v=w\) 是关系字 \(vw^{-1}\) 的简写。

\ProseParagraphBreak
任何群都有呈示。给定生成集 \(S\)，自由群的泛性质给出满同态 \(F(S)\to G\)；取其核的所有元素为关系字，就有 \(G\cong\langle S\mid R\rangle\)。这并不保证关系集有限，也没有自动提供判别两个字是否相等的算法。

\subsection{呈示与同态构造}
\label{b1:subsec:0802:02}

\begin{proposition}\label{b1:prop:presentation-homomorphism}
从 \(\langle S\mid R\rangle\) 到群 \(H\) 的同态，与满足所有关系 \(r\in R\) 的集合映射 \(f:S\to H\) 一一对应。
\end{proposition}
\begin{proof}
先将 \(f\) 唯一延拓为 \(F(S)\to H\) 的同态。满足关系恰好表示 \(R\) 包含于其核；核正规，等价于 \(\langle\!\langle R\rangle\!\rangle\) 包含于核。因此同态唯一地通过该商群。反过来，商群上的同态在生成元上的限制自然满足每个关系。两个步骤互为逆过程。
\end{proof}
验证一个群的呈示通常需要两方面工作。首先，在具体群 \(G\) 中找出满足关系的生成元，由此得到呈示群 \(P\) 到 \(G\) 的满同态；其次，证明 \(P\) 没有多余元素，或直接构造逆同态。只检验关系成立，只能证明 \(G\) 是 \(P\) 的商群。

\ProseParagraphBreak
例如 \(P=\langle a,b\mid ab=ba\rangle\) 中，每个字可整理为 \(a^mb^n\)。映射 \(a\mapsto(1,0),b\mapsto(0,1)\) 给出 \(P\to\mathbb Z^2\)；反向映射 \((m,n)\mapsto a^mb^n\) 因交换关系而保持加法。两者复合都是恒等，故 \(P\cong\mathbb Z^2\)，且所列标准形唯一。若再加入 \(a^2=b^3=1\)，则得到 \(C_2\times C_3\)，而不是一个含有任意多交错字的非交换群。

\subsection{循环群、二面体群与多面体群的呈示}
\label{b1:subsec:0802:03}

循环群最简单的呈示为
\[
C_n\cong\langle a\mid a^n=1\rangle\qquad(n\ge1).
\]
关系使每个字成为 \(1,a,\ldots,a^{n-1}\) 之一；映到已知的 \(n\) 阶循环群，说明这 \(n\) 个元素彼此不同。无限循环群的呈示则是一个生成元而没有关系。

\begin{proposition}\label{b1:prop:dihedral-presentation}
对 \(n\ge3\)，有
\[
D_{2n}\cong\langle r,s\mid r^n=1,\ s^2=1,\ srs^{-1}=r^{-1}\rangle.
\]
\end{proposition}
\begin{proof}
正 \(n\) 边形的旋转与反射满足所列关系，且生成全部对称，所以存在呈示群到 \(D_{2n}\) 的满同态。呈示中的关系给出 \(sr^k=r^{-k}s\)，可把每个 \(s\) 推到字的最右端，再约去 \(s^2\) 及 \(r^n\)。于是每个元素都具有形式 \(r^i s^\varepsilon\)，其中 \(0\le i<n\)、\(\varepsilon\in\{0,1\}\)，呈示群至多有 \(2n\) 个元素。满像已有 \(2n\) 个元素，故同态必为同构，标准形也唯一。
\end{proof}
多面体的旋转群可以采用类似办法，只是适合的生成元未必只有两个。以正四面体的旋转群 \(A_4\) 为例，取
\[
x=(12)(34),\qquad y=(14)(23),\qquad t=(123).
\]
其呈示为
\[
\left\langle x,y,t\ \middle|\
\begin{gathered}
x^2=y^2=t^3=1,\quad xy=yx,\\
txt^{-1}=y,\quad tyt^{-1}=xy
\end{gathered}\right\rangle.
\]
先核对置换满足关系；\(x,y\) 生成 Klein 四元群，\(t\) 循环置换其三个非单位元，所以它们在 \(A_4\) 中生成十二个元素。反过来，在抽象呈示中用两个共轭关系把 \(t\) 全推到右边，再用阶与交换关系把每个字化为
\[
x^\varepsilon y^\delta t^j,
\qquad \varepsilon,\delta\in\{0,1\},\quad 0\le j<3.
\]
这样至多有十二个元素，满同态到 \(A_4\) 因而为同构。几何上，绕顶点与对面中心的三分之一周转动给出三轮换，绕一对对边中点连线的半周转动给出双换位，故这些确实是四面体的旋转。

\ProseParagraphBreak
同一个群可以有不同的呈示。上例保留 \(V_4\) 的两个生成元，使半直积结构清楚可见；为了追求最少的生成元，可以用 \(y=txt^{-1}\) 消去 \(y\)，但关系会变长。呈示的简短程度、标准形的易算程度及结构的可见程度，往往需要分别考虑。

% !TeX root = ../../Book1.tex
% 纲要标题：### 8.3 自由积
% 纲要位置：计划纲要.md，第 342 行。

\section{自由积}
\label{b1:sec:0803}

% 纲要内容（仅作写作计划，尚非教材正文）：
%
% * 自由积的约化字
% * 自由积的泛性质
% * 直积与自由积的区别
%

\subsection{自由积的约化字}
\label{b1:subsec:0803:01}

直积把两个群放在一起，并强制来自不同因子的元素交换。如果只想保留两个群各自已有的乘法，而不添加跨因子的关系，就应采用自由积。这与自由群相似，但一个字母现在可以是某个因子的任意非单位元。

\begin{definition}\label{b1:def:free-product-word}
设 \(G,H\) 为群，先把它们的非单位元作不相交标记。由这些元素组成的字称为自由积字；若相邻字母总来自不同因子，就称为约化字。空字也约化。
\end{definition}
从所有有限字出发，允许在同一因子内把相邻的 \(x,y\) 换为它们的乘积 \(xy\)，乘积为一时删去这两个字母；也允许这些操作的逆操作。由此生成的等价关系与串接相容。

\begin{lemma}\label{b1:lem:free-product-normal-form}
每个等价类中有唯一约化字。
\end{lemma}
\begin{proof}
从左到右读取字母，始终保存一个约化字。如果新字母与末尾来自不同因子，就接在末尾；若来自相同因子，就把末尾替换为两者乘积，乘积为一时删去末尾。这给出一个与原字等价的约化字。

验证输出不随允许的操作改变。设 \(x,y\) 来自同一因子 \(G\)，比较连续读取 \(x,y\) 与一次读取 \(xy\)。若原字为空或末尾来自 \(H\)，两次读取先添 \(x\) 再合并成 \(xy\)，与直接读取相同。若原末尾为 \(z\in G\)，且 \(zx\ne1\)，两次读取把末尾变为 \((zx)y=z(xy)\)；若 \(xy=1\)，这就恢复 \(z\)。若 \(zx=1\)，第一次删去 \(z\)，前一字母如存在必来自 \(H\)，第二次便接上 \(y\)；直接读取的结果由 \(z(xy)=y\) 给出，也一样。单位乘积的删除包含在这些情形中。交换 \(G,H\) 得到另一因子的同样结论。

因此一次合并或拆分同因子字母不改变算法输出。约化字的输出是自身，所以同一等价类不能包含两个不同约化字。
\end{proof}
等价类在串接下组成群：结合律从串接继承，空字为单位，逆由字母倒序取逆得到。称此群为 \(G,H\) 的\term[ziyou-ji-product]{自由积}[free product]，记为 \(G*H\)。长度为一的字分别给出 \(G,H\) 到 \(G*H\) 的单同态，因而可把原群看作其子群；两者的交只有单位元。
\symidx{\(G*H\)：两个群的自由积}

\ProseParagraphBreak
例如在 \(C_2*C_3\) 中，设 \(a^2=b^3=1\)，则
\[
(ab^2a)(ab)=ab^2b=a,
\]
但 \((ab)^k\) 对每个 \(k\ge1\) 都是长度 \(2k\) 的约化字，所以 \(ab\) 具有无限阶。自由积即使由有限群构成，也通常无限。

\subsection{自由积的泛性质}
\label{b1:subsec:0803:02}

\begin{theorem}\label{b1:thm:free-product-universal}
给定同态 \(f:G\to K\)、\(g:H\to K\)，存在唯一同态 \(u:G*H\to K\)，使它在两个因子上的限制分别为 \(f,g\)。
\end{theorem}
\begin{proof}
把一个字映为各字母像的乘积。同因子字母的合并不改变这个乘积，因为 \(f,g\) 各自保持乘法，所以映射在等价类上良定义，并保持串接乘法。两个因子共同生成 \(G*H\)，故这样的同态唯一。
\end{proof}
自由积的泛性质没有要求 \(f(G)\) 与 \(g(H)\) 交换。例如从 \(C_2*C_3\) 到 \(S_3\)，可以指定 \(a\mapsto(12)\)、\(b\mapsto(123)\)，得到满同态。若再令 \((ab)^2=1\)，商群就同构于 \(S_3\)：该关系给出 \(aba=b^{-1}\)，所以每个字化为 \(b^i a^\varepsilon\)，至多六个元素；相应置换像已经有六个元素。

\ProseParagraphBreak
与自由群相同，泛性质也确定了带指定因子映射的自由积，直到唯一同构。它还给出
\[
F(S\sqcup T)\cong F(S)*F(T)
\]
（这里取不相交副本）。具体地，任意指定 \(S\sqcup T\) 的像，等价于分别指定两个集合的像，再分别延拓为两个自由群上的同态。也可以在两个方向直接发送各个字母，利用生成性检查复合为恒等。

\ProseParagraphBreak
若 \(G=\langle S\mid R\rangle\)、\(H=\langle T\mid U\rangle\)，且生成符号互不相交，则
\[
G*H\cong\langle S\sqcup T\mid R\cup U\rangle.
\]
因为到任意群的同态，只须分别满足两个因子内部的关系。这个描述没有额外加入 \([s,t]=1\)；加入这些关系后才得到直积。

\subsection{直积与自由积的区别}
\label{b1:subsec:0803:03}

\begin{proposition}\label{b1:prop:product-versus-free-product}
自然满同态
\[
\pi:G*H\longrightarrow G\times H,\qquad
g\longmapsto(g,1),\quad h\longmapsto(1,h)
\]
的核是所有 \([g,h]\)（\(g\in G,h\in H\)）的正规闭包。
\end{proposition}
\begin{proof}
记该正规闭包为 \(N\)。跨因子交换子在直积中为单位，故 \(\pi\) 通过 \(Q=(G*H)/N\)。在 \(Q\) 中两个因子的像逐元素交换，因此
\[
G\times H\longrightarrow Q,\qquad
(g,h)\longmapsto \overline g\,\overline h
\]
是同态。它与 \(Q\to G\times H\) 的复合在两个因子上均为恒等，因子又生成相应群，所以两次复合都是恒等。由此 \(Q\cong G\times H\)，得到核的描述。
\end{proof}
当两个因子都非平凡时，选非单位元 \(g\in G,h\in H\)，字 \(ghg^{-1}h^{-1}\) 长四且约化，故自然同态不单射。直积中的相应交换子却必为单位。

\ProseParagraphBreak
特别地，\(C_2*C_2\) 是无限二面体群。设两个二阶生成元为 \(a,b\)，令 \(r=ab,s=a\)。则
\[
srs^{-1}=r^{-1},
\]
而 \(r\) 具有无限阶。每个交错字都能化为唯一的 \(r^k\) 或 \(r^ks\)，其中 \(k\in\mathbb Z\)。它可具体作用在整数直线上：令 \(a\) 作用为 \(m\mapsto-m\)，\(b\) 作用为 \(m\mapsto1-m\)，则 \(r\) 作用为平移 \(m\mapsto m-1\)。这同时展示了反射的几何来源与约化字的代数判别。对比 \(C_2\times C_2\)，后者只有四个元素，两次反射的乘积也只有二阶。

% !TeX root = ../../Book1.tex
% 纲要标题：### 8.4* 自由群的进一步性质
% 纲要位置：计划纲要.md，第 348 行。

\OptionalSection{自由群的进一步性质}{b1:sec:0804}

% 纲要内容（仅作写作计划，尚非教材正文）：
%
% * Nielsen–Schreier 定理的图论证明路线
% * 子群秩与 Schreier 公式的有限指数情形
% * 字问题与几何群论的入口
%
% ---
%

\subsection{Nielsen–Schreier 定理的图论证明}
\label{b1:subsec:0804:01}

本节给出\term[nielsen-schreier-dingli]{Nielsen--Schreier 定理}[Nielsen--Schreier theorem]的组合图论证明。所需的路径只是一串相接的有向边，不需要拓扑空间或基本群的先修知识。图论方法还会实际产生子群的自由基。

\ProseParagraphBreak
先约定图允许重边和环。每条几何边有两个相反方向 \(e,\bar e\)，即使是环也区分方向。一条边路径是首尾相接的有向边序列；插入或删除 \(e\bar e\) 称为消去一次往返。固定顶点 \(o\)，从 \(o\) 回到 \(o\) 的路径在往返消去的等价关系下，以串接为乘法组成群，记为 \(L(\Gamma,o)\)。结合律由串接给出，逆是倒序反向的路径。

\begin{lemma}\label{b1:lem:graph-loop-free}
设 \(\Gamma\) 是连通图，\(T\) 是包含全部顶点的树。为每条不在 \(T\) 中的几何边选一个方向，组成集合 \(E_0\)。从 \(o\) 到顶点 \(v\) 的唯一树路径记为 \(p_v\)。则 \(L(\Gamma,o)\) 是自由群，自由基为
\[
\ell_e=p_v\,e\,p_w^{-1}\qquad
(e\in E_0,\ e:v\longrightarrow w).
\]
\end{lemma}
\begin{proof}
树的两顶点间有唯一路径：存在性来自连通，若有两条不同的无往返路径，取首次分离与再次汇合的区段便得到回路，违背树的无回路性。树内的任何闭路径均可消为常路径；否则有限条边的路径中保留一条无往返闭路径，便给出回路。

把 \(E_0\) 当作自由字母，由泛性质得到 \(F(E_0)\to L(\Gamma,o)\)，字母 \(e\) 映为 \(\ell_e\)。反向定义一个映射：沿闭路径读取边，树边略过，非树边按所选方向读作 \(e\)，反向读作 \(e^{-1}\)，最后约化。插入往返时，若是树边不会读出字母，若是非树边则增加相邻逆字母，因此映射良定义且为同态。

在一个基环路 \(\ell_e\) 上，反向读取只得到字母 \(e\)，故一个方向的复合为恒等。为检查另一方向，设闭路径依次走边 \(e_i:v_{i-1}\to v_i\)，其中 \(v_0=v_k=o\)。在每个相邻位置插入 \(p_{v_i}^{-1}p_{v_i}\)，得到它等价于
\[
\prod_{i=1}^k
\bigl(p_{v_{i-1}}e_i p_{v_i}^{-1}\bigr).
\]
树边对应的括号项可在树中消去；非树边对应基环路或其逆。因此读出字母再送回环路，恰好恢复原来的等价类。两个同态互逆。
\end{proof}
连通图总可选一棵这样的生成树。有限图可以从根开始，每次添加通向新顶点的边。一般无限图在通常的选择公理下，对包含全部顶点的无回路子图用 Zorn 引理取极大者：链的并仍无回路，因为回路只用有限条边；极大者若不连通，可沿原图连接两个分支的路径加进一条跨分支边而不产生回路，矛盾。

\begin{theorem}[Nielsen--Schreier]\label{b1:thm:nielsen-schreier}
自由群的每个子群都是自由群。
\end{theorem}
\begin{proof}
设 \(H\le F(S)\)。构造以右陪集 \(Hg\) 为顶点的\term[schreier-peiji-tu]{Schreier 陪集图}[Schreier coset graph]：对每个顶点 \(Hg\) 和 \(s\in S\)，画一条指定正向为
\[
Hg\xrightarrow{s}Hgs
\]
的几何边，其反向标记为 \(s^{-1}\)。右乘使边与代表元选择无关。即使两条边有同样的端点，也保留其不同身份。任意字从顶点 \(H\) 出发有唯一的带相应字母标签的路径，其终点为该字代表的陪集，所以图连通。

以 \(H\) 为根。闭路径的标签乘积属于 \(H\)，并给出同态 \(L(\Gamma,H)\to H\)。每个 \(h\in H\) 的字都有回到根的路径，故此同态满射。若闭路径的标签在自由群中为单位，则其字可以相邻逆字母消为空字。图中每个顶点都有唯一的给定字母出边，所以路径上相邻逆字母必是同一条边的正反向。字的每次消去因此对应路径的一次往返消去，最终路径也消为空路径，同态单射。

选生成树并应用前一引理，得到 \(H\cong L(\Gamma,H)\cong F(E_0)\)。若树路径 \(p_v\) 的标签字为 \(t_v\)，基元素可显式写为
\[
t_v s\,t_w^{-1}
\]
其中 \(v\xrightarrow{s}w\) 取遍不在树中的正向边。这些字属于 \(H\)，因为两条到达 \(w\) 的路径 \(t_vs,t_w\) 表示相同右陪集。
\end{proof}
证明同时解释了自由群子群为何不会凭空产生新关系：子群中的字关系，转化为图中闭路径的往返消去；树消除了连接不同顶点的冗余路径，剩下的每条边恰好提供一个自由字母。

\subsection{有限指数子群与 Schreier 公式}
\label{b1:subsec:0804:02}

\term[schreier-zhi-gongshi]{Schreier 秩公式}[Schreier rank formula]把有限指数子群的自由基个数化为图的边数与顶点数之差。

\begin{theorem}[Schreier 秩公式]\label{b1:thm:schreier-rank}
设 \(F_r\) 的秩 \(r\) 有限，且子群 \(H\le F_r\) 的指数为有限数 \(d\)。则
\[
\operatorname{rank}H=1+d(r-1).
\]
\end{theorem}
\begin{proof}
Schreier 图有 \(d\) 个顶点，每个顶点对于每个基字母恰有一条正向边，故共有 \(dr\) 条几何边；反向边不重复计数。一棵含全部顶点的树有 \(d-1\) 条边。这一计数可对有限树逐次删去末端顶点归纳得到。非树边数于是为 \(dr-(d-1)\)，而它就是前面所构造自由基的大小。\(r=0\) 时只能有 \(d=1\)，公式仍成立。
\end{proof}
例如 \(F_2=F(\{a,b\})\) 到加法群 \(C_2\) 的同态 \(a\mapsto1,b\mapsto0\)，核 \(H\) 的指数为二。图有顶点 \(H,Ha\)，选从 \(H\) 到 \(Ha\) 的那条正向 \(a\) 边为树。余下正向边分别是 \(Ha\) 到 \(H\) 的另一条 \(a\) 边，以及两个顶点上的 \(b\) 环。树代表字为 \(1,a\)，所以所得自由基为
\[
a^2,\qquad b,\qquad aba^{-1}.
\]
公式也给出 \(\operatorname{rank}H=3\)。它表明非交换自由群子群的秩可以大于母群的秩，不能用向量空间子空间的维数直觉判断。

\subsection{字问题与几何群论}
\label{b1:subsec:0804:03}

给定一个有限生成群的呈示，\term[zi-wenti]{字问题}[word problem]要求判定输入字是否表示单位。对自由群，约化算法读入每个字母时至多进栈或出栈一次，因而在通常按字母计费的模型下，可用与字长成正比的步骤判定。对前面二面体群及四面体群的具体呈示，也可以按已证明的标准形进行判断。

\ProseParagraphBreak
一般呈示的关系闭包没有约化字那样的唯一代表。把关系字及其共轭反复插入、删除，可以验证某些相等式；若无法找到变换过程，不能据此断定两个字不相等，更不能保证搜索总会停止。本节的算法结论只针对已经给出标准形及终止论证的情形，不把它推广到任意有限呈示。

\ProseParagraphBreak
取 \(H=1\) 的 Schreier 图称为这里的\term[cayley-tu]{Cayley 图}[Cayley graph]：顶点为群元素，边记录右乘生成元。自由群的 Cayley 图是一棵树，因为无往返闭路径的标签会是非空约化字，不可能表示单位。给每条边长度一以后，从单位到一个元素的距离恰为其约化字长。于是代数乘法给出的字与图上的路径长度发生联系，这构成几何群论研究的一个基本出发点。更多关于自由群、呈示与组合方法的材料可参阅 Lyndon 与 Schupp 的《Combinatorial Group Theory》\cite{LyndonSchupp2001}。


% TODO: 撰写本章小结、习题与逐题详解。

% !TeX root = ../../Book1.tex
% 纲要标题：## 第9章　有限生成阿贝尔群
% 纲要位置：计划纲要.md，第 356 行。

\chapter{有限生成阿贝尔群}
\label{b1:ch:09}
\BookChapterNumber{b1:ch:09}

\begin{chapterabstract}
有限生成阿贝尔群可以完全分类。本章从整数格出发，证明自由阿贝尔群的子群仍自由，分离挠与无挠部分，再通过整数关系矩阵的 Smith 化简建立不变因子分解。唯一性由各素数的挠层数恢复，不借助后续模论。最后给出商格、指数和坐标变换的具体计算。
\end{chapterabstract}
\begin{learninggoals}
\begin{itemize}
\item 分清整数线性无关、自由基、有限生成与无挠之间的关系。
\item 会把有限阿贝尔群在初等因子与不变因子形式之间转换，并解释唯一性。
\item 能由整数关系矩阵计算商群，保留需要追踪元素时的坐标变换。
\end{itemize}
\end{learninggoals}
本章仅需前面群同态、商群、循环群的基础，以及整数的带余除法、Bézout 等式和线性代数中的行列式。阿贝尔群统一使用加法。一般无限秩子群自由的证明使用良序与选择公理；以有限生成分类为目标时，只需该证明的有限秩部分。有关生成关系与结构定理的伴读可参阅 Dummit 与 Foote 的《Abstract Algebra》\cite{DummitFoote2003}，本章的初等证明不依赖第二卷的模论。

% !TeX root = ../../Book1.tex
% 纲要标题：### 9.1 自由阿贝尔群
% 纲要位置：计划纲要.md，第 358 行。

\section{自由阿贝尔群}
\label{b1:sec:0901}

% 纲要内容（仅作写作计划，尚非教材正文）：
%
% * 整数格与自由阿贝尔群
% * 子群仍自由及有限秩情形
% * 秩及其不变性
%

\subsection{整数格与自由阿贝尔群}
\label{b1:subsec:0901:01}

阿贝尔群的乘法可以改写为加法，生成元之间的关系因而成为整数线性关系。本章统一用 \(0\) 表示单位元，用 \(mx\) 表示元素 \(x\) 的整数倍。先从没有非平凡整数关系的群开始。

\begin{definition}\label{b1:def:free-abelian}
设 \(I\) 为集合，令
\[
\mathbb Z^{(I)}
=\{\,(n_i)_{i\in I}\mid n_i\in\mathbb Z,\ %
n_i\ne0\text{ 的指标只有有限个}\,\}.
\]
逐分量加法使它成为阿贝尔群。与某个 \(\mathbb Z^{(I)}\) 同构的群称为\term[ziyou-abeier-qun]{自由阿贝尔群}[free abelian group]。其中标准向量 \(e_i\) 的对应元素构成一组基：每个元素均可唯一写为有限和 \(\sum n_i e_i\)。
\end{definition}
\symidx{\(\mathbb Z^{(I)}\)：有限支集整数族}
有限的 \(I\) 给出熟悉的整数格 \(\mathbb Z^n\)。无限指标时必须限定有限支集；允许任意整数族所得的全直积是不同的构造，本章不把它当作自由阿贝尔群的定义。

\begin{proposition}\label{b1:prop:free-abelian-universal}
对任意阿贝尔群 \(A\) 和映射 \(f:I\to A\)，存在唯一同态
\(\widetilde f:\mathbb Z^{(I)}\to A\)，使 \(\widetilde f(e_i)=f(i)\)。
\end{proposition}
\begin{proof}
定义 \(\widetilde f(\sum n_i e_i)=\sum n_i f(i)\)。表示的唯一性保证良定义；和只有有限项，交换性保证加法相容。生成元的像已经指定，所以这样的同态唯一。
\end{proof}
这与自由群的泛性质相似，区别在于目标限为阿贝尔群。比如从 \(\mathbb Z^2\) 到群 \(G\) 指定两个基向量的像 \(x,y\)，必须要求 \(xy=yx\)；从 \(F_2\) 出发则没有这个要求。

\ProseParagraphBreak
在 \(\mathbb Z^2\) 中，\((1,0),(1,1)\) 构成一组基，因为
\((u,v)=(u-v)(1,0)+v(1,1)\)，系数唯一。而 \((2,0),(0,1)\) 虽然整数线性无关，却只生成第一坐标为偶数的子群。整数格的“无关”与“生成”必须分别检查，不能把实向量空间中的一组基自动当作整数格的基。

\subsection{自由阿贝尔群的子群}
\label{b1:subsec:0901:02}

\begin{theorem}\label{b1:thm:free-abelian-subgroup}
自由阿贝尔群的每个子群仍是自由阿贝尔群。若 \(H\le\mathbb Z^n\)，则 \(H\) 有一组至多 \(n\) 个元素的基。
\end{theorem}
\begin{proof}
先证明有限秩情形，对 \(n\) 归纳。\(n=0\) 显然。对 \(n\ge1\)，把 \(H\) 投影到最后一个坐标，其像是 \(\mathbb Z\) 的子群，所以为 \(d\mathbb Z\)，其中 \(d\ge0\)。核
\[
K=H\cap(\mathbb Z^{n-1}\times\{0\})
\]
由归纳假设自由且秩至多 \(n-1\)。若 \(d=0\)，则 \(H=K\)。若 \(d>0\)，选 \(h\in H\) 的末坐标为 \(d\)。任意 \(x\in H\) 的末坐标为 \(md\)，于是 \(x-mh\in K\)；而 \(mh\in K\) 迫使 \(md=0\)，进而 \(m=0\)。所以 \(H=K\oplus\mathbb Zh\)，在 \(K\) 的基后加上 \(h\) 即可。

一般情形在通常的选择公理下，把母群的一组基良序为 \((e_\alpha)_{\alpha<\kappa}\)。令 \(F_\alpha\) 由所有 \(e_\beta\)（\(\beta<\alpha\)）生成，令 \(H_\alpha=H\cap F_\alpha\)。从 \(H_{\alpha+1}\) 取 \(e_\alpha\) 的坐标，像为 \(d_\alpha\mathbb Z\)。若 \(d_\alpha>0\)，选 \(h_\alpha\) 的此坐标为 \(d_\alpha\)，完全相同的减去整数倍步骤给出
\[
H_{\alpha+1}=H_\alpha\oplus\mathbb Zh_\alpha.
\]
若 \(d_\alpha=0\)，两项相等。对极限指标 \(\lambda\)，有限支集保证
\(H_\lambda=\bigcup_{\alpha<\lambda}H_\alpha\)。

所选全部 \(h_\alpha\) 生成 \(H\)：逐级使用上式，并在极限步骤取并即可；任何给定元素只涉及有限项。它们也整数线性无关，因为任何有限关系中指标最大的 \(h_\alpha\) 在 \(e_\alpha\) 位置的系数为其整数系数乘 \(d_\alpha\)，更早的项在该位置全为零，所以该系数必须为零；依次向前得到所有系数为零。这便是一组自由基。
\end{proof}
证明并未断言子群的基总能扩充为母群的基。子群 \(2\mathbb Z\le\mathbb Z\) 的基为 \(2\)，却不能成为 \(\mathbb Z\) 的基的一部分；商群在这里出现了二阶挠元。作为可跟算的二维例子，
\[
H=\{\,(u,v)\in\mathbb Z^2\mid u+v\equiv0\pmod3\,\}
\]
在第二坐标的投影为 \(\mathbb Z\)，可选 \(h=(-1,1)\)；核为 \(3\mathbb Z\times\{0\}\)。所以 \((3,0),(-1,1)\) 是 \(H\) 的一组基。

\subsection{秩及其不变性}
\label{b1:subsec:0901:03}

基的个数必须与选择无关，才可用来比较群。

\begin{proposition}\label{b1:prop:free-abelian-rank}
自由阿贝尔群任意两组基的基数相同，称为\term[ziyou-abeier-qun-de-zhi]{自由阿贝尔群的秩}[rank of a free abelian group]。
\end{proposition}
\begin{proof}
若一组基有有限个 \(n\) 元，则每个元素模 \(2A=\{2x\mid x\in A\}\) 后，可唯一把各坐标减为 \(0\) 或 \(1\)，所以 \(|A/2A|=2^n\)。另一组有限基因而具有相同大小。若 \(A\) 有有限生成集，则这些生成元相对于任意一组基只涉及有限多个基元素；其余基元素不可能由它们生成，所以每组基都有限。

对两组无限基 \(B,C\)，把 \(C\) 中的每个元素写为 \(B\) 的有限线性组合。这些支集的并必须覆盖 \(B\)，故 \(|B|\le |C|\)；交换角色得到反向不等式。这与自由群的无限秩证明使用相同的有限支集和基数比较。
\end{proof}
整数线性无关最多只能给出秩的下界。例如前面的指数三子群与 \(\mathbb Z^2\) 都有秩二，但并不相等。另一方面，秩为一的自由阿贝尔群必为无限循环群，因此加法群 \(\mathbb Q\) 若自由就不可能仅凭“嵌在一条实直线中”断言它秩为一；下一节会直接证明它根本不自由。

\ProseParagraphBreak
有限生成阿贝尔群都有从某个 \(\mathbb Z^n\) 出发的满同态。其核由本节定理仍是有限秩自由阿贝尔群，于是关系可以选为有限多个整数向量。这是后面使用有限整数矩阵进行分类的依据；有限生成的关系并非一开始就应当默认为已知。

% !TeX root = ../../Book1.tex
% 纲要标题：### 9.2 挠与无挠
% 纲要位置：计划纲要.md，第 364 行。

\section{挠与无挠}
\label{b1:sec:0902}

% 纲要内容（仅作写作计划，尚非教材正文）：
%
% * 挠子群与商群
% * 有限阿贝尔群的 p 主分解
% * 有限生成条件不可省略的例子
%

\subsection{挠子群与商群}
\label{b1:subsec:0902:01}

整数格中非零向量没有有限阶，但一般阿贝尔群可以同时含有无限阶和有限阶元素。先把后一部分提取出来，可以把后续分类拆成自由方向与有限阶方向。

\begin{definition}\label{b1:def:torsion-subgroup}
阿贝尔群 \(A\) 中满足某个正整数 \(m\) 使 \(mx=0\) 的元素称为\term[naoyuan]{挠元}[torsion element]。全部挠元组成的集合记为
\[
A_{\mathrm{tor}}\coloneqq\{\,x\in A\mid mx=0
\text{ 对某个整数 }m\ge1\,\}.
\]
若 \(A_{\mathrm{tor}}=A\)，称为\term[naoqun]{挠群}[torsion group]；若 \(A_{\mathrm{tor}}=0\)，称为\term[wunao-qun]{无挠群}[torsion-free group]。
\end{definition}
\symidx{\(A_{\mathrm{tor}}\)：挠子群}

\begin{proposition}\label{b1:prop:torsion-quotient}
\(A_{\mathrm{tor}}\) 是 \(A\) 的特征子群，称为\term[naoziqun]{挠子群}[torsion subgroup]；商群 \(A/A_{\mathrm{tor}}\) 无挠。
\end{proposition}
\begin{proof}
若 \(mx=0,ny=0\)，则交换性给出 \(mn(x-y)=0\)，故挠元对差封闭。同态保持整数倍，所以自同构保持挠元，给出特征性。若 \(k(x+A_{\mathrm{tor}})=0\)，则 \(kx\) 为挠元，存在 \(m\ge1\) 使 \(mkx=0\)。于是 \(x\in A_{\mathrm{tor}}\)，原陪集就是零陪集。
\end{proof}
交换性不可随意去掉。在无限二面体群中，两次反射都有二阶，而其乘积可以是无限阶平移。因此一般群的有限阶元素集合未必为子群。

\begin{proposition}\label{b1:prop:fg-abelian-subgroups}
有限生成阿贝尔群的每个子群都有限生成；其中挠子群为有限群。
\end{proposition}
\begin{proof}
取满同态 \(\pi:\mathbb Z^n\to A\)。对子群 \(B\le A\)，原像 \(\pi^{-1}(B)\le\mathbb Z^n\) 由前节定理有有限基，其像生成 \(B\)。特别地，\(A_{\mathrm{tor}}\) 有有限生成元 \(t_1,\ldots,t_s\)。若 \(t_i\) 的阶为 \(m_i\)，每个挠元都能写成
\(\sum a_i t_i\)，其中 \(0\le a_i<m_i\)，故元素数至多 \(\prod m_i\)，为有限群。
\end{proof}
在 \(A=\mathbb Z\oplus C_6\) 中，\((u,\bar v)\) 为挠元当且仅当 \(u=0\)，所以挠子群为 \(C_6\)，无挠商群为 \(\mathbb Z\)。一般有限生成阿贝尔群是否都能像这样拆成直和，仍需下一节的结构定理。

\subsection{有限阿贝尔群的主分解}
\label{b1:subsec:0902:02}

不同素数的挠可以独立拆开。对素数 \(p\)，定义
\[
A_p=\{\,x\in A\mid p^k x=0\text{ 对某个 }k\ge0\,\},
\]
称为 \(A\) 的\term[p-zhufenliang]{\(p\) 主分量}[p-primary component]。它是特征子群，因为两个元素可取同一个较大的 \(p\) 幂来消去，且同态保持其消去关系。
\symidx{\(A_p\)：阿贝尔群的 \(p\) 主分量}

\begin{theorem}\label{b1:thm:abelian-primary}
设 \(A\) 为有限阿贝尔群，\(|A|=p_1^{e_1}\cdots p_t^{e_t}\)。则加法映射给出同构
\[
A_{p_1}\oplus\cdots\oplus A_{p_t}\ \cong\ A.
\]
\end{theorem}
\begin{proof}
若 \(A=0\)，素因子列表为空，空直和为零群，结论成立。以下设 \(A\ne0\)。令 \(N=|A|\)、\(N_i=N/p_i^{e_i}\)。这些 \(N_i\) 的最大公因数为一，故 Bézout 等式给出整数 \(c_i\) 使 \(\sum_i c_iN_i=1\)。对任意 \(x\in A\)，Lagrange 定理给出 \(Nx=0\)，所以
\[
x=\sum_i x_i,\qquad x_i=c_iN_i x,\qquad
p_i^{e_i}x_i=0.
\]
这说明各主分量之和为 \(A\)。

为证唯一性，设 \(\sum_i y_i=0\)，其中 \(y_i\in A_{p_i}\)。每个 \(y_i\) 的阶既为 \(p_i\) 幂又整除 \(N\)，所以 \(p_i^{e_i}y_i=0\)。对等式施以整数倍 \(c_jN_j\)：当 \(i\ne j\) 时，\(N_j\) 含因子 \(p_i^{e_i}\)，故对应项为零；当 \(i=j\) 时，由 Bézout 等式模 \(p_j^{e_j}\) 得 \(c_jN_j\equiv1\)，对应项仍为 \(y_j\)。于是 \(y_j=0\)，对每个 \(j\) 都成立。
\end{proof}
这称为有限阿贝尔群的\term[zhufenjie]{主分解}[primary decomposition]。例如在 \(C_{12}\) 中，对 \(N_2=3,N_3=4\)，可取 \(c_2=-1,c_3=1\)。所以每个 \(\bar x\) 唯一写为
\[
\bar x=(-3\bar x)+(4\bar x),
\]
前项被 \(4\) 消去，后项被 \(3\) 消去，对应 \(C_{12}\cong C_4\oplus C_3\)。这种投影是由群上的整数倍映射给出的，不需要预先选择每个主分量的生成元。

\subsection{有限生成条件与反例}
\label{b1:subsec:0902:03}

有限生成条件的作用至少有两处：它让关系由有限矩阵描述，也让挠子群有限。两个具体例子可以防止把有限生成结构定理误用于一般阿贝尔群。

\ProseParagraphBreak
首先，\((\mathbb Q,+)\) 无挠，但不有限生成。若有限多个有理数的分母都整除 \(N\)，它们的整数线性组合都在 \(\frac1N\mathbb Z\) 中；有理数 \(\frac1{2N}\) 不在其中。它也不是自由阿贝尔群：在 \(\mathbb Q\) 中每个元素都可除以二，故 \(\mathbb Q=2\mathbb Q\)；在任何非零自由阿贝尔群中，一个基元素都不属于二倍子群。这个区别排除了任意有限或无限秩的自由基。

\ProseParagraphBreak
其次，固定素数 \(p\)，考虑
\[
C_{p^\infty}
=\mathbb Z[1/p]/\mathbb Z,\qquad
\mathbb Z[1/p]=\{\,a/p^k\mid a\in\mathbb Z,\ k\ge0\,\}.
\]
这里 \(\mathbb Z[1/p]\) 只作为有理数加法的子群使用。商群中 \(\frac1{p^k}+\mathbb Z\) 的阶恰为 \(p^k\)，所以它是无限挠群。有限多个元素都在某一个由 \(\frac1{p^k}+\mathbb Z\) 生成的有限循环群中，因而整个群不有限生成。这个群称为\term[prufer-qun]{Prüfer 群}[Prüfer group]。

\ProseParagraphBreak
上述两个群还共享一种现象：有限生成的小部分很简单，全部小部分的并却未必具有同样的有限描述。因此后续每次使用“有限个不变因子”或“挠子群有限”时，都应保留有限生成这一条件。下一节证明的无挠推出自由，也明确限于有限生成阿贝尔群。

% !TeX root = ../../Book1.tex
% 纲要标题：### 9.3 结构定理
% 纲要位置：计划纲要.md，第 370 行。

\section{结构定理}
\label{b1:sec:0903}

% 纲要内容（仅作写作计划，尚非教材正文）：
%
% * 不变因子分解与初等因子分解
% * 结构定理的初等证明
% * 分解的唯一性
%

\subsection{不变因子分解与初等因子分解}
\label{b1:subsec:0903:01}

本节的目标是把任意有限生成阿贝尔群化为循环群的有限直和，而且给出可由群本身恢复的分类数据。有限循环因子有两种常用排列方式：按整除关系排列，或按素数幂拆开。

\begin{theorem}\label{b1:thm:fga-structure}
每个有限生成阿贝尔群 \(A\) 都有分解
\[
A\cong\mathbb Z^r\oplus C_{d_1}\oplus\cdots\oplus C_{d_s},
\qquad 1<d_1\mid d_2\mid\cdots\mid d_s,
\]
其中 \(r,s\ge0\)。整数 \(r\) 及序列 \(d_1,\ldots,d_s\) 由 \(A\) 唯一确定。也可以唯一地写为一个自由部分与若干素数幂阶循环群的直和，唯一性允许重新排列这些循环因子。
\end{theorem}
第一种称为\term[bubian-yinzi-fenjie]{不变因子分解}[invariant factor decomposition]，其中 \(d_i\) 为不变因子；第二种称为\term[chudeng-yinzi-fenjie]{初等因子分解}[elementary divisor decomposition]，各循环因子的素数幂阶称为初等因子。下文依次证明存在性与唯一性。

\ProseParagraphBreak
例如 \(C_4\oplus C_6\) 还没有按整除关系排列，因为 \(4\nmid6\)。先按素数拆开，
\[
C_4\oplus C_6\cong C_4\oplus C_2\oplus C_3
\cong C_2\oplus C_{12}.
\]
初等因子是 \(4,2,3\)，不变因子是 \(2,12\)。这两种写法表达同一个群，但各自必须遵守相应的归一化条件；不能在任意循环因子拆分之间直接谈唯一性。

\subsection{结构定理的初等证明}
\label{b1:subsec:0903:02}

生成元与关系构成一个整数矩阵。要整理关系，只能用可逆的整数行、列操作，不能像实数线性代数那样任意除以主元；除法会改变整数格。

\begin{lemma}[整数矩阵的 Smith 化简]\label{b1:lem:smith-integer}
对任意 \(m\times n\) 整数矩阵 \(M\)，可以通过交换两行或两列、把一行或一列乘 \(-1\)、把某行或列的整数倍加到另一行或列，化为矩形对角阵
\[
\operatorname{diag}(d_1,\ldots,d_t,0,\ldots),
\qquad 0<d_1\mid d_2\mid\cdots\mid d_t.
\]
\end{lemma}
\begin{proof}
若矩阵全零则结束。否则把一个非零元素移到左上角，并使其为正，记为 \(a\)。对第一列的某个元素 \(b\)，作整数除法 \(b=qa+r\)，其中 \(0\le r<a\)，再把第一行的 \(q\) 倍从该行减去。如果 \(r>0\)，交换该行与第一行，使正主元严格变小；若 \(r=0\)，就消去了该元素。对第一行作同样的列操作。每当出现非零余数，主元严格下降，因此这种下降只可能发生有限次；主元不下降时，有限步内可把第一行和第一列的其余元素都清零。

此时矩阵形如
\[
\begin{pmatrix}a&0\\0&B\end{pmatrix}.
\]
如果 \(a\) 不整除 \(B\) 中的某个元素 \(b_{ij}\)，把该元素所在的行加到第一行。左上角仍为 \(a\)，而第一行现在含有 \(b_{ij}\)。再次用列的整数除法就产生小于 \(a\) 的非零余数，使主元严格下降，然后重新清理首行首列。正整数不能无限下降，所以最终获得这样的分块形式，且 \(a\) 整除 \(B\) 的全部元素。

固定首行首列，对 \(B\) 递归进行相同步骤。行列数都减少，递归必停。由于 \(B\) 的所有元素一开始均为 \(a\) 的倍数，后续整数行列操作保持此性质，所以它所得的非零对角元也都是 \(a\) 的倍数。再用递归中的整除链，便得到全部 \(d_i\mid d_{i+1}\)。所用每一步操作都有整数逆操作，没有改变相应整数格的可逆坐标性质。
\end{proof}
上述结果称为整数矩阵的\term[smith-biaozhunxing]{Smith 标准形}[Smith normal form]。证明给出的步骤是一个终止算法，而不只是允许“选取适当基”的存在性叙述。

\begin{proof}[结构定理的存在性证明]
取 \(A\) 的有限生成元，得到满同态 \(\pi:\mathbb Z^m\to A\)。由\cref{b1:thm:free-abelian-subgroup}，核 \(K\) 有有限基 \(k_1,\ldots,k_t\)，其中 \(t\le m\)。把它们在标准基下的坐标作为列，组成 \(m\times t\) 整数矩阵 \(M\)。

整数可逆列操作只改变 \(K\) 的生成基，因而不改变其生成子群；整数可逆行操作是母群 \(\mathbb Z^m\) 的自同构，因而把商群变成同构的商群。对 \(M\) 作 Smith 化简后，核变成由
\[
d_1e_1,\ldots,d_te_t
\]
生成的子群。这里恰有 \(t\) 个非零对角元：原列在整数上无关，也就在有理数上无关，因为有理数关系乘以公共分母就成为整数关系；可逆行列操作保持这个性质。于是
\[
A\cong
\mathbb Z^m/K
\cong C_{d_1}\oplus\cdots\oplus C_{d_t}\oplus\mathbb Z^{m-t}.
\]
等于一的对角元只贡献平凡群，删去它们便得到不变因子形式。

最后把每个 \(d_i\) 作素因数分解。对互素的 \(u,v\)，映射
\(C_{uv}\to C_u\oplus C_v\)、\(\bar x\mapsto(\bar x,\bar x)\)
的核平凡，且两边元素数相同，所以为同构。反复应用便把每个有限循环因子拆成素数幂阶循环群，得到初等因子形式。
\end{proof}
证明同时得到：有限生成无挠阿贝尔群自由；一般有限生成阿贝尔群的挠子群可以分离成一个直和因子。不过具体选出的自由补群依赖基的选择，结构定理没有宣称这个补群唯一。

\subsection{分解的唯一性}
\label{b1:subsec:0903:03}

先恢复自由部分，再分别恢复每个素数的幂阶。挠子群 \(T=A_{\mathrm{tor}}\) 是由群内条件确定的，所以同构必把它映到另一个群的挠子群。上面的分解说明 \(A/T\cong\mathbb Z^r\)，于是\cref{b1:prop:free-abelian-rank}确定了 \(r\)。

\ProseParagraphBreak
对有限群 \(T\) 的 \(p\) 主分量 \(B=T_p\)，考虑全由整数倍定义的子群及商群
\[
B\supseteq pB\supseteq p^2B\supseteq\cdots,\qquad
p^{j-1}B/p^jB.
\]
若 \(B\cong\bigoplus_{i=1}^k C_{p^{e_i}}\)，每个因子对第 \(j\) 层的贡献在 \(e_i\ge j\) 时有 \(p\) 个元素，在 \(e_i<j\) 时只有一个。因此
\[
\bigl|p^{j-1}B/p^jB\bigr|=p^{c_j},
\qquad c_j=\#\{\,i\mid e_i\ge j\,\}.
\]
左边是由 \(B\) 本身确定的数，所以每个 \(c_j\) 都是同构不变量。恰好等于 \(j\) 的指数出现 \(c_j-c_{j+1}\) 次，因而全部 \(e_i\) 及其重数唯一。这证明初等因子的唯一性，没有引入模或张量积。

\ProseParagraphBreak
由初等因子恢复不变因子也没有歧义。若 \(T=0\)，有限循环因子为空，取 \(s=0\) 即可。以下设 \(T\ne0\)。对每个素数 \(p\)，把正指数按非降次序排列；令 \(s\) 为各素数指数列表长度的最大值，把较短列表的左端补零至长度 \(s\)。逐位置把相应的 \(p\) 幂相乘，便得到
\[
d_1\mid d_2\mid\cdots\mid d_s.
\]
反过来，任何满足整除链且 \(d_1>1\) 的不变因子分解，其每个素数的指数必非降，去掉左端的零以后恰为该素数的初等因子指数。至少有一个素数整除 \(d_1\)，所以某个列表的长度恰为 \(s\)，不会额外留下全为零的位置。由此不变因子也唯一，\cref{b1:thm:fga-structure}得证。

\ProseParagraphBreak
例如若一个二主群的各层商群阶依次为 \(2^3,2^2,2^2\)，再下一层为一，则指数至少为 \(1,2,3\) 的因子分别有 \(3,2,2\) 个。因此指数一出现一次，指数二不出现，指数三出现两次，该群为 \(C_2\oplus C_8\oplus C_8\)。只知道群的阶为 \(2^7\) 无法作出这一判断，逐层取整数倍才恢复了各个循环因子的长度。

% !TeX root = ../../Book1.tex
% 纲要标题：### 9.4 结构定理的计算与模论解释
% 纲要位置：计划纲要.md，第 383 行。

\section{结构定理的计算与模论解释}
\label{b1:sec:0904}

结构定理给出了分类形式，实际计算还须把生成元、关系和坐标变换一一对应。下面先固定关系矩阵的行列约定，再用一个完整例子说明怎样从化简结果读出循环因子并追踪具体元素。随后讨论整数格的商群、子格指数与子式判据。节末的模论解释只是对第二卷的预告，本节全部计算仍只依赖整数矩阵、子群与商群。

% 纲要内容（仅作写作计划，尚非教材正文）：
%
% * 由整数关系矩阵确定交换群
% * 整数格的商群与有限指数
% * 模论解释与后续推广预告
%
% ---
%

\subsection{整数关系矩阵与交换群}
\label{b1:subsec:0904:01}

使用结构定理计算时，第一步是固定矩阵的方向。本节一律把生成元列为坐标轴，把一条整数关系写成一个列向量。给定 \(m\times n\) 整数矩阵 \(M\)，可以用它的各列施加关系，并定义交换群
\[
A_M\coloneqq\mathbb Z^m/M\mathbb Z^n.
\]
这里 \(M\mathbb Z^n\) 是由这些列关系的全部整数线性组合组成的子群。列操作改变关系生成集，行操作改变群的生成坐标；两者都必须是可逆的整数操作。

\ProseParagraphBreak
若要用已知关系分析一个事先给定的交换群 \(A\)，还须核对关系是否完整。设 \(A\) 由 \(a_1,\ldots,a_m\) 生成，并令
\[
\pi:\mathbb Z^m\longrightarrow A,
\qquad e_i\longmapsto a_i
\]
为自然满同态。若 \(M\) 的各列只记录若干已经验证的关系，则只能得到 \(M\mathbb Z^n\subseteq\ker\pi\)，因而 \(\pi\) 诱导自然满同态
\[
\overline\pi:\mathbb Z^m/M\mathbb Z^n\longrightarrow A.
\]
只有进一步证明 \(M\mathbb Z^n=\ker\pi\)，这个满同态才是同构，Smith 化简所得结构才是 \(A\) 本身的分类。例如 \(A=C_2=\langle a\rangle\) 确实满足关系 \(4a=0\)，但只记录这一关系得到的是 \(\mathbb Z/4\mathbb Z\cong C_4\)，它只自然满射到 \(C_2\)；遗漏的关系 \(2a=0\) 正说明所列关系并不完整。

\ProseParagraphBreak
Smith 计算本身的输入是一张有限整数矩阵 \(M\)。在 \(\mathbb Z\) 上逐次作可逆行列操作，并同时累计这些操作，输出 \(m\times m\) 与 \(n\times n\) 整数可逆矩阵 \(U,V\)，以及 Smith 标准形 \(D=UMV\)。只求同构类型时，从 \(D\) 读取自由秩与非平凡不变因子即可；还要追踪原生成元或具体元素时，则须保留 \(U,V\) 及其逆矩阵。\cref{b1:lem:smith-integer}的证明给出了这一过程的终止性，但本章不讨论它的复杂度。若 \(M\) 用来分析一个事先给定的群，求得生成完整关系核的矩阵是 Smith 化简之前另须解决的问题。

\ProseParagraphBreak
设 Smith 标准形有 \(t\) 个非零对角元 \(d_1\mid\cdots\mid d_t\)。\cref{b1:tab:smith-coordinate-types}区分新坐标中的三种情况：前 \(t\) 个坐标受到非零关系约束，其余 \(m-t\) 个坐标没有这样的约束。

\begin{table}[htbp]
\centering
\caption{Smith 标准形中的坐标与商群因子}
\label{b1:tab:smith-coordinate-types}
\begin{tabularx}{\linewidth}{@{}>{\raggedright\arraybackslash}p{0.44\linewidth}>{\raggedright\arraybackslash}X@{}}
\toprule
新坐标中的情形 & 对商群的贡献 \\
\midrule
非单位主元 \(d_i>1\) & 非平凡有限循环因子 \(C_{d_i}\) \\
单位主元 \(d_i=1\) & 该坐标在商群中为零，平凡因子 \(C_1\) 删去 \\
没有非零主元约束的坐标 & 每个贡献一个无限循环因子，共有 \(m-t\) 个 \\
\bottomrule
\end{tabularx}
\end{table}

矩阵的非零 Smith 对角元允许为 \(1\)；群结构定理所列的不变因子则删去这些 \(1\)，只保留大于 \(1\) 的项。标准形中的零列表示经过可逆整数重组后出现的零关系，可以在这一阶段删去；这不意味着原关系组中必有某条关系可以直接删除。关系矩阵可以有任意多个相互依赖的列；自由部分的秩由目标格的坐标数减去非零主元数得到，为 \(m-t\)，不是 \(n-t\)。

\ProseParagraphBreak
例如对一生成元群施加关系 \(2x=0\)、\(3x=0\)，关系矩阵为 \(M=\begin{pmatrix}2&3\end{pmatrix}\)。取整数可逆矩阵
\[
V=\begin{pmatrix}-1&3\\1&-2\end{pmatrix},
\qquad \det V=-1,
\]
便有
\[
MV=\begin{pmatrix}1&0\end{pmatrix}.
\]
零列是在两个原关系经过可逆重组后产生的；原关系中却没有一条可以直接删去，因为两条关系合起来定义平凡群，只保留 \(2x=0\) 或只保留 \(3x=0\) 时分别得到 \(C_2\) 或 \(C_3\)。

\paragraph*{求同构类型：读取 Smith 对角元}
例如群由 \(x,y\) 生成，关系为 \(4x+2y=0\)、\(6x+8y=0\)，矩阵为
\[
M=\begin{pmatrix}4&6\\2&8\end{pmatrix}.
\]
一次完整化简为
\[
\begin{pmatrix}4&6\\2&8\end{pmatrix}
\longrightarrow
\begin{pmatrix}2&8\\4&6\end{pmatrix}
\longrightarrow
\begin{pmatrix}2&8\\0&-10\end{pmatrix}
\longrightarrow
\begin{pmatrix}2&0\\0&-10\end{pmatrix}
\longrightarrow
\begin{pmatrix}2&0\\0&10\end{pmatrix}.
\]
四步分别为交换两行、第二行减第一行的两倍、第二列减第一列的四倍、第二行变号。因此 \(A\cong C_2\oplus C_{10}\)。

\paragraph*{追踪原元素：保存行变换矩阵}
若还需要指出商群元素怎样对应，不能只记住对角数字。本例行操作的总矩阵为
\[
U=\begin{pmatrix}0&1\\-1&2\end{pmatrix},\qquad
U^{-1}=\begin{pmatrix}2&-1\\1&0\end{pmatrix}.
\]
原坐标向新坐标转换使用 \(x\mapsto Ux\)；若要把新坐标轴对应的生成元写回原坐标，则使用 \(U^{-1}e_i\)。先在整数坐标上定义同态
\[
\psi:\mathbb Z^2\longrightarrow C_2\oplus C_{10},\qquad
\psi(u,v)=\Bigl([v]_2,[-u+2v]_{10}\Bigr).
\]
前述列操作的矩阵及其整数逆矩阵为
\[
V=\begin{pmatrix}1&-4\\0&1\end{pmatrix},\qquad
V^{-1}=\begin{pmatrix}1&4\\0&1\end{pmatrix}.
\]
实际相乘得到 \(UMV=\operatorname{diag}(2,10)\)。因为 \(V\mathbb Z^2=\mathbb Z^2\)，有 \(U(M\mathbb Z^2)=2\mathbb Z\oplus10\mathbb Z\)。因此
\[
\ker\psi=U^{-1}(2\mathbb Z\oplus10\mathbb Z)=M\mathbb Z^2.
\]
\(U\) 在整数上可逆，随后分别取模的映射满射，所以 \(\psi\) 满射。

\ProseParagraphBreak
由\cref{b1:thm:first-isomorphism}，\(\psi\) 诱导同构 \(\mathbb Z^2/M\mathbb Z^2\rightarrow C_2\oplus C_{10}\)，把陪集 \((u,v)+M\mathbb Z^2\) 送到 \(\psi(u,v)\)。原坐标中的 \(U^{-1}e_1=(2,1)\) 对应 \(C_2\) 的标准生成元，\(U^{-1}e_2=(-1,0)\) 对应 \(C_{10}\) 的标准生成元；它们在商群中的阶分别为二与十。

\ProseParagraphBreak
若生成元为 \(x,y,z\)，关系只有 \(2x=0\)、\(x+3y=0\)，则先用第二条消去 \(x=-3y\)，第一条成为 \(6y=0\)，而 \(z\) 没有关系。因此群为 \(C_6\oplus\mathbb Z\)。没有非零对角关系约束的位置贡献自由生成元，而不是有限阶循环因子。

\subsection{整数格的商群与有限指数}
\label{b1:subsec:0904:02}

\begin{proposition}[子格指数与商群结构]\label{b1:prop:lattice-index}
设 \(m,n\) 为正整数，\(M\) 为 \(m\times n\) 整数矩阵，\(L=M\mathbb Z^n\le\mathbb Z^m\)。若 \(M\) 的 Smith 标准形有 \(t\) 个非零对角元 \(d_1,\ldots,d_t\)，则
\[
\mathbb Z^m/L\cong\mathbb Z^{m-t}\oplus
\bigoplus_{i=1}^t C_{d_i}.
\]
特别地，\(L\) 的指数有限当且仅当 \(t=m\)，此时指数为 \(\prod_i d_i\)。若 \(M\) 为非奇异的 \(m\times m\) 矩阵，则指数为 \(|\det M|\)。
\end{proposition}

\begin{proof}
由\cref{b1:lem:smith-integer}，存在整数可逆矩阵 \(U,V\) 使 \(UMV=D\)，其中 \(D\) 是所述对角阵。由于 \(V\mathbb Z^n=\mathbb Z^n\)，有 \(U(L)=D\mathbb Z^n=\langle d_1e_1,\ldots,d_te_t\rangle\)。因此 \(x+L\mapsto Ux+U(L)\) 给出商群同构。再把前 \(t\) 个坐标分别模 \(d_i\)，其余坐标保留，所得满同态的核恰为 \(U(L)\)；\cref{b1:thm:first-isomorphism}遂给出所述商群公式。

指数就是商群的元素数。当 \(m-t>0\) 时有无限循环直和因子，故指数无限；当 \(m=t\) 时指数为有限循环因子阶的乘积。若 \(M\) 是非奇异方阵，则 \(t=m\)，而 \(\det U,\det V\in\{1,-1\}\)，所以从 \(UMV=D\) 得 \(|\det M|=\det D=\prod_i d_i\)。
\end{proof}

例如上一小节的二维关系子格指数为二十，等于 \(|4\cdot8-6\cdot2|\)。但行列式只给出总阶，并不决定商群：对角矩阵 \(\operatorname{diag}(1,20)\) 与 \(\operatorname{diag}(2,10)\) 的行列式同为二十，商群分别为 \(C_{20}\) 与 \(C_2\oplus C_{10}\)，两者中元素的最大阶不同。

\begin{proposition}[行列式因子]\label{b1:prop:determinantal-divisors}
设 \(m,n\) 为正整数，\(M\) 为 \(m\times n\) 整数矩阵。对 \(1\le k\le\min(m,n)\)，令 \(\Delta_k(M)\) 为 \(M\) 的全部 \(k\) 阶子式的非负最大公因数，并将它称为第 \(k\) 个\term[xinglieshi-yinzi]{行列式因子}[determinantal divisor]；全部相应子式都为零时约定其最大公因数为零，并令 \(\Delta_0(M)=1\)。令
\[
t=\operatorname{rank}_{\mathbb Q}M
\]
为 \(M\) 在有理数域上的矩阵秩。这个数也等于像子群 \(M\mathbb Z^n\) 的自由秩，以及 Smith 标准形中非零对角元的个数。将这些非零对角元依次记为 \(d_1,\ldots,d_t\)，则
\[
\begin{aligned}
\Delta_k(M)&=d_1\cdots d_k &&(1\le k\le t),\\
\Delta_k(M)&=0 &&(t<k\le\min(m,n)).
\end{aligned}
\]
因此比值公式为
\[
d_k=\frac{\Delta_k(M)}{\Delta_{k-1}(M)}\qquad(1\le k\le t).
\]
\end{proposition}

\symidx{\(\Delta_k(M)\)：矩阵的行列式因子}

\begin{proof}
整数可逆矩阵在有理数域上也可逆，所以 Smith 化简 \(UMV=D\) 保持 \(\operatorname{rank}_{\mathbb Q}\)，而 \(D\) 的有理秩就是其非零对角元个数。另一方面，\(U\) 把像子群 \(M\mathbb Z^n\) 同构地送到由 \(d_1e_1,\ldots,d_te_t\) 生成的自由交换群，所以该像子群的自由秩也等于 \(t\)。

整数初等行列操作保持各阶子式的最大公因数。以某行加另一行的整数倍为例，每个受影响子式由行列式的线性性写成原子式与另一原子式的整数线性组合；如果所加行已包含在子式中，额外项有重复行而为零。逆操作给出反向整除，所以最大公因数不变。换行、变号及列操作同样成立。

对 Smith 对角阵，当 \(1\le k\le t\) 时，非零 \(k\) 阶子式是所选 \(k\) 个对角元的乘积。若指标为 \(i_1<\cdots<i_k\)，则 \(i_j\ge j\)，整除链使 \(d_1\cdots d_k\) 整除该乘积；而前 \(k\) 个对角元的子式恰好等于此乘积。因此它就是全部子式的最大公因数。各项为正，结合 \(\Delta_0(M)=1\)，便得到所写的比值公式。

若 \(k>t\)，对角阵的任意 \(k\) 阶子式都为零，故 \(\Delta_k(M)=0\)。特别地，零矩阵的所有正阶行列式因子均为零，不适用上述比值公式。
\end{proof}

\begin{corollary}[满行秩矩形矩阵的子格指数]\label{b1:cor:full-row-rank-index}
设 \(m,n\) 为正整数，\(M\) 为满行秩的 \(m\times n\) 整数矩阵，即 \(\operatorname{rank}_{\mathbb Q}M=m\)。则 \(m\le n\)，并且
\[
[\mathbb Z^m:M\mathbb Z^n]=\Delta_m(M).
\]
即子格 \(M\mathbb Z^n\) 在 \(\mathbb Z^m\) 中的指数等于全部 \(m\) 阶子式绝对值的最大公因数。若 \(m=n\)，这就是 \(|\det M|\)。
\end{corollary}

\begin{proof}
满行秩意味着非零 Smith 对角元的个数 \(t=m\)。又因 \(t\leq\min(m,n)\)，必有 \(m\leq n\)。由\cref{b1:prop:lattice-index}，指数为 \(d_1\cdots d_m\)；再由\cref{b1:prop:determinantal-divisors}，此乘积等于 \(\Delta_m(M)\)。方阵只有一个 \(m\) 阶子式，即其行列式。
\end{proof}

这里必须是满行秩，而不能笼统地只说矩形矩阵“满秩”。若 \(n<m\) 且 \(M\) 满列秩，则 \(t=n<m\)，商群仍有自由秩 \(m-n\)，因而子格指数无限；此时全部最大非零子式是 \(n\) 阶子式，它们的最大公因数并不表示 \(M\mathbb Z^n\) 在 \(\mathbb Z^m\) 中的有限指数。

\ProseParagraphBreak
前面用于说明零列的矩阵 \(M=\begin{pmatrix}2&3\end{pmatrix}\) 有满行秩，而
\[
M\mathbb Z^2=2\mathbb Z+3\mathbb Z=\mathbb Z,
\]
因为 \(3-2=1\)。其指数为 \(1=\gcd(2,3)\)，不能任选一个非零最大阶子式 \(2\) 或 \(3\) 作为指数。

\ProseParagraphBreak
行列式因子也适合核验小矩阵的 Smith 结果。本节的二维算例中，全部元素的最大公因数为二，唯一二阶子式为二十，因此 \(d_1=2,d_2=10\)。大矩阵的子式很多，可以先做整数行列化简，再抽查少量子式以发现计算失误。不过，若所抽取子式的最大公因数为 \(g\)，只能直接得到 \(\Delta_k(M)\mid g\)，不能据此断言相等。例如 \(\operatorname{diag}(2,3)\) 的全部一阶子式最大公因数为 \(1\)，只检查条目 \(2\) 就不能恢复它。

\ProseParagraphBreak
若要独立确认正的候选值 \(\delta=\Delta_k(M)\)，须证明全部 \(k\) 阶子式均被 \(\delta\) 整除，并找到最大公因数恰为 \(\delta\) 的一组子式。这两个方向合起来才给出等号。另一种充分验证是完整保存 \(U,V\)，核验 \(UMV=D\) 及 \(U,V\) 的整数逆矩阵，并检查 \(D\) 的非零对角元为正且形成整除链；这样便确认了 Smith 结果，无须枚举全部子式。

\subsection{模论解释与后续推广预告}
\label{b1:subsec:0904:03}

本章始终只用整数倍、子群与商群。下一卷将把这种结构重新解释为模：整数 \(n\) 对交换群元素 \(x\) 的作用就是 \(nx\)。这个预告不作为当前证明的输入；模的定义与公理将在第二卷第二章“模、子模与商模”正式建立。

\ProseParagraphBreak
从这一后续观点看，\(\mathbb Z^m\) 是有限秩自由对象，关系矩阵记录了从一个自由对象到另一个自由对象的映射，商群记录生成元在关系约束下留下的信息。整数的任意理想由一个元素生成，这一性质将被抽象为主理想整环条件，从而把本章的循环分解推广到更一般的系数系统。在本章，支持消元的具体工具是整数除法以及 Bézout 等式；推广时必须重新核对系数系统是否提供相应工具。

\ProseParagraphBreak
例如对有理数系数多项式也可以作带余除法，第二卷第四章“主理想整环上的模与标准形”会用类似方法研究线性变换的有理标准形。届时矩阵中的“整数关系”变为“多项式关系”，循环因子表达一个向量及其在反复作用下的生成规律。这个类比解释了有限生成交换群为什么安排在群论之后：它既完成一项独立分类，也提供了从群的生成关系过渡到模的生成关系的具体经验。

\ProseParagraphBreak
实际使用本章结果时，可按三步整理资料。第一步，确认群交换且有限生成，选择有限生成元，并取得一张列向量生成整个关系核的有限整数矩阵。第二步，明确行列约定，在整数上作可逆变换，同时累计 \(U,V\)，直至得到 \(D=UMV\)。第三步，核验 \(U,V\) 都有整数逆矩阵，且 \(D\) 的非零对角元为正并形成整除链，再从 \(D\) 读取自由秩及非平凡不变因子。

\ProseParagraphBreak
有限秩子群定理只保证完整关系核存在有限基，并不自动给出可供计算的生成组。若矩阵只记录部分已知关系，只能得到相应展示群到原群的自然满同态，不能把展示群的分类结果直接当作原群的分类。若需要追踪特定元素或子群，应使用保存的坐标变换；若只需判断两个有限生成交换群是否同构，比较自由秩与归一化后的不变因子即可。


\begin{chaptersummary}
有限生成阿贝尔群的关系核是有限秩自由阿贝尔群，因此可用有限整数矩阵表示。Smith 化简只作可逆整数行列操作，得到自由部分与按整除次序排列的循环因子。自由秩由无挠商群确定，每个素数幂因子的长度由 \(p^{j-1}A_p/p^jA_p\) 的阶确定。矩阵行列式给出满秩商格的指数，各阶子式的最大公因数进一步恢复全部不变因子；只有总阶通常不足以区分群。
\end{chaptersummary}
\BookExercises

\begin{exercise}\label{b1:ex:09-congruence-lattice}
求子群
\[
H=\{\,(x,y,z)\in\mathbb Z^3\mid
x+2y+3z\equiv0\pmod6\,\}
\]
的一组基，以及 \(\mathbb Z^3/H\) 的结构。
\end{exercise}
\begin{solution}
每个元素唯一写为
\[
k(6,0,0)+y(-2,1,0)+z(-3,0,1),
\qquad k=(x+2y+3z)/6\in\mathbb Z.
\]
因此所列三个向量为一组基。同态
\(\mathbb Z^3\to C_6\)、\((x,y,z)\mapsto\overline{x+2y+3z}\)
满射且核为 \(H\)，故商群为 \(C_6\)，指数为六。
\end{solution}

\begin{exercise}\label{b1:ex:09-explicit-quotient}
求 \(A=\mathbb Z^2/\langle(6,0),(4,10)\rangle\) 的不变因子，并在原坐标中选出对应两个循环直和因子的生成元。
\end{exercise}
\begin{solution}
关系矩阵为 \(\begin{pmatrix}6&4\\0&10\end{pmatrix}\)。一阶子式最大公因数为二，行列式为六十，所以不变因子为 \(2,30\)。令 \(u=(3,0)+L\)、\(v=(0,1)+L\)，其中 \(L\) 为关系格。\(u\) 的阶为二；若 \(kv=0\)，则
\[
(0,k)=a(6,0)+b(4,10),
\]
所以 \(k=10b\)、\(6a+4b=0\)，迫使 \(3\mid b\)，最小正 \(k\) 为三十。若 \(u\) 属于 \(\langle v\rangle\)，某个关系向量的第一坐标应为三，但所有关系向量的第一坐标均为偶数，不可能。故两个循环子群交为零，乘积有六十个元素，等于全群，给出 \(A=\langle u\rangle\oplus\langle v\rangle\)。
\end{solution}

\begin{exercise}\label{b1:ex:09-order-360}
列出所有阶为 \(360\) 的阿贝尔群，写成不变因子形式，并说明没有遗漏或重复。
\end{exercise}
\begin{solution}
\(360=2^3\cdot3^2\cdot5\)。二主分量对应三的分拆 \(3,\ 2+1,\ 1+1+1\)，三主分量对应二的分拆 \(2,\ 1+1\)，五主分量只能为 \(C_5\)。共六种组合，按素数指数左端补零并逐列相乘，得到
\[
\begin{gathered}
C_{360},\qquad C_3\oplus C_{120},\qquad
C_2\oplus C_{180},\\
C_6\oplus C_{60},\qquad
C_2\oplus C_2\oplus C_{90},\qquad
C_2\oplus C_6\oplus C_{30}.
\end{gathered}
\]
每条整除链都满足不变因子条件，阶均为三百六十。主分解及素数幂循环因子的唯一性保证没有遗漏，互不相同的不变因子链保证没有重复。
\end{solution}

\begin{exercise}\label{b1:ex:09-free-quotient-splitting}
设 \(B\le A\)，其中 \(A\) 阿贝尔，且 \(A/B\) 为自由阿贝尔群。证明存在子群 \(C\) 使 \(A=B\oplus C\)。举例说明 \(C\) 未必唯一。
\end{exercise}
\begin{solution}
对 \(A/B\) 的每个基元素选一个提升 \(c_i\in A\)。自由阿贝尔群的泛性质给出同态 \(s:A/B\to A\)，把基元素送到对应提升；商映射 \(\pi\) 满足 \(\pi s=1\)。令 \(C=s(A/B)\)。任意 \(a\in A\) 可写成 \((a-s\pi(a))+s\pi(a)\)，首项属于 \(B\)，故 \(A=B+C\)。若 \(s(x)\in B\)，施以 \(\pi\) 得 \(x=0\)，故交为零。对 \(A=\mathbb Z\oplus C_2\)、\(B=0\oplus C_2\)，子群 \(\langle(1,\bar0)\rangle\) 与 \(\langle(1,\bar1)\rangle\) 都是这样的补群，且彼此不同。
\end{solution}

\begin{exercise}\label{b1:ex:09-cyclic-hom}
把同态集合按逐点相加组成阿贝尔群，证明
\[
\operatorname{Hom}(C_m,C_n)\cong C_{\gcd(m,n)}.
\]
进而计算从 \(C_{12}\oplus C_{18}\) 到 \(C_6\) 的同态数与满同态数。
\end{exercise}
\begin{solution}
同态由 \(\bar1\) 的像 \(x\in C_n\) 决定，条件是 \(mx=0\)。令 \(d=\gcd(m,n)\)，满足条件的元素恰为 \(\overline{(n/d)t}\)，其中 \(t\) 模 \(d\) 取值，在加法下构成 \(C_d\)。本题两生成元在 \(C_6\) 中都可以任取像，所以共有三十六个同态。若像由 \(\bar a,\bar b\) 生成，满射当且仅当 \(\gcd(a,b,6)=1\)。非满情况是两者都为二的倍数或都为三的倍数：分别有九种与四种，两者都满足有一种。故满同态数为 \(36-9-4+1=24\)。
\end{solution}

\begin{exercise}\label{b1:ex:09-minimal-generators}
设
\[
A\cong\mathbb Z^r\oplus C_{d_1}\oplus\cdots\oplus C_{d_s},
\qquad 1<d_1\mid\cdots\mid d_s.
\]
证明 \(A\) 的最少生成元数为 \(r+s\)。
\end{exercise}
\begin{solution}
各因子的标准生成元给出 \(r+s\) 个生成元。若 \(s>0\)，选素数 \(p\mid d_1\)，则 \(p\) 整除全部 \(d_i\)，所以 \(A/pA\cong C_p^{\,r+s}\)。任何 \(k\) 个生成元的像至多产生 \(p^k\) 个元素，因为每个整数系数可模 \(p\) 约减；而该商群有 \(p^{r+s}\) 个元素，故 \(k\ge r+s\)。若 \(s=0\)，任选素数 \(p\)，同样用 \(A/pA\cong C_p^r\) 得到下界 \(r\)。单位群对应 \(r=s=0\)，无需生成元。
\end{solution}

\begin{exercise}\label{b1:ex:09-torsion-layers}
设有限阿贝尔二群 \(A\) 满足
\[
|A/2A|=8,\quad |2A/4A|=4,\quad
|4A/8A|=2,\quad 8A=0.
\]
确定 \(A\)，并求其中阶恰为四的元素个数。
\end{exercise}
\begin{solution}
指数至少为一、二、三的循环因子分别有三、二、一个，因此指数恰为一、二、三的各有一个，故
\(A\cong C_2\oplus C_4\oplus C_8\)。被四消去的元素数为 \(2\cdot4\cdot4=32\)，被二消去的元素数为 \(2\cdot2\cdot2=8\)。阶恰为四的元素就是两者之差，共二十四个。
\end{solution}

\begin{exercise}\label{b1:ex:09-fg-endomorphism}
证明有限生成阿贝尔群的每个满自同态都是同构，并举出一个单自同态不是满射的例子。有限生成条件能否从第一个结论中删除？
\end{exercise}
\begin{solution}
设 \(f:A\to A\) 满射，\(T=A_{\mathrm{tor}}\)。它诱导 \(\bar f:A/T\to A/T\) 的满自同态。写 \(A/T\cong\mathbb Z^r\)，其整数矩阵 \(M\) 有整数右逆：为每个标准基选原像并列为矩阵 \(N\)，得 \(MN=I\)。行列式给出 \(\det M=\pm1\)，所以 \(\bar f\) 为同构。于是若 \(f(a)\in T\)，则 \(a\in T\)。满射性因而限制为 \(f|_T:T\to T\) 的满射；\(T\) 有限，所以这个限制为同构。若 \(f(a)=0\)，先得 \(a\in T\)，再得 \(a=0\)，故 \(f\) 同构。

单而不满的例子为 \(\mathbb Z\to\mathbb Z\)、\(n\mapsto2n\)。有限生成条件不能删除：在 \(\bigoplus_{i\ge1}C_2\) 上，把第一个标准生成元送零，把第 \(i+1\) 个送到第 \(i\) 个，便是有非平凡核的满自同态。
\end{solution}

